%%%%%%%%%%%%%%%%%%%%%%%%%%%%%%%%%%%%%%%%%%%%%%%%%%%%%%%%%%%%%%%%%%%%%%%%
%                                                                      %
%                       Chapter 10 - Section 1                         %
%                                                                      %
%%%%%%%%%%%%%%%%%%%%%%%%%%%%%%%%%%%%%%%%%%%%%%%%%%%%%%%%%%%%%%%%%%%%%%%%

\chapter{Series and Approximations}

An important theme in this book is to give {\bf constructive}
definitions\index{constructive definition} of mathematical objects.
Thus, for instance, if you needed to evaluate
     $$
        \int_0^1 e^{-x^2} \, dx,
        $$
you could set up a Riemann sum to evaluate this expression to any
desired degree of accuracy.  Similarly, if you wanted to evaluate a
quantity like $e^{.3}$ from first principles, you could apply Euler's
method to approximate the solution to the differential equation
\[        
y'(t) = y(t), \mbox{ with initial condition } y(0) = 1,
\]
using small enough intervals to get a value for $y(.3)$ to the number
of decimal places you needed.  You might pause for a moment to think
how you would get $\sin(5)$ to 7 decimal places---you wouldn't do it
by drawing a unit circle and measuring the $y$-coordinate of the point
where this circle is intersected by the line making an angle of 5
radians with the $x$-axis!  Defining the sine function to be the
solution to the second-order differential equation $y''=-y$ with
initial conditions $y=0$ and $y'=1$ when $t=0$ is much better if we
actually want to construct values of the function with more than two
decimal accuracy.

What these examples illustrate 
%
\mar{Ordinary arithmetic\\ lies at the heart\\ of all calculations}
%
is the fact that the only functions our brains or digital computers
can evaluate directly are those involving the arithmetic operations of
addition, subtraction, multiplication, and division. Anything else we
or computers evaluate must ultimately be reducible to these four
operations.  But the only functions directly expressible in such terms
are polynomials and rational functions (i.e., quotients of one
polynomial by another).  When you use your calculator to evaluate $\ln
2$, and the calculator shows .69314718056, it is really doing some
additions, subtractions, multiplications, and divisions to compute
this 11-digit {\em approximation\/} to $\ln 2$\,.  There are no
obvious connections to logarithms at all in what it does.  One of the
triumphs of calculus is the development of techniques for calculating
highly accurate approximations of this sort quickly.  In this chapter
we will explore these techniques and their applications.

%%%%%%%%%%%%%%%%%%%%%%%%%%%%%%%%%%%%%%%%%%%%%%%%%%%%%%%%%%%%%%%%%%%%%%%%
%                                                                      %
%                               Section 1                              %
%                                                                      %
%%%%%%%%%%%%%%%%%%%%%%%%%%%%%%%%%%%%%%%%%%%%%%%%%%%%%%%%%%%%%%%%%%%%%%%%

        
\section{Approximation Near a Point or\\ Over an Interval}

Suppose we were interested in approximating the sine function---we
might need to make a quick estimate and not have a calculator handy,
or we might even be designing a calculator.  In the next section we
will examine a number of other contexts in which such approximations
are helpful.  Here is a third degree polynomial that is a good
approximation in a sense which will be made clear shortly:
\[
P(x) = x - \dfrac{x^3}{6}.
\]
(You will see in section~2 where $P(x)$ comes from.)

If we compare the values of $\sin(x)$ and $P(x)$ over the interval
$[0, 1]$ we get the following:
\[
\begin{array}{cccc}
        x  & \sin{x} & P(x) &\sin{x}-P(x)\\ [2pt] \hline
        \begin{array}{r}
        \rule{0pt}{14pt} 0.0\\ .2 \\ .4 \\ .6 \\ .8 \\ 1.0 
        \end{array}
        &
        \begin{array}{r}
        \rule{0pt}{14pt} 0.0 \\ .198669 \\ .389418 \\ .564642 \\ .717356 \\ .841471  
        \end{array}
        &
        \begin{array}{r}
        \rule{0pt}{14pt} 0.0 \\ .198667 \\ .389333 \\ .564000 \\ .714667 \\ .833333
        \end{array}
        &
        \begin{array}{r}
        \rule{0pt}{14pt} 0.0 \\ .000002 \\ .000085 \\ .000642 \\ .002689 \\ .008138  
        \end{array}
\end{array}
\]
The fit is good, with the largest difference occurring at $x=1.0$,
where the difference is only slightly greater than $.008$.

If we plot $\sin(x)$ and $P(x)$ together over the interval $[0, \pi]$
we see the ways in which $P(x)$ is both very good and not so good.
Over the initial portion of the graph---out to around $x=1$---the
graphs of the two functions seem to coincide. As we move further from
the origin, though, the graphs separate more and more.  Thus if we
were primarily interested in approximating $\sin(x)$ near the origin,
$P(x)$ would be a reasonable choice.  If we need to approximate
$\sin(x)$ over the entire interval, $P(x)$ is less useful.

\vspace{125pt}

\special{psfile=../ps/calc2/sin1.ps}

On the other hand, consider the second degree polynomial 
$$Q(x) = -.4176977 x^2 +1.312236205 x - .050465497$$ (You will see how
to compute these coefficients in section~6.)  When we graph $Q(x)$ and
$\sin(x)$ together we get the following:

\vspace{110pt}

\special{psfile=../ps/calc2/sin2.ps}

While $Q(x)$ does not fits the graph of $\sin(x)$ as well as $P(x)$
does near the origin, it is a good fit overall.  In fact, $Q(x)$
exactly equals $\sin(x)$ at 4 values of $x$, and the greatest
separation between the graphs of $Q(x)$ and $\sin(x)$ over the
interval $[0,\pi]$ occurs at the endpoints, where the distance between
the graphs is .0505 units.

What we have here, then, are two kinds of approximation of the sine
function by polynomials: we have a polynomial $P(x)$ that behaves very
much like the sine function near the origin, and we have another
%
\mar{There's more than one way to make the "best fit" to a given curve}
%
polynomial $Q(x)$ that keeps close to the sine function over the
entire interval $[0, \pi]$. Which one is the ``better" approximation
depends on our needs.  Each solves an important problem.  Since
finding approximations near a point has a neater solution---Taylor
polynomials---we will start with this problem.  We will turn to the
problem of finding approximations over an interval in section~6.


\section{Taylor Polynomials}

{\bf The general setting.} In chapter 3 we discovered that functions
were locally linear at most points---when we zoomed in on them they
looked more and more like straight lines.  This fact was central to
the development of much of the subsequent material.  It turns out that
this is only the initial manifestation of an even deeper phenomenon:
Not only are functions locally linear, but, if we don't zoom in quite
so far, they look locally like parabolas.  From a little further back
still they look locally like cubic polynomials, etc.  Later in this
section we will see how to use the computer to visualize these ``local
parabolizations'', ``local cubicizations'', etc.  Let's summarize the
idea and then explore its significance:

\noindent
\begin{center}
\framebox[4.8in]
{\parbox{4.45in}
{\vspace{5pt}
The functions of interest to calculus look locally like polynomials at
most points of their domain.  The higher the degree of the polynomial,
the better typically will be the fit.
\vspace{5pt}
}
}
\end{center}

\noindent
{\bf Comments}\quad The ``at most points'' qualification is because of
exceptions like those we ran into when we explored locally linearity.
The function $|x|$, for instance, was not locally linear at
$x=0$---it's not locally like any polynomial of higher degree at that
point either.  The issue of what ``goodness of fit'' means and how it
is measured is a subtle one which we will develop over the course of
this section.  For the time being, your intuition is a reasonable
guide---one fit to a curve is better than another near some point if
it ``shares more phosphor'' with the curve when they are graphed on a
computer screen centered at the given point.

The fact that functions look locally like polynomials has profound
implications conceptually and
%
\mar{The behavior of a function can often be inferred from the
  behavior of a local polynomialization}
%
computationally.  It means we can often determine the behavior of a
function locally by examining the corresponding behavior of what we
might call a ``local polynomialization'' instead.  In particular, to
find the values of a function near some point, or graph a function
near some point, we can deal with the values or graph of a local
polynomialization instead.  Since we can actually evaluate polynomials
directly, this can be a major simplification.

There is an extra feature to all this which makes the concept
particularly attractive: not only are functions locally polynomial, 
%
\mar{We want the\\ best fit at $x=0$}
%
it
is easy to find the coefficients of the polynomials.  Let's see how
this works.  Suppose we had some function $f(x)$ and we wanted to find
the fifth degree polynomial that best fit this function at $x=0$.
Let's call this polynomial
\[
P(x) = a_0+a_1x+a_2x^2+a_3x^3+a_4x^4+a_5x^5.
\]
To determine $P$, we need to find values for the six coefficients
$a_0$, $a_1$, $a_2$, $a_3$, $a_4$, $a_5$.

Before we can do this, we need to define what we mean by the ``best''
fit to $f$ at $x=0$.  Since we have six unknowns, we need six
conditions.  One obvious 
%
\mar{The best fit should pass through the point $(0, f(0))$}
%
condition is that the graph of $P$ should pass through the point $(0,
f(0))$.  But this is equivalent to requiring that $P(0)=f(0)$.  Since
$P(0)=a_0$, we thus must have $a_0=f(0)$, and we have found one of the
coefficients of $P(x)$.  Let's summarize the argument so far:

\noindent
\begin{center}
\framebox[4.8in]
{\parbox{4.45in}{\vspace{8pt}
The graph of a polynomial passes through the point $(0,f(0))$
   if and only if the polynomial is of the form
    $$f(0)+a_1x+a_2x^2+\cdots .$$}
}
\end{center}

But we're 
%
\mar{The best fit should have the right slope at $(0, f(0))$}
%
not interested in just any polynomial passing through the right point;
it should be headed in the right direction as well.  That is, we want
the slope of $P$ at $x=0$ to be the same as the slope of $f$ at this
point---we want $P'(0)=f'(0).$ But
\[
P'(x) = a_1 + 2a_2x + 3a_3x^2 + 4a_4x^3 + 5a_5x^4,
\]
so $P'(0)=a_1$.  Our second condition therefore must be that
$a_1=f'(0)$.  Again, we can summarize this as

\noindent
\begin{center}
\framebox[4.8in]
{\parbox{4.45in}{\vspace{8pt}
The graph of a polynomial passes through the point $(0, f(0))$
  and has slope $f'(0)$ there if and only if it is of the form
  $$f(0) + f'(0)x + a_2x^2 + \cdots .$$}
}
\end{center}

\noindent
Note that at this point we have recovered the general form for the
local linear approximation to $f$ at $x=0$: $L(x) = f(0) + f'(0)x$.

But there is no reason to stop with the first derivative.  Similarly,
we would want the way in which the slope of $P(x)$ is changing---we
are now talking about $P''(0)$---to behave the way the slope of $f$ is
changing at $x=0$, etc.  Each higher derivative controls a more subtle
feature of the shape of the graph.  We now see how we could formulate
reasonable additional conditions which would determine the remaining
coefficients of~$P(x)$:

\medskip

\noindent
\begin{center}
\framebox[4.9in]
{\parbox{4.5in}{\vspace{4pt}
Say that $P(x)$ is the {\bf best fit}\index{best fit at a point} to
$f(x)$ at the point $x=0$ if\vspace{-8pt}
\[
P(0)=f(0),\ P'(0)=f'(0),\ P''(0)=f''(0),\ldots ,P^{(5)}(0)=f^{(5)}(0).
\]\vspace{-20pt}
}
}
\end{center}

\medskip

\enlargethispage*{20pt} Since $P(x)$ is a fifth degree polynomial, all
the derivatives of $P$ beyond the fifth will be identically 0, so we
can't control their values by altering the values of the $a_k$. What
we are saying, then, is that we are using as our
%
\mar{The final criterion\\ for best fit at $x=0$}
%
criterion for the best fit that all the derivatives of $P$ as high as
we can control them have the same values at $x=0$ as the corresponding
derivatives of $f$.

While this is a reasonable definition for something we might call the
``best fit'' at the point $x=0$, it gives us no direct way to tell how
good the fit really is.  This is a serious shortcoming---if we want to
approximate function values by polynomial values, for instance, we
would like to know how many decimal places in the polynomial values
are going to be correct.  We will take up this question of goodness of
fit later in this section; we'll be able to make measurements that
allow us to to see how well the polynomial fits the function.  First,
though, we need to see how to determine the coefficients of the
approximating polynomials and get some practice manipulating them.

\smallskip

{\bf Note on Notation:}  We have used the
%
\mar{Notation for\\ higher derivatives}
%
notation $f^{(5)}(x)$ to denote the fifth derivative of $f(x)$ as a
convenient shorthand for $f'''''(x)$, which is harder to read.  We
will use this throughout.

\smallskip

{\bf Finding the coefficients}\quad We first observe that the
derivatives of $P$ at $x=0$ are easy to express in terms of
$a_1,a_2,\ldots \, .$ We have
\begin{align*}
P'(x) &= a_1+2\,a_2x+3\,a_3x^2+4\,a_4x^3+5\,a_5x^4, \\
P''(x)&= 2\,a_2+3\cdot2\,a_3x+4\cdot3\,a_4x^2+5\cdot4\,a_5x^3, \\
P^{(3)}(x)&= 3\cdot2\,a_3+4\cdot3\cdot2\,a_4x+5\cdot4\cdot3\,a_5x^2, \\
P^{(4)}(x)&= 4\cdot3\cdot2\,a_4+5\cdot4\cdot3\cdot2\,a_5x, \\
P^{(5)}(x)&= 5\cdot4\cdot3\cdot2\,a_5.
\end{align*}
Thus $P''(0)=2\,a_2$, $P^{(3)}(0)= 3\cdot2\,a_3$, $P^{(4)}(0)=
4\cdot3\cdot2\,a_4$, and $P^{(5)}(0)=5\cdot4\cdot3\cdot2\,a_5 \, .$

We can simplify this a bit by introducing the {\bf
  factorial}\index{factorial} notation, in which we write
$n!=n\cdot(n-1)\cdot(n-2)\cdots3\cdot2\cdot1 \, .$ This is called
``$n$ factorial".
%
\mar{Factorial notation}
%
Thus, for example, $7!=7\cdot6\cdot5\cdot4\cdot3\cdot2\cdot1 = 5040.$
It turns out to be convenient to extend the factorial notation to 0 by
defining $0!=1.$ (Notice, for instance, that this makes the formulas
below work out right.)  In the exercises you will see why this
extension of the notation is not only convenient, but reasonable as
well!

With this notation 
%
\mar{The desired rule for finding the coefficients}
%
we can express compactly the equations above as $P^{(k)}(0)=k!\,a_k$
for $k=0$, 1, 2, \ldots 5 \,.  Finally, since we want
$P^{(k)}(0)=f^{(k)}(0)$, we can solve for the coefficients of $P(x)$:
\[
a_k = \frac{f^{(k)}(0)}{k!} \quad \text{for $k=0,1,2,3,4,5$}.
\]

We can now write down an explicit formula for the fifth degree
polynomial which best fits $f(x)$ at $x=0$ in the sense we've put
forth:
\[
P(x) = f(0) + f'(0)x + \frac{f^{(2)}(0)}{2!}x^2 
   + \frac{f^{(3)}(0)}{3!}x^3 + \frac{f^{(4)}(0)}{4!}x^4 
   + \frac{f^{(5)}(0)}{5!}x^5.
\]
We can express this more compactly using the $\Sigma$--notation we
introduced in the discussion of Riemann sums in chapter~6:
\[
P(x)=\sum_{k=0}^{5} \frac{f^{(k)}(0)}{k!}\,x^k.
\]
We call this the {\bf fifth degree Taylor polynomial for $f(x)$}.  It
is sometimes also called the fifth {\bf order} Taylor polynomial.

It should be obvious to you that we can generalize what we've done
above to get a
%
\mar{\vspace{30pt}General rule for\\ the Taylor polynomial\\ at $x=0$}
%
best-fitting polynomial\index{polynomial!best-fitting} of any degree.
Thus
\begin{center}
\framebox[4.4in]
{\parbox{3.85in}{\vspace{8pt}
The {\bf Taylor polynomial\index{Taylor polynomial!definition} of degree} {\boldmath $n$} approximating the function $f(x)$ at $x=0$ is given by the formula
\[
P_n(x)=\sum_{k=0}^{n} \frac{f^{(k)}(0)}{k!}\,x^k.
\]}
}
\end{center}
We also speak of the Taylor polynomial {\em centered at} $x=0$.

\medskip

\noindent
{\bf Example}.  \quad Consider $f(x) = \sin(x)$. Then for $n=7$ we have 
\begin{alignat*}{2}
f(x)       &= \sin(x),  &\qquad f(0)       &= 0,   \\
f'(x)      &= \cos(x),  &       f'(0)      &= +1,  \\
f^{(2)}(x) &= -\sin(x), &       f^{(2)}(0) &= 0,   \\
f^{(3)}(x) &= -\cos(x), &       f^{(3)}(0) &= -1,  \\
f^{(4)}(x) &= \sin(x),  &       f^{(4)}(0) &= 0,   \\
f^{(5)}(x) &= \cos(x),  &       f^{(5)}(0) &= +1,  \\
f^{(6)}(x) &= -\sin(x), &       f^{(6)}(0) &= 0,   \\
f^{(7)}(x) &= -\cos(x), &       f^{(7)}(0) &= -1. 
\end{alignat*}
From this we can see that the pattern 0, $+1$, 0 , $-1$, \ldots will
repeat forever.  Substituting these values into the formula we get
that for any odd integer $n$ the $n$-th degree Taylor polynomial for
$\sin(x)$ is
\[
        P_n(x) =  x - \dfrac{x^3}{3!} + \dfrac{x^5}{5!} 
                - \dfrac{x^7}{7!} + \cdots \pm {x^n}{n!}.
\]

Note that $P_3(x)=x-x^3/6$, which is the polynomial we met in
section~1. We saw there that this polynomial seemed to fit the graph
of the sine function only out to around $x=1$.  Now, though, we have a
way to generate polynomial approximations of higher degrees, and we
would expect to get better fits as the degree of the approximating
polynomial is increased.  To see how closely these polynomial
approximations follow $\sin(x)$, here's the graph of $\sin(x)$
together with the Taylor polynomials of degrees $n=1,3,5,\ldots ,17$
plotted over the interval $[0, 7.5]$:\label{sin3}

\vspace*{140pt}

\special{psfile=../ps/calc2/sin3.ps}\index{Taylor polynomial!graphs for $\sin{x}$}

\newpage

While each polynomial eventually wanders off
%
\mar{The higher the degree of the polynomial,\\ the better the fit}
%
to infinity, successive polynomials stay close to the sine function
for longer and longer intervals---the Taylor polynomial of degree 17
is just beginning to diverge visibly by the time $x$ reaches $2\pi$.
We might expect that if we kept going, we could find Taylor
polynomials that were good fits out to $x=100$, or $x=1000$.  This is
indeed the case, although they would be long and cumbersome
polynomials to work with.  Fortunately, as you will see in the
exercises, with a little cleverness we can use a Taylor polynomial of
degree 9 to calculate $\sin(100)$ to 5 decimal place accuracy.

\medskip
\noindent
{\bf Other Taylor Polynomials:}  
%
\mar{Approximating polynomials for\\ other basic functions}
%
In a similar fashion, we can get Taylor polynomials for other
functions.  You should use the general formula to verify the Taylor
polynomials for the following basic functions.  (The Taylor polynomial
for $\sin(x)$ is included for convenient reference.)
$$
        \begin{array}{ccl}
        f(x)  &\rule{30pt}{0pt}& \rule{60pt}{0pt} P_n(x) \\ [2pt] \hline
        & \\ [-6pt]
        \sin(x) & & x - \dfrac{x^3}{3!} + \dfrac{x^5}{5!}
                - \dfrac{x^7}{7!} + \cdots \pm \dfrac{x^n}{n!} \quad
        (n \mbox{ odd})\\ [.2in]
        \cos(x) & & 1 - \dfrac{x^2}{2!} + \dfrac{x^4}{4!} 
                - \dfrac{x^6}{6!} + \cdots \pm \dfrac{x^n}{n!} \quad
                (n \mbox{ even})\\ [.2in]
        e^x & & 1 + x + \dfrac{x^2}{2!} + \dfrac{x^3}{3!} 
                + \dfrac{x^4}{4!} + \cdots + \dfrac{x^n}{n!} \\ [.2in]
        \ln{(1-x)} & & -\left(x+ \dfrac{x^2}{2} + \dfrac{x^3}{3} 
                + \dfrac{x^4}{4} + \cdots + \dfrac{x^n}{n}\right) \\ [.2in]
         \dfrac{1}{1-x} & &{\displaystyle 1+x+x^2+x^3+\cdots +x^n}
        \end{array}
$$

\medskip
\noindent
{\bf Taylor polynomials at points other than \boldmath $x=0$.}  
%
\mar{General rule for\\ the Taylor polynomial\\ at $x=a$}
%
Using exactly the same arguments we used to develop the best-fitting
polynomial at $x=0$, we can derive the more general formula for the
best-fitting polynomial at any value of $x$.  Thus, if we know the
behavior of $f$ and its derivatives at some point $x=a$, we would like
to find a polynomial $P_n(x)$ which is a good approximation to $f(x)$
for values of $x$ close to $a$.

Since the expression $x-a$ tells us how close $x$ is to $a$, we use it
(instead of the variable $x$ itself) to construct the polynomials
approximating $f$ at $x=a$:
$$
P_n(x)=b_0+b_1(x-a)+b_2(x-a)^2+b_3(x-a)^3+\cdots+b_n(x-a)^n.
$$
You should be able to apply the reasoning we used above to derive the
following:
\begin{center}
\framebox[5.3in]
{\parbox{4.85in}{\vspace{8pt}
The {\bf Taylor polynomial of degree} {\boldmath $n$} {\bf centered at} {\boldmath $x=a$} approximating the function $f(x)$  is given by the formula
\begin{align*}
P_n(x)&= f(a) +f'(a)(x-a) + \frac{f''(a)}{2!}(x-a)^2+\cdots 
               + \frac{f^{n}(a)}{n!}(x-a)^n \\
      &= \sum_{k=0}^{n} \frac{f^{(k)}(a)}{k!}(x-a)^k.
\end{align*}\vspace{-6pt}}
}
\end{center}\index{Taylor polynomial!centered at a point other than the origin}


%This still needs fixing -- I can't get this program to run

\begin{center}
{\bf Program: TAYLOR }
\end{center}
\noindent
{\sl Set up GRAPHICS}
\vspace{-8pt}
\begin{verbatim}
DEF fnfact(m)
     P = 1
     FOR r = 2 TO m
         P = P * r
     NEXT r
     fnfact = P
END DEF
DEF fnpoly(x)
     Sum = x
     Sign = -1
     FOR k = 3 TO 17 STEP 2
         Sum = Sum + Sign * x^k/fnfact(k)
         Sign = (-1) * Sign
     NEXT k
     fnpoly = Sum
END DEF
FOR x = 0 TO 3.14 STEP .01
\end{verbatim}
\vspace{-8pt}
\rule{24pt}{0pt} {\sl Plot the line from 
                {\tt (x, fnpoly(x))}} {\sl to} 
                {\tt (x + .01, fnpoly(x + .01))}\\
\vspace{-26pt}
\begin{verbatim}
NEXT x
\end{verbatim}

\noindent
{\bf A computer program for graphing Taylor polynomials} Shown above
is a program that evaluates the 17-th degree Taylor polynomial for
$\sin(x)$ and graphs it over the interval $[0, 3.14]$.  The first
seven lines of the program constitute a subroutine for evaluating
factorials.  The syntax of such subroutines varies from one computer
language to another, so be sure to use the format that's appropriate
for you.  You may even be using a language that already knows how to
compute factorials, in which case you can omit the subroutine.  The
second set of 9 lines defines the function {\tt poly} which evaluates
the 17-th degree Taylor polynomial.  Note the role of the variable
{\tt Sign}---it simply changes the sign back and forth from positive
to negative as each new term is added to the sum.  As usual, you will
have to put in commands to set up the graphics and draw lines in the
format your computer language uses.  You can modify this program to
graph other Taylor polynomials.  \index{program!TAYLOR}



\subsection{New Taylor Polynomials from Old}

Given a function we want to approximate by Taylor polynomials, we
could always go straight to the general formula for deriving such
polynomials.  On the other hand, it is often possible to avoid a lot
of tedious calculation of derivatives by using a polynomial we've
already calculated. It turns out that any manipulation on Taylor
polynomials you might be tempted to try will probably work.  Here are
some examples to illustrate the kinds of manipulations that can be
performed on Taylor polynomials.

\enlargethispage*{10pt}
\medskip
\noindent
{\bf Substitution in Taylor Polynomials.}  Suppose we wanted the
Taylor polynomial for $e^{x^2}$.  We know from what we've already done
that for any value of $u$ close to 0,
$$
e^u \approx 1 + u + \dfrac{u^2}{2!} + \dfrac{u^3}{3!} 
                + \dfrac{u^4}{4!} + \cdots + \dfrac{u^n}{n!}.
$$
In this expression $u$ can be anything, including another variable
expression.  For instance, if we set $u=x^2$, we get the Taylor
polynomial
\begin{align*}
e^{x^2} &= e^u \\[-3pt]
        &\approx  1 + u + \dfrac{u^2}{2!} + \dfrac{u^3}{3!} 
                + \dfrac{u^4}{4!} + \cdots + \dfrac{u^n}{n!} \\[3pt]
        &= 1 + (x^2) + \dfrac{(x^2)^2}{2!} + \dfrac{(x^2)^3}{3!} 
          + \dfrac{(x^2)^4}{4!} + \cdots + \dfrac{(x^2)^n}{n!} \\[3pt]
        &= 1 + x^2 + \dfrac{x^4}{2!} + \dfrac{x^6}{3!} + \dfrac{x^8}{4!} 
             + \cdots + \dfrac{x^{2n}}{n!}.
\end{align*}
You should check to see that this is what you get if you apply the
general formula for computing Taylor polynomials to the function
$e^{x^2}$.


\pagebreak

Similarly, suppose we wanted a Taylor polynomial for $1/(1+x^2)$.  We
could start with the approximation given earlier:
$$
\dfrac{1}{1-u} \approx 1+u+u^2+u^3+\cdots +u^n.
$$
If we now replace $u$ everywhere by $-x^2$, we get the desired expansion:
\begin{align*}
\dfrac {1}{1+x^2} =  \dfrac{1}{1-(-x^2)}
        &=  \dfrac{1}{1-u}  \\
        &\approx  1+u+u^2+u^3+\cdots +u^n \\[3pt]
        &=  1+(-x^2)+(-x^2)^2+(-x^2)^3+\cdots +(-x^2)^n \\[3pt]
        &=  1-x^2+x^4-x^6+\cdots \pm x^{2n}.
\end{align*}
Again, you should verify that if you start with $f(x)=1/(1+x^2)$ and
apply to $f$ the general formula for deriving Taylor polynomials, you
will get the preceding result.  Which method is quicker?
 
\medskip
\noindent
{\bf Multiplying Taylor Polynomials.}  Suppose we wanted the 5-th
degree Taylor polynomial for $e^{3x}\cdot \sin(2x)$.  We can use
substitution to write down polynomial approximations for $e^{3x}$ and
$\sin(2x)$, so we can get an approximation for their product by
multiplying the two polynomials:
\begin{align*}
\lefteqn{e^{3x} \cdot \sin(2x)}\quad \\
&\approx \left(1 + (3x)+ \dfrac{(3x)^2}{2!} + \dfrac{(3x)^3}{3!} +
      \dfrac{(3x)^4}{4!} + \dfrac{(3x)^5}{5!}\right)
    \left((2x) - \dfrac{(2x)^3}{3!} + \dfrac{(2x)^5}{5!}\right) \\
&\approx  2x + 6x^2 + \dfrac{23}{3}x^3 + 5x^4 - \dfrac{61}{60}x^5.
\end{align*}
Again, you should try calculating this polynomial directly from the
general rule, both to see that you get the same result, and to
appreciate how much more tedious the general formula is to use in this
case.

In the same way, we can also divide Taylor polynomials, raise them to
powers, and chain them by composition. The exercises provide examples
of some of these operations.

\medskip
\noindent
{\bf Differentiating Taylor Polynomials.}  Suppose we know a Taylor
polynomial for some function $f$. If $g$ is the derivative of $f$, we
can immediately get a Taylor polynomial for $g$ (of degree one less)
by differentiating the polynomial we know for $f$. You should review
the definition of Taylor polynomial to see why this is so.  For
instance, suppose $f(x) = 1/(1-x)$ and $g(x) =1/(1-x)^2 $.  Verify
that $f'(x) = g(x)$. It then follows that
\begin{align*}
\dfrac{1}{(1-x)^2} 
   = \dfrac{d}{dx} \left(\dfrac{1}{1 - x}\right)
   &\approx \dfrac{d}{dx} (1 + x + x^2 + \cdots + x^n) \\[3pt]
   &= 1 + 2x + 3x^2 + \cdots + nx^{n-1}.
\end{align*} 

\medskip
\noindent
{\bf Integrating Taylor Polynomials.}  Again suppose we have functions
$f(x)$ and $g(x)$ with $f'(x) = g(x)$, and suppose this time that we
know a Taylor polynomial for $g$.  We can then get a Taylor polynomial
for $f$ by antidifferentiating term by term.  For instance, we find in
chapter 11 that the derivative of $\arctan(x)$ is $1/(1+x^2)$, and we
have seen above how to get a Taylor polynomial for $1/(1+x^2)$.
Therefore we have
\begin{align*}
\arctan{x} = \int_0^x \frac{1}{1+t^2} dt
  &\approx \int_0^x \left( 1-t^2+t^4-t^6+ \cdots \pm t^{2n}\right) dt \\[3pt]
  &=  \left. t-\dfrac{1}{3}t^3 + \dfrac{1}{5} t^5 
          - \cdots \pm \dfrac{1}{2n+1}t^{2n+1} \; \right|_0^x \\[3pt]
  &= x - \dfrac{1}{3}x^3 + \dfrac{1}{5} x^5 - \cdots \pm 
            \dfrac{1}{2n+1}x^{2n+1}.
\end{align*}


\subsection{Goodness of fit}

Let's turn to the question of {\em measuring\/}
%
\mar{Graph the difference between a function and its Taylor polynomial}
%
the fit between a function and one of its Taylor polynomials.  The
ideas here have a strong geometric flavor, so you should use a
computer graphing utility to follow this discussion.  Once again,
consider the function $\sin(x)$ and its Taylor polynomial $P(x) = x -
x^3/6$.  According to the table in section~1, the difference $\sin(x) -
P(x)$ got smaller as $x$ got smaller.  Stop now and graph the function
$y = \sin(x) - P(x)$ near $x = 0$.  This will show you exactly how
$\sin(x) - P(x)$ depends on $x$.  If you choose the interval $-1 \leq
x \leq 1$ (and your graphing utility allows its vertical and
horizontal scales to be set independently of each other), your graph
should resemble this one.\label{startsine}


\vspace{70pt}
\special{psfile=../ps/calc2/j101.ps}

\newpage

This graph looks very much like a cubic polynomial.  
%
\mar{The difference looks like a power of $x$} 
%
If it really is a cubic, we can figure out its formula, because we
know the value of $\sin(x) - P(x)$ is about .008 when $x = 1$.
Therefore the cubic should be $y = .008\, x^3$ (because then $y =
.008$ when $x = 1$).  However, if you graph $y = .008 \, x^3$ together
with $y = \sin(x) - P(x)$, you should find a poor match (the left-hand
figure, below.)  Another possibility is that $\sin(x) - P(x)$ is more
like a {\em fifth\/} degree polynomial.  Plot $y = .008 \, x^5$; it's
so close that it ``shares phosphor'' with $\sin(x) - P(x)$ near $x =
0$.


 
\special{psfile=../ps/calc2/j102.ps}
\special{psfile=../ps/calc2/j103.ps}

\vspace{120pt}

If $\sin(x) - P(X)$ were {\em exactly\/} a multiple of $x^5$, 
%
\mar{Finding the multiplier}
%
then $(\sin{x} - P(x))/x^5$ would be constant and would equal the value
of the multiplier.  What we actually find is this:
$$
{\small
        \begin{array}{lc}
        \multicolumn{1}{c}{x} & 
                \dfrac{\sin{x} - P(x)}{x^5} \\ [8pt]
        \hline 
        1.0\rule{0pt}{14pt}     & .0081377 \\
        0.5     & .0082839 \\
        0.1     & .0083313 \\ 
        0.05    & .0083328 \\
        0.01    & .0083333
        \end{array}
}       
        \qquad \mbox{suggesting} \quad
        \lim_{x \to 0} \dfrac{\sin{x} - P(x)}{x^5} = .008333\ldots .
$$
Thus, although the ratio is not constant,  
%
\mar{How $P(x)$ fits $\sin(x)$}
%
it appears to converge to a definite value---which we can take to be
the value of the multipier:
        $$
        \sin{x} - P(x) \approx .008333 \, x^5 \quad 
                \mbox{when} \quad x \approx 0.
        $$
We say that $\sin(x) - P(x)$ {\em has the same order of magnitude
  as\/} $x^5$ as $x \to 0$.  So $\sin(x) - P(x)$ is about as small as
$x^5$.  Thus, if we know the size of $x^5$ we will be able to tell how
close $\sin(x)$ and $P(x)$ are to each other.

A rough way to measure how close two numbers are 
%
\mar{Comparing\\ two numbers}     
%
is to count the number of decimal places to which they agree.  But
there are pitfalls here; for instance, none of the decimals of 1.00001
and 0.99999 agree, even though the difference between the two numbers
is only 0.00002.  This suggests that a good way to compare two numbers
is to look at their difference.  Therefore, we say
\begin{center}
        $A = B$ to $k$ decimal places 
        \quad {\em means\/}  \quad
        $A - B = 0$ to $k$ decimal places
\end{center}
Now, a number equals 0 to k decimal places precisely when it {\em
  rounds off to\/} 0 (when we round it to k decimal places).  Since
$X$ rounds to 0 to $k$ decimal places if and only $|X| < .5 \times
10^{-k}$, we finally have a precise way to compare the size of two
numbers:
\begin{center}\framebox[360pt]
{
\parbox{324pt}
        {\vspace{5pt}
        \hfill 
        $A = B$ to $k$ decimal places 
        \quad {\em means\/}  \quad
        $|A - B| < .5 \times 10^{-k}$.
        \hfill
        \vspace{5pt}
        }
}
\end{center}

Now we can say how close $P(x)$ is to $\sin(x)$.  
%      
\mar{What the fit means computationally}        
%
Since $x$ is small, we can take this to mean {\em $x = 0$ to $k$
  decimal places}, or $|x| < .5 \times 10^{-k}$.  But then,
\[
|x^5 - 0| = |x - 0|^5 < (.5 \times 10^{-k})^5 < .5 \times 10^{-5k - 1}
\]
(since $.5^5 = .03125 < .5 \times 10^{-1}$).  In other words, if $x =
0$ to $k$ decimal places, then $x^5 = 0$ to $5k + 1$ places.  Since
$\sin(x) - P(x)$ has the same order of magnitude as $x^5$ as $x \to
0$, $\sin(x) = P(x)$ to $5k + 1$ places as well.  In fact, because the
multiplier in the relation
\[
\sin{x} - P(x) \approx .008333 \, x^5 \qquad  (x \approx 0)
\]
is .0083\ldots , we gain two more decimal places of accuracy.  (Do you
see why?)  Thus, finally, we see how reliable the polynomial $P(x) = x
- x^3/6$ is for calculating values of $\sin(x)$:
\begin{center}\framebox[360pt]
{
\parbox{324pt}
        {
        \vspace{5pt}
        When $x = 0$ to $k$ decimal places of accuracy, we can use $P(x)$ to calculate the first $5k + 3$ decimal places of the value of $\sin(x)$. 
        \vspace{-3pt}
        }
}
\end{center}
Here are a few examples comparing $P(x)$ to the {\em exact\/} value of $\sin(x)$:
{\small
\addtolength{\arraycolsep}{8pt}
        $$
        \begin{array}{lll}
        \multicolumn{1}{c}{x} & \multicolumn{1}{c}{P(x)} &
                \multicolumn{1}{c}{\sin(x)} \\[2pt]
        \hline
        .0372   &  .\underline{03719142}01920 & 
             .\underline{0371942}07856\ldots \rule{0pt}{14pt} \\
        .0086   & .\underline{00859989}39907 &
             .\underline{00859989}3991\ldots \\
        .0048   & .\underline{0047999815680}000 &
                 .\underline{0047999815680}212\ldots
        \end{array}
        $$ 
} 
The underlined digits are guaranteed to be correct, based on the
number of decimal places for which $x$ agrees with 0.  (Note that,
according to our rule, $.0086 = 0$ to {\em one\/} decimal place, not
two.)\label{endsine}
%\addtolength{\arraycolsep}{-8pt}

\subsection{Taylor's theorem}

Taylor's theorem 
%
\mar{Order of magnitude}   
%
is the generalization of what we have just seen; it describes the
goodness of fit between an arbitrary function and one of its Taylor
polynomials.  We'll state three versions of the theorem, gradually
uncovering more information.  To get started, we need a way to compare
the order of magnitude of {\em any\/} two functions.

\medskip

\begin{center}\framebox[360pt]
{
\parbox{320pt}
        {
        \vspace{5pt}
        We say that $\varphi(x)$ has {\bf the same order of magnitude} as $q(x)$ as $x \to a$, and we write $\varphi(x) = O(q(x))$ as $x \to a$, if there is a constant $C$ for which
        $$
        \lim_{x \to a} \dfrac{\varphi(x)}{q(x)} = C.
        $$
        \vspace{-10pt}
        }
}
\end{center}

\medskip

\noindent
Now, when $\lim_{x\to a} \varphi(x)/q(x)$ is $C$, we have \label{order}
\[
\varphi(x) \approx C q(x) \quad \mbox{when $x \approx a$}.
\]
We'll frequently use this relation to express the idea that
$\varphi(x)$ has the same order of magnitude as $q(x)$ as $x \to a$.

The symbol $O$ is an upper case ``oh''.  
%
\mar{`Big oh' notation} 
%
When $\varphi(x) = O(q(x))$ as $x \to a$, we say {\em $\varphi(x)$ is
  `big oh' of $q(x)$ as $x$ approaches $a$}.  Notice that the equal
sign in $\varphi(x) = O(q(x))$ does {\em not\/} mean that $\varphi(x)$
and $O(q(x))$ are equal; $O(q(x))$ isn't even a function.  Instead,
the equal sign and the $O$ together tell us that $\varphi(x)$ stands
in a certain relation to $q(x)$.

\medskip

\begin{center}\framebox[360pt]
{
\parbox{330pt}
        {
        \vspace{5pt}
{\bf Taylor's theorem, version~1}.  If $f(x)$ has derivatives up to order $n$ at $x = a$, then
        $$
        f(x) = f(a) + \dfrac{f'(a)}{1!}(x - a) + \cdots + 
                \dfrac{f^{(n)}(a)}{n!} (x - a)^n + R(x), 
        $$
where $R(x) = O((x - a)^{n+1})$ as $x \to a$.  The term $R(x)$ is called the {\bf remainder}.
\vspace{8pt}
        }
}
\end{center}

\medskip

\enlargethispage*{20pt}
This version of Taylor's theorem focusses 
%
\mar{Informal language} 
%
on the general shape of the remainder function.  Sometimes we just say
the remainder has ``order $n + 1$'', using this short phrase as an
abbreviation for ``the order of magnitude of the function $(x -
a)^{n+1}$''.  In the same way, we say that {\em a function and its
  $n$-th degree Taylor polynomial at $x = a$ agree to order $n+1$ as
  $x \to a$}.

\newpage

\aside{Notice that, if $\varphi(x) = O(x^3)$ as $x \to 0$, then it is also true that $\varphi(x) = O(x^2)$ (as $x \to 0$).  This implies that we should take $\varphi(x) = O(x^n)$ to mean ``$\varphi$ has {\em at least\/} order $n$'' (instead of simply ``$\varphi$ has order $n$'').  In the same way, it would be more accurate (but somewhat more cumbersome) to say that $\varphi = O(q)$ means ``$\varphi$ has {\em at least\/} the order of magnitude of $q$''.}
 
As we saw in our example, we can translate the order  
%
\mar{Decimal places\\ of accuracy}        
%
of agreement between the function and the polynomial into information
about the number of decimal places of accuracy in the polynomial
approximation.  In particular, if $x - a = 0$ to $k$ decimal places,
then $(x - a)^n = 0$ to $nk$ places, at least.  Thus, as the order of
magnitude $n$ of the remainder increases, the fit increases, too.
(You have already seen this illustrated with the sine function and its
various Taylor polynomials, in the figure on page~\pageref{sin3}.)

While the first version of Taylor's theorem tells us that $R(x)$ 
%
\mar{A formula for\\ the remainder}       
%
looks like $(x - a)^{n+1}$ in some general way, the next gives us a
concrete formula.  At least, it {\em looks\/} concrete.  Notice,
however, that $R(x)$ is expressed in terms of a number $c_x$ (which
depends upon $x$), but the formula doesn't tell us {\em how\/} $c_x$
depends upon $x$.  Therefore, if you want to use the formula to {\em
  compute\/} the value of $R(x)$, you can't.  The theorem says only
that $c_x$ exists; it doesn't say how to find its value.
Nevertheless, this version provides useful information, as you will
see.
\begin{center}\framebox[360pt]
{
\parbox{330pt}
        {
        \vspace{5pt}
{\bf Taylor's theorem, version~2}.  Suppose $f$ has continuous derivatives up to order $n+1$ for all $x$ in some interval containing~$a$.  Then, for each $x$ in that interval, there is a number $c_x$ between $a$ and $x$ for which
        $$
        R(x) =  \dfrac{f^{(n+1)}(c_x)}{(n+1)!} \, (x - a)^{n+1}. 
        $$
This is called {\bf Lagrange's form of the remainder}. 
\vspace{8pt}
        }
}
\end{center}

We can use the Lagrange form 
%
\mar{Another formula for the remainder}   
%
as an aid to computation.  To see how, return to the formula 
        $$
        R(x) \approx C (x - a)^{n+1} \quad (x \approx a)
        $$
that expresses $R(x) = O((x - a)^{n+1})$ as $x \to a$ (see page~\pageref{order}).  The constant here is the limit
        $$
        C = \lim_{x \to a} \dfrac{R(x)}{(x - a)^{n+1}}.
        $$
If we have a good estimate for the value of $C$, then $R(x) \approx C (x - a)^{n+1}$ gives us a good way to estimate $R(x)$.   Of course, we could just evaluate the limit to determine $C$.  In fact, that's what we did in the example; knowing $C \approx .008$ there gave us two more decimal places of accuracy in our polynomial approxmation to the sine function.   

But the Lagrange form of the remainder 
%
\mar{Determining $C$\\ from $f$ at $x = a$}       
%
gives us another way to determine~$C$:
\begin{align*}
C = \lim_{x \to a} \dfrac{R(x)}{(x - a)^{n+1}} 
   &= \lim_{x \to a} \dfrac{f^{(n+1)}(c_x)}{(n+1)!} \\[3pt]
   &= \dfrac{f^{(n+1)}(\lim_{x \to a} c_x)}{(n+1)!} \\[3pt]
   &= \dfrac{f^{(n+1)}(a)}{(n+1)!}.
\end{align*}
In this argument, we are permitted to take the limit ``inside'' $f^{(n+1)}$ because $f^{(n+1)}$ is a continuous function.   (That is one of the hypotheses of version~2.)  Finally, since $c_x$ lies between $x$ and $a$, it follows that $c_x  \to a$ as $x \to a$; in other words, $\lim_{x \to a}c_x = a$.  Consequently, we get $C$ {\em directly\/} from the function $f$ itself, and we can therefore write
        $$
        R(x) \approx \dfrac{f^{(n+1)}(a)}{(n+1)!} \, (x - a)^{n+1} \qquad
                (x \approx a).
        $$

The third version of Taylor's theorem 
%
\mar{An error bound}
%
uses the Lagrange form of the remainder in a similar way to get an
{\em error bound\/} for the polynomial approximation based on the size
of $f^{(n+1)}(x)$.
\begin{center}\framebox[360pt]
{
\parbox{330pt}
        {
        \vspace{5pt}
{\bf Taylor's theorem, version~3}.  Suppose that $|f^{(n+1)}(x)| \leq M$ for all $x$ in some interval containing~$a$.  Then, for each $x$ in that interval,
        $$
        |R(x)| \leq  \dfrac{M}{(n+1)!} \, |x - a|^{n+1}. 
        $$ 
        }
}
\end{center}
With this error bound, which is derived from knowledge of $f(x)$ near
$x = a$, we can determine quite precisely how many decimal places of
accuracy a Taylor polynomial approximation achieves.  The following
example illustrates the different versions of Taylor's theorem.


\medskip
\noindent
{\bf Example}.  Consider $\sqrt{x}$ near $x = 100$.  The second degree Taylor polynomial for $\sqrt{x}$, centered at $x = 100$, is
        $$
        Q(x) = 10 + \dfrac{(x - 100)}{20} - \dfrac{(x - 100)^2}{8000}.
        $$

\vspace{105pt}
\special{psfile=../ps/calc2/j104.ps}
\special{psfile=../ps/calc2/j105.ps}

\noindent
Plot $y = Q(x)$ and $y = \sqrt{x}$ together; 
%
\mar{Version 1:\\ the remainder is $O((x-100)^3)$}
%
the result should look like the figure on the left, above.  Then plot
the remainder $y = \sqrt{x} - Q(x)$ near $x = 100$.  This graph should
suggest that $\sqrt{x} - Q(x) = O((x-100)^3)$ as $x \to 100$.  In
fact, this is what version 1 of Taylor's theorem asserts.
Furthermore,
        $$
        \lim_{x \to 100} \dfrac{\sqrt{x} - Q(x)}{(x - 100)^3} 
                \approx 6.25 \times 10^{-7};
        $$
check this yourself by constructing a table of values.  Thus
        $$
        \sqrt{x} - Q(x) \approx C (x - 100)^3 \quad 
                \mbox{where $C \approx 6.25 \times 10^{-7}$}.
        $$ 

We can use the Lagrange form of the remainder  
%
\mar{Version 2:\\ determining $C$ in terms of $\sqrt{x}$ at $x = 100$}    
%
(in version 2 of Taylor's theorem) to get the value of $C$ another
way---directly from the third derivative of $\sqrt{x}$ at $x = 100$:
\[
C = \left. \dfrac{(x^{1/2})'''}{3!} \, \right|_{x = 100} = 
 \dfrac{\tfrac{1}{2} \cdot \tfrac{-1}{2} \cdot \tfrac{-3}{2} 
                     \cdot (100)^{-5/2}}{6} = 
   \dfrac{1}{2^4 \cdot 10^5} = 6.25 \times 10^{-7}.
\]
This is the {\em exact\/} value, confirming the estimate obtained above.

Let's see what the equation  
%
\mar{Accuracy of\\ the polynomial approximation}
%
$\sqrt{x} - Q(x) \approx 6.25 \times 10^{-7} (x - 100)^3$ tells us
about the accuracy of the polynomial approximation.  If we assume $|x
- 100| < .5 \times 10^{-k}$, then
\begin{align*}
|\sqrt{x} - Q(x)| 
  &< 6.25 \times 10^{-7} \times  (.5 \times 10^{-k})^3 \\
  &= .78125 \times 10^{-(3k+7)} \ < \ .5 \times 10^{-(3k+6)}.
\end{align*}
Thus
        \begin{center}
        $x = 100$ to $k$ decimal places 
        $\Longrightarrow$
        $\sqrt{x} = Q(x)$ to $3k + 6$ places.
        \end{center}
For example, if $x = 100.47$, then $k = 0$, so $Q(100.47) =
\sqrt{100.47}$ to 6 decimal places.  We find
\[
Q(100.47) = \underline{10.023472}3875, 
\]
and the underlined digits should be correct.  In fact, 
\[
\sqrt{100.47} =  \underline{10.023472}4521\ldots .
\]
Here is a second example.  If $x = 102.98$, then we can take $k = -1$,
so $Q(102.98) = \sqrt{102.98}$ to $3 (-1) + 6 = 3$ decimal places.  We
find
        $$
        Q(102.98) = \underline{10.147}88995, \qquad
                \sqrt{102.98} = \underline{10.147}906187\ldots .
        $$

Let's see what additional light version 3 
%
\mar{Version 3: \\ an explicit \\ error bound}  
%
sheds on our investigation.  Suppose we assume $x = 100$ to $k = 0$
decimal places.  This means that $x$ lies in the open interval $(99.5,
100.5)$.  Version 3 requires that we have a bound on the size of the
third derivative of $f(x) = \sqrt{x}$ over this interval.  Now
$f'''(x) = \tfrac{3}{8} x^{-5/2}$, and this is a decreasing function.
(Check its graph; alternatively, note that its derivative is
negative.)  Its maximum value therefore occurs at the left endpoint of
the (closed) interval $[99.5, 100.5]$:
        $$
        |f'''(x)| \leq f'''(99.5) = \tfrac{3}{8} \, (99.5)^{-5/2}
                < 3.8 \times 10^{-6}.
        $$
Therefore, from version 3 of Taylor's theorem,
        $$
        |\sqrt{x} - Q(x)| < \dfrac{3.8 \times 10^{-6}}{3!} \, |x - 100|^3
        $$
Since $|x - 100| < .5$, $|x - 100|^3 < .125$, so
        $$
        |\sqrt{x} - Q(x)| < \dfrac{3.8 \times 10^{-6} \times .125}{6} = 
                .791667 \times 10^{-7} < .5 \times 10^{-6}.
        $$
This proves $\sqrt{x} = Q(x)$ to 6 decimal places---confirming what we
found earlier.


\newpage


\subsection{Applications}

{\bf Evaluating Functions.}  An obvious use 
%
\mar{Now you can do anything your calculator can!}
%
of Taylor polynomials is to evaluate functions.  In fact, whenever you
ask a calculator or computer to evaluate a function---trigonometric,
exponential, logarithmic---it is typically giving you the value of an
appropriate polynomial (though not necessarily a Taylor polynomial).

\medskip
\noindent {\bf Evaluating Integrals.} The fundamental theorem of
calculus gives us a quick way of evaluating a definite integral provided
we can find an antiderivative for the function under the integral (cf.
chapter~6.4).  Unfortunately, many common functions, like $e^{-x^2}$
or $(\sin{x} )/x$, don't have antiderivatives that can be expressed as
finite algebraic combinations of the basic functions.  Up until now,
whenever we encountered such a function we had to rely on a Riemann sum
to estimate the integral.  But now we have Taylor polynomials, and it's
easy to find an antiderivative for a polynomial!  Thus, if we have an
awkward definite integral to evaluate, it is reasonable to expect that
we can estimate it by first getting a good polynomial approximation to
the integrand, and then integrating this polynomial.  As an example,
consider the {\bf error function}\index{error function}, erf$(t)$,
\mar{\vspace{20pt}The error function} defined by
\[
\mbox{erf$(t)$} = \dfrac{2}{\sqrt{\pi}}\int_0^t e^{-x^2} \, dx\, .
\]
This is perhaps the most important integral in statistics.  It is the
basis of the so-called ``normal distribution" and is widely used to
decide how good certain statistical estimates are.  It is important to
have a way of obtaining fast, accurate approximations for erf$(t)$.
We have already seen that
        $$
        e^{-x^2} \approx 1 - x^2 + \dfrac{x^4}{2!} - 
                \dfrac{x^6}{3!} + \dfrac{x^8}{4!} - \cdots 
               \pm \dfrac{x^{2n}}{n!}.
        $$
Now, if we antidifferentiate term by term:
\begin{align*}
\int e^{-x^2}\, dx 
  &\approx \int \left( 1 - x^2 + \dfrac{x^4}{2!} - \dfrac{x^6}{3!} 
     + \dfrac{x^8}{4!} - \cdots \pm \dfrac{x^{2n}}{n!}\right) \, dx \\[3pt]
  &= \int 1 \, dx - \int x^2 \, dx + \int \dfrac{x^4}{2!} \, dx 
     - \int \dfrac{x^6}{3!} \, dx + \cdots  
     \pm \int \dfrac{x^{2n}}{n!} \, dx\\[3pt]
  &= x - \dfrac{x^3}{3} + \dfrac{x^5}{5 \cdot 2!} - \dfrac{x^7}{7 \cdot 3!} 
     + \cdots \pm \dfrac{x^{2n+1}}{(2n+1)\cdot n!}.
\end{align*}
Thus,
        $$
        \int_0^t e^{-x^2} \, dx\, \approx \left.
                x - \dfrac{x^3}{3} + \dfrac{x^5}{5 \cdot 2!} -
                \dfrac{x^7}{7 \cdot 3!} + \cdots\ 
                \pm \dfrac{x^{2n+1}}{(2n+1) \cdot n!}\, \right|_0^t,
        $$
giving us, finally, 
%
\mar{\vspace{16pt}A formula for approximating the \\error function}
%
an approximate formula for erf$(t)$:
\[
\mbox{erf$(t)$} \approx \dfrac{2}{\sqrt{\pi}} 
     \left(t - \dfrac{t^3}{3} + \dfrac{t^5}{5 \cdot 2!} -
     \dfrac{t^7}{\cdot 3!} + \cdots \pm \dfrac{t^{2n+1}}{(2n+1) \cdot n!}
      \right).
\]
Thus if we needed to know, say, erf(1), we could quickly approximate
it.  For instance, letting $n=6$, we have
\begin{align*}
\mbox{erf$(1)$}
   &\approx  \dfrac{2}{\sqrt{\pi}} 
        \left(1 - \dfrac{1}{3} + \dfrac{1}{5 \cdot 2!} -
        \dfrac{1}{7 \cdot 3!} + \dfrac{1}{9 \cdot 4!} 
        - \dfrac{1}{11 \cdot 5!} + \dfrac{1}{13 \cdot 6!}\right)\\[3pt]
   &\approx \dfrac{2}{\sqrt{\pi}} \left(1 - \dfrac{1}{3} + 
          \dfrac{1}{10} - \dfrac{1}{42} + \dfrac{1}{216} - 
          \dfrac{1}{1320} + \dfrac{1}{9360}\right) \\[3pt]
   &\approx .746836\,\dfrac{2}{\sqrt{\pi}} \approx .842714,
\end{align*}
a value accurate to 4 decimals.  If we had needed greater accuracy, we
could simply have taken a larger value for $n$.  For instance, if we
take $n=12$, we get the estimate .8427007929\ldots , where all 10
decimals are accurate (i.e., they don't change as we take larger
values~$n$).

\vspace*{14pt}
\noindent
{\bf Evaluating Limits.}  Our final application of Taylor polynomials
makes explicit use of the order of magnitude of the remainder.
Consider the problem of evaluating a limit like
        $$
        \lim_{x \to 0} \dfrac{1 - \cos(x)}{x^2}.
        $$
Since both numerator and denominator approach $0$ as $x \to 0$, it
isn't clear what the quotient is doing.  If we replace $\cos(x)$ by
its third degree Taylor polynomial with remainder, though, we get
        $$
        \cos(x) = 1 - \dfrac{1}{2!} x^2 + R(x),
        $$
and $R(x) = O(x^4)$ as $x \to 0$.  Consequently, if $x \neq 0$ but $x
\to 0$, then
\begin{align*}
\dfrac{1 - \cos(x)}{x^2} 
  &= \dfrac{1-\left(1 - \tfrac{1}{2}x^2 + R(x)\right)}{x^2}\\[3pt]
  &= \dfrac{\tfrac{1}{2}x^2 - R(x)}{x^2} 
      = \dfrac{1}{2} - \dfrac{R(x)}{x^2}.
\end{align*}
Since $R(x) = O(x^4)$, we know that there is some constant $C$ for
which $R(x)/x^4 \to C$ as $x \to 0$.  Therefore,
\begin{align*}
 \lim_{x \to 0}\dfrac{1 - \cos(x)}{x^2}
   &= \dfrac{1}{2} - \lim_{x \to 0}\dfrac{R(x)}{x^2}
        = \dfrac{1}{2} - \lim_{x \to 0}\dfrac{x^2 \cdot R(x)}{x^4} \\[3pt]
   &= \frac{1}{2} -
        \lim_{x \to 0} x^2 \cdot \lim_{x \to 0}\dfrac{R(x)}{x^4} 
        = \frac{1}{2} -  0 \cdot C 
        = \frac{1}{2}.
\end{align*}

There is a way to shorten        
%
\mar{Extending the\\ `big oh' notation}    
%
these calculations---and to make them more transparent---by extending
the way we read the `big oh' notation.  Specifically, we will read
$O(q(x))$ as ``some (unspecified) function that is the same order of
magnitude as $q(x)$''.

Then, instead of writing $\cos(x) = 1 - \tfrac{1}{2} x^2 +
R(x)$, and then noting $R(x) = O(x^4)$ as $x \to 0$, we'll just write
        $$
        \cos(x) = 1 - \tfrac{1}{2} x^2 + O(x^4) \qquad (x \to 0).
        $$
In this spirit,
\begin{align*}
\dfrac{1 - \cos(x)}{x^2} 
  &= \dfrac{1-\left(1-\tfrac{1}{2}x^2 + O(x^4)\right)}{x^2}\\[3pt]
  &= \dfrac{\tfrac{1}{2}x^2 - O(x^4)}{x^2}
     =  \tfrac{1}{2} +O(x^2) \qquad (x \to 0).
\end{align*}
We have used the fact that $\pm O(x^4)/x^2 = O(x^2)$.  Finally, since
$O(x^2) \to 0$ as $x \to 0$ (do you see why?), the limit of the last
expression is just 1/2 as $x \to 0$.  Thus, once again we arrive at
the result
        $$
        \lim_{x \to 0} \dfrac{1 - \cos(x)}{x^2} = \dfrac{1}{2}.
        $$

\subsection{Exercises}


\num Find a seventh degree Taylor polynomial centered at $x=0$ for the
indicated antiderivatives.

\medskip

\sub $\dint \dfrac{\sin(x)}{x} \, dx$.

\ans{$\dint \dfrac{\sin(x)}{x} \, dx \approx x - \dfrac{x^3}{3 \cdot 3!} +
\dfrac{x^5}{5 \cdot 5!} - \dfrac{x^7}{7 \cdot 7!}$.}

\sub $\dint e^{x^2} \, dx$.
\sub $\dint \sin(x^2) \, dx$.

\num Plot the 7-th degree polynomial you found in part (a) above over
the interval $[0, 5]$.  Now plot the 9-th degree approximation on the
same graph.  When do the two polynomials begin to differ visibly?


\num  Using the seventh degree Taylor approximation 
        $$
        E(t) \approx \int_0^t e^{-x^2} \, dx\, 
                = t - \dfrac{t^3}{3} + \dfrac{t^5}{5 \cdot 2!} 
                - \dfrac{t^7}{7 \cdot 3!},
        $$
calculate the values of $E(.3)$ and $E(-1)$.  Give only the
significant digits---that is, report only those decimals of your
estimates that you think are fixed.  (This means you will also need to
calculate the ninth degree Taylor polynomial as well---do you see
why?)

\num Calculate the values of $\sin(.4)$ and $\sin(\pi/12)$ using the
seventh degree Taylor polynomial centered at $x=0$
        $$
        \sin(x) \approx x - \dfrac{x^3}{3!} + \dfrac{x^5}{5!} 
                -  \dfrac{x^7}{7!}.
        $$
Compare your answers with what a calculator gives you.

\num Find the third degree Taylor polynomial for $g(x) = x^3 - 3x$ at
$x=1$.  Show that the Taylor polynomial is actually equal to
$g(x)$---that is, the remainder is 0.  What does this imply about the
     {\em fourth} degree Taylor polynomial for $g$ at $x=1$ ?

\num Find the seventh degree Taylor polynomial centered at $x=\pi$
for\\ (a) $\sin(x)$; \quad (b) $\cos(x)$; \quad (c) $\sin(3x)$.

\num In this problem you will compare computations using Taylor
polynomials centered at $x=\pi$ with computations using Taylor
polynomials centered at $x=0$.

\sub Calculate the value of $\sin(3)$ using a seventh degree Taylor
polynomial centered at $x=0$.  How many decimal places of your estimate
appear to be fixed?

\sub Now calculate the value of $\sin(3)$ using a seventh degree
Taylor polynomial centered at $x=\pi$.  Now how many decimal places of
your estimate appear to be fixed?


\num Write a program which evaluates a Taylor polynomial to print out
$\sin(5^\circ)$, $\sin(10^\circ)$, $\sin(15^\circ)$, \ldots,
$\sin(40^\circ)$, $\sin(45^\circ)$ accurate to 7 decimals.  (Remember
to convert to radians before evaluating the polynomial!)

\num {\bf Why \boldmath $0!=1$}. \quad When you were first introduced to
exponential notation in expressions like $2^n$, $n$ was restricted to
being a positive integer, and $2^n$ was defined to be the product of 2
multiplied by itself $n$ times.  Before long, though, you were working
with expressions like $2^{-3}$ and $2^{1/4}$.  These new expressions
weren't defined in terms of the original definition.  For instance, to
calculate $2^{-3}$ you wouldn't try to multiply 2 by itself $-3$
times---that would be nonsense!  Instead, $2^{-m}$ is defined by
looking at the key {\em properties} of exponentiation for positive
exponents, and extending the definition to other exponents in a way
that preserves these properties.  In this case, there are two such
properties, one for adding exponents and one for multiplying them:
\begin{alignat*}{3}
\text{Property A:}& &\qquad 2^m \cdot 2^n &= 2^{m+n} &\quad 
         &\text{for all positive $m$ and $n$}, \\
\text{Property M:}& &\qquad \left(2^m\right)^n &= 2^{mn} & 
         &\text{for all positive $m$ and $n$}. \\
\end{alignat*}

\sub Show that to preserve property A we have to define $2^0=1$.

\sub Show that we then have to define $2^{-3}=1/2^3$ if we are to
continue to preserve property A.

\sub Show why $2^{1/4}$ must be $\sqrt[4]{2}$.

\sub In the same way, you should convince yourself that a basic
property of the factorial notation is that $(n+1)!=(n+1)\cdot n!$ for
any positive integer~$n$.  Then show that to preserve this property,
we have to define $0!=1$.

\sub Show that there is no way to define $(-1)!$ which preserves this
property.

\num Use the general rule to derive the 5-th degree Taylor polynomial
centered at $x=0$ for the function
\[
f(x) = (1+x)^{\frac{1}{2}}.
\]
Use this approximation to estimate $\sqrt{1.1}$. How accurate is
this?

\newpage

\num Use the general rule to derive the formula for the $n$-th degree
Taylor polynomial centered at $x=0$ for the function
\[
f(x) = (1+x)^c \, \mbox{ where $c$ is a constant}.
\]

\num Use the result of the preceding problem to get the 6-th degree
Taylor polynomial centered at $x=0$ for $1/\sqrt[3]{1+x^2}$.

\ans{$1 - \dfrac{1}{3}x^2 + \dfrac{2}{9}x^4 - \dfrac{14}{81}x^6$.}

\num Use the result of the preceding problem to approximate
\[
\int_0^1 \dfrac{1}{\sqrt[3]{1+x^2}}\, dx.
\]

\num Calculate the first 7 decimals of erf$(.3)$. Be sure to show why
you think all 7 decimals are correct.  What degree Taylor polynomial
did you need to produce these 7 decimals?

\ans{erf$(.3)=.3286267\ldots .$}

\numa Apply the general formula for calculating Taylor polynomials
centered at $x=0$ to the tangent function to get the 5-th degree
approximation.

\ans{$\tan(x)\approx x+x^3/3+2x^5/15$.}

\sub Recall that $\tan(x) =\sin(x)/\cos(x)$.  Multiply the 5-th degree
Taylor polynomial for $\tan(x)$ from part a) by the 4-th degree Taylor
polynomial for $\cos(x)$ and show that you get the fifth degree
polynomial for $\sin(x)$ (discarding higher degree terms).

\num Show that the $n$-th degree Taylor polynomial centered at $x=0$
for $1/(1-x)$ is $1+x+x^2+\cdots+x^n$.

\num Note that
\[
\int \frac{1}{1-x}\, dx = -\ln(1-x).
\]
Use this observation, together with the result of the previous
problem, to get the $n$-th degree Taylor polynomial centered at $x=0$
for $\ln(1-x)$.

\numa Find a formula for the $n$-th degree Taylor polynomial centered
at $x=1$ for $\ln(x)$.

\sub Compare your answer to part (a) with the Taylor polynomial
centered at $x=0$ for $\ln(1-x)$ you found in the previous problem.
Are your results consistent?

\numa The first degree Taylor polynomial for $e^x$ at $x=0$ is $1+x$.
Plot the remainder $R_1(x) = e^x - (1+x)$ over the interval $-.1 \leq
x \leq .1$.  How does this graph demonstrate that $R_1(x) = O(x^2)$ as
$x \rightarrow 0$?

\sub There is a constant $C_2$ for which $R_1(x) \approx C_2 x^2$ when
$x \approx 0$.  Why?  Estimate the value of $C_2$.

\num This concerns the second degree Taylor polynomial for $e^x$ at
$x=0$.  Plot the remainder $R_2(x) = e^x - (1 + x + x^2/2)$ over the
interval $-.1 \leq x \leq .1$.  How does this graph demonstrate that
$R_2(x) = O(x^3)$ as $x \rightarrow 0$?

\sub There is a constant $C_3$ for which $R_2(x) \approx C_3 x^3$ when
$x \approx 0$.  Why?  Estimate the value of $C_3$.

\num Let $R_3(x) = e^x - P_3(x)$, where $P_3(x)$ is the third degree
Taylor polynomial for $e^x$ at $x=0$.  Show $R_3(x) = O(x^4)$ as $x
\rightarrow 0$.

\num At first glance, Taylor's theorem says that
$$
\sin(x) = x - \dfrac{1}{6} x^3 + O(x^4) \quad \mbox{as } x \rightarrow 0.
$$ 
However, graphs and calculations done in the text
(pages~\pageref{startsine}--\pageref{endsine}) make it clear that
$$
\sin(x) = x - \dfrac{1}{6} x^3 + O(x^5) \quad \mbox{as } x \rightarrow 0.
$$
Explain this.  Is Taylor's theorem wrong here?

\num  Using a suitable formula (that is, a Taylor polynomial with
remainder) for each of the functions involved, find the indicated limit.

\sub ${\displaystyle \lim_{x \to 0} \dfrac{\sin(x)}{x}}$ \hfill [Answer: 1]

\sub ${\displaystyle \lim_{x \to 0} \dfrac{e^x - (1 + x)}{x^2}}$ \hfill
[Answer: 1/2]

\sub ${\displaystyle \lim_{x \to 1} \dfrac{\ln{x}}{x - 1}}$

\sub ${\displaystyle \lim_{x \to 0} \dfrac{x - \sin(x)}{x^3}}$ \hfill
[Answer: 1/6]

\sub ${\displaystyle \lim_{x \to 0} \dfrac{\sin(x^2)}{1 - \cos(x)}}$

\num Suppose $f(x) = 1 + x^2 + O(x^4)$ as $x \rightarrow 0$.  Show
that
$$
(f(x))^2 = 1 + 2x^2 + O(x^4) \quad \mbox{as } x \rightarrow 0.
$$

\numa Using $\sin x = x - \tfrac{1}{6} x^3 + O(x^5)$ as $x \rightarrow
0$, show
$$
(\sin x)^2 = x^2 - \tfrac{1}{3} x^4 + O(x^6) \quad \mbox{as } x \rightarrow 0.
$$

\sub Using $\cos x = 1 - \tfrac{1}{2} x^2 + \tfrac{1}{24} x^4 +
O(x^5)$ as $x \rightarrow 0$, show
$$
(\cos x)^2 = 1 - x^2 + \tfrac{1}{3} x^4 + O(x^5) \quad \mbox{as } 
x \rightarrow 0.
$$

\sub Using the previous parts, show $(\sin x)^2 + (\cos x)^2 = 1 +
O(x^5)$ as $x \rightarrow 0$.  (Of course, you already know $(\sin
x)^2 + (\cos x)^2 = 1$ {\em exactly}.)

\numa Apply the general formula for calculating Taylor polynomials to
the tangent function to get the 5-th degree approximation.

% ans:  $\tan(x) \approx x + x^3/3 + 2x^5/15$

\sub Recall that $\tan(x) = \sin(x)/\cos(x)$, so $\tan(x) \cdot
\cos(x) = \sin(x)$.  Multiply the fifth degree Taylor polynomial for
$\tan(x)$ from part a) by the fifth degree Taylor polynomial for
$\cos(x)$ and show that you get the fifth degree Taylor polynomial for
$\sin(x)$ plus $O(x^6)$---that is, plus terms of order 6 and higher.

\numa  Using the formulas
\begin{align*}
e^u &= 1 + u + \tfrac{1}{2}u^2  + 
                \tfrac{1}{6} u^3 + O(u^4) \qquad (u \to 0), \\[4pt]
\sin{x} &= x - \tfrac{1}{6} x^3 + O(x^5) \qquad (x \to 0),
\end{align*}
show that $e^{\sin{x}} = 1 + x + \tfrac{1}{2} x^2 + O(x^4)$ as              $x \to 0$.

\sub Apply the general formula to obtain the third degree Taylor
polynomial for $e^{\sin{x}}$ at $x = 0$, and compare your result with
the formula in part (a).

\num Using $\dfrac{e^x -1}{x} = 1 + \tfrac{1}{2} x + \tfrac{1}{6} x^2 +
\tfrac{1}{24} x^3 + O(x^4)$ as $x \to 0$, show that
\[
\dfrac{x}{e^x - 1} = 1 - \tfrac{1}{2} x + 
    \tfrac{1}{12} x^2 + O(x^4) \qquad (x \to 0).
\]

\num  Show that the following are true as $x \rightarrow \infty$.

\sub $x + 1/x = O(x)$.

\sub  $5x^7 - 12x^4 + 9 = O(x^7)$.

\sub  $\sqrt{1 + x^2} = O(x)$.

\sub  $\sqrt{1 + x^p} = O(x^{p/2})$.

\numa  Let $f(x) = \ln(x)$.  Find the smallest bound $M$ for which
        $$
        |f^{(4)}(x)| \leq M \quad \mbox{when $|x - 1| \leq .5$}.
        $$

\sub Let $P_3(x)$ be the degree 3 Taylor polynomial for $\ln(x)$ at $x
= 1$, and let $R_3(x)$ be the remainder $R_3(x) = \ln(x) - P_3(x)$.
Find a number $K$ for which
        $$
        |R(x)| \leq K \, |x - 1|^4 
        $$
for all $x$ satisfying $|x - 1| \leq .5$.

\sub If you use $P_3(x)$ to approximate the value of $\ln(x)$ in the
interval $.5 \leq x \leq 1.5$, how many digits of the approximation
are correct?

\sub Suppose we restrict the interval to $|x - 1| \leq .1$.  Repeat
parts (a) and (b), getting {\em smaller\/} values for $M$ and $K$.
Now how many digits of the polynomial approximation $P_3(x)$ to
$\ln(x)$ are correct, if $.9 \leq x \leq 1.1$?

\bigskip
\noindent
{\bf ``Little oh" notation}.  Similar to the ``big oh" notation is
another, called the ``little oh": if
$$
\lim_{x \rightarrow a} \dfrac{\phi(x)}{q(x)} = 0,
$$
then we write $\phi(x) = o(q(x))$ and say $\phi$ is `{\em little oh}'
of $q$ as $x \rightarrow a$.

\num  Suppose $\phi(x) = O(x^6)$ as $x \rightarrow 0$.  Show the following.

\sub  $\phi(x) = O(x^5)$ as $x \rightarrow 0$.

\sub  $\phi(x) = o(x^5)$ as $x \rightarrow 0$.

\sub It is false that $\phi(x) = O(x^7)$ as $x \rightarrow 0$.  (One
way you can do this is to give an explicit example of a function
$\phi(x)$ for which $\phi(x) = O(x^6)$ but for which you can show
$\phi(x) = O(x^7)$ is false.)

\sub  It is false that $\phi(x) = o(x^6)$ as $x \rightarrow 0$.

\num Sketch the graph $y = x \ln(x)$ over the interval $0 < x \leq 1$.
Explain why your graph shows $\ln(x) = o(1/x)$ as $x \rightarrow 0$.

\newpage
%%%%%%%%%%%%%%%%%%%%%%%%%%%%%%%%%%%%%%%%%%%%%%%%%%%%%%%%%%%%%%%%%%%%%%%%
%                                                                      %
%                               Section 3                              %
%                                                                      %
%%%%%%%%%%%%%%%%%%%%%%%%%%%%%%%%%%%%%%%%%%%%%%%%%%%%%%%%%%%%%%%%%%%%%%%%


\section{Taylor Series}

In the previous section we have been talking about approximations to
functions by their Taylor polynomials.  Thus, for instance, we were
able to write statements like
\[
\sin(x) \approx x - \dfrac{x^3}{3!} + \dfrac{x^5}{5!} - \dfrac{x^7}{7!},
\]
where the approximation was a good one for values of $x$ not too far
from 0.  On the other hand, when we looked at Taylor polynomials of
higher and higher degree, the approximations were good for larger and
larger values of $x$. We are thus tempted to write
\[
\sin(x) = x - \dfrac{x^3}{3!} + \dfrac{x^5}{5!} 
                - \dfrac{x^7}{7!} + \dfrac{x^9}{9!} - \cdots ,
\]
indicating that the sine function is equal to this ``infinite degree"
polynomial.  This infinite sum is called the {\bf Taylor
  series}\index{Taylor series!definition} centered at $x=0$ for
$\sin(x)$.  But what would we even mean by such an infinite sum?
%
\mar {You have seen\\ infinite sums before}
%
We will explore this question in detail in section~5, but you should
already have some intuition about what it means, for it can be
interpreted in exactly the same way we interpret a more familiar
statement like
\begin{align*}
\dfrac{1}{3} 
    &= .33333 \ldots \\
    &= \dfrac{3}{10} + \dfrac{3}{100} + \dfrac{3}{1000}
           + \dfrac{3}{10000} + \dfrac{3}{100000} + \cdots .
\end{align*}
Every decimal number is a sum of fractions whose denominators are
powers of 10; 1/3 is a number whose decimal expansion happens to need
an infinite number of terms to be completely precise.  Of course, when
a practical matter arises (for example, typing a number like 1/3 or
$\pi$ into a computer) just the beginning of the sum is used---the
``tail'' is dropped.  We might write 1/3 as 0.33, or as 0.33333, or
however many terms we need to get the accuracy we want.  Put another
way, we are saying that 1/3 is the {\em limit} of the finite sums of
the right hand side of the equation.

Our new formulas for Taylor series 
%
\mar{Infinite degree polynomials are to be viewed like infinite decimals}
%
are meant to be used exactly the same way: when a computation is
involved, take only the beginning of the sum, and drop the tail.  Just
where you cut off the tail depends on the input value $x$ and on the
level of accuracy needed.  Look at what happens when we we approximate
the value of $\cos(\pi/3)$ by evaluating Taylor polynomials of
increasingly higher degree: 

\newpage

\mbox{}\vspace{-18pt}
\small
\begin{align*}
1 \quad &= 1.0000000 \\[3pt]
1 - \dfrac{1}{2!} \left( \dfrac{\pi}{3} \right)^2 
        &\approx 0.4516887 \\
1 - \dfrac{1}{2!} \left( \dfrac{\pi}{3} \right)^2 
                + \dfrac{1}{4!} \left( \dfrac{\pi}{3} \right)^4
        &\approx  0.5017962 \\
1 - \dfrac{1}{2!} \left( \dfrac{\pi}{3} \right)^2 
                + \dfrac{1}{4!} \left( \dfrac{\pi}{3} \right)^4
                - \dfrac{1}{6!} \left( \dfrac{\pi}{3} \right)^6
        &\approx  0.4999646 \\
1 - \dfrac{1}{2!} \left( \dfrac{\pi}{3} \right)^2 
                + \dfrac{1}{4!} \left( \dfrac{\pi}{3} \right)^4
                - \dfrac{1}{6!} \left( \dfrac{\pi}{3} \right)^6
                + \dfrac{1}{8!} \left( \dfrac{\pi}{3} \right)^8
        &\approx  0.5000004 \\
1 - \dfrac{1}{2!} \left( \dfrac{\pi}{3} \right)^2 
                + \dfrac{1}{4!} \left( \dfrac{\pi}{3} \right)^4
                - \dfrac{1}{6!} \left( \dfrac{\pi}{3} \right)^6
                + \dfrac{1}{8!} \left( \dfrac{\pi}{3} \right)^8
                - \dfrac{1}{10!} \left( \dfrac{\pi}{3} \right)^{10}
        &\approx  0.5000000 
\end{align*}
\normalsize

These sums were evaluated by setting $\pi = 3.141593$.  As you can
see, at the level of precision we are using, a sum that is six terms
long gives the correct value.  However, five, four, or even three
terms may have been adequate for the needs at hand.  The crucial fact
is that these are all honest calculations using {\em only\/} the four
operations of elementary arithmetic.

Note that if we had wanted to get the same 6 place accuracy for
$\cos(x)$ for a larger value of $x$, we might need to go further out
in the series.  For instance $\cos(7\pi/3)$ is also equal to .5, but
the tenth degree Taylor polynomial centered at $x=0$ gives \small
\[
1 - \dfrac{1}{2!} \left( \dfrac{7\pi}{3} \right)^2 
                + \dfrac{1}{4!} \left( \dfrac{7\pi}{3} \right)^4
                - \dfrac{1}{6!} \left( \dfrac{7\pi}{3} \right)^6
                + \dfrac{1}{8!} \left( \dfrac{7\pi}{3} \right)^8
                - \dfrac{1}{10!} \left( \dfrac{7\pi}{3} \right)^{10}
                = -37.7302,
\]
\normalsize% 
which is not even close to .5 .  In fact, to get $\cos(7\pi/3)$ to 6
decimals, we need to use the Taylor polynomial centered at $x=0$ of
degree 30, while to get $\cos(19\pi/3)$ (also equal to .5) to 6
decimals we need the Taylor polynomial centered at $x=0$ of degree 66!

The key fact, though, is that, for any value of $x$, if we go out in
the series far enough (where what constitutes ``far enough" will
depend on $x$), we can approximate $\cos(x)$ to any number of decimal
places desired. For any $x$, the value of $\cos(x)$ is the limit of
the finite sums of the Taylor series, just as 1/3 is the limit of the
finite sums of its infinite series representation.

In general, given a function $f(x)$, its Taylor series\index{Taylor
  series!definition} centered at $x=0$ will be
\[
f(0) + f'(0)x + \frac{f^{(2)}(0)}{2!}x^2 + \frac{f^{(3)}(0)}{3!}x^3 
   + \frac{f^{(4)}(0)}{4!}x^4 + \cdots 
   = \sum_{k=0}^{\infty} \frac{f^{(k)}(0)}{k!}x^k.
\]
We have the following Taylor series\index{Taylor series} centered at
$x=0$ for some common functions:
$$
        \begin{array}{ccl}
        f(x) &\hspace{.5in} & \hspace{.5in}\mbox{Taylor series for $f(x)$} \\ [2pt] \hline
        & & \\ [-6pt] 
        \sin(x) & &  x - \dfrac{x^3}{3!} + \dfrac{x^5}{5!} 
                - \dfrac{x^7}{7!} + \cdots \\[.15in]
        \cos(x) & & 1 - \dfrac{x^2}{2!} + \dfrac{x^4}{4!} 
                - \dfrac{x^6}{6!} + \cdots \\[.15in]
        e^x & & 1 + x + \dfrac{x^2}{2!} + \dfrac{x^3}{3!} 
                + \dfrac{x^4}{4!} + \cdots  \\[.15in]
        \ln{(1-x)}& & -\left(x + \dfrac{x^2}{2} + \dfrac{x^3}{3} 
                + \dfrac{x^4}{4} + \cdots \right)\\[.15in]
        \dfrac{1}{1-x} & &  1 + x + x^2 + x^3 + \cdots \\ [.15in]
        \dfrac{1}{1+x^2}& & 1 - x^2 + x^4 - x^6 + \cdots \\ [.15in]
        (1+x)^c & & 1 +cx+ \dfrac{c(c-1)}{2!}x^2 
           + \dfrac{c(c-1)(c-2)}{3!}x^3 + \cdots
        \end{array}
$$\index{Taylor series!table of}

While it is true that $\cos(x)$ and $e^x$ equal their Taylor series,
just as $\sin(x)$ did, we have to be more careful with the last four
functions.  To see why this is, let's graph $1/(1+x^2)$ and its Taylor
polynomials $P_n(x) = 1-x^2+ x^4-x^6+\cdots \pm x^n$ for $n=2$, 4, 6,
8, 10, 12, 14, 16, 200, and 202.  Since all the graphs are symmetric
about the $y$-axis (why is this?), we draw only the graphs for
positive~$x$:

\vspace*{2.3in}
\special{psfile=../ps/calc2/Darct.ps}\index{Taylor polynomial!graphs for $1/(1+x^2)$}

It appears that the 
%
\mar{A Taylor series\\ may not converge\\ for all values of $x$}
%
graphs of the Taylor polynomials $P_n(x)$ approach the graph of
$1/(1+x^2)$ very nicely {\em so long as $x < 1$}.  If $x\geq1$,
though, it looks like there is no convergence, no matter how far out
in the Taylor series we go.  We can thus write
\[
\dfrac{1}{1+x^2} = 1 - x^2 + x^4 - x^6 + \cdots \qquad 
        \text{for $|x| < 1$},
\]
where the restriction on $x$ is essential if we want to use the $=$
sign. We say that the interval $-1<x<1$ is the {\bf interval of
  convergence}\index{interval of convergence} for the Taylor series
centered at $x=0$ for $1/(1+x^2)$.  Some Taylor series, like those for
$\sin(x)$ and $e^x$, converge for all values of $x$---their interval
of convergence is $(-\infty, \infty)$.  Other Taylor series, like
those for $1/(1+x^2)$ and $\ln(1-x)$, have finite intervals of
convergence.

\aside{Brook Taylor (1685--1731)\index{Taylor, Brook} was
  an English mathematician who developed the series that bears his
  name in his book {\em Methodus incrementorum} (1715).  He did not
  worry about questions of convergence, but used the series freely to
  attack many kinds of problems, including differential equations.}

\noindent
{\bf Remark} On the one hand it is perhaps not too surprising that a
function should equal its Taylor series---after all, with more and
more coefficients to fiddle with, we can control more and more of the
behavior of the associated polynomials.  On the other hand, we are
saying that a function like $\sin(x)$ or $e^x$ has its behavior for
all values of $x$ completely determined by the value of the function
and all its derivatives at a single point, so perhaps it is surprising
after all!

\subsection{Exercises}

\numa Suppose you wanted to use the Taylor series centered at $x=0$ to
calculate $\sin(100)$.  How large does $n$ have to be before the term
$(100)^n/{n!}$ is less than 1?

\sub If we wanted to calculate $\sin(100)$ directly using this Taylor
series, we would have to go very far out before we began to approach a
limit at all closely.  Can you use your knowledge of the way the
circular functions behave to calculate $\sin(100)$ much more rapidly
(but still using the Taylor series centered at $x=0$)?  Do it.

\sub Show that we can calculate the sine of any number by using a
Taylor series centered at $x=0$ {\em either} for $\sin(x)$ {\em or}
for $\cos(x)$ to a suitable value 0f $x$ between 0 and $\pi/4$.

\numa Suppose we wanted to calculate $\ln{5}$ to 7 decimal places. An
obvious place to start is with the Taylor series centered at $x=0$ for
$\ln(1-x)$:
\[
-\left(x + \dfrac{x^2}{2} + \dfrac{x^3}{3} + \dfrac{x^4}{4} 
       + \cdots \right)
\]
with $x=-4$.  What happens when you do this, and why?  Try a few more
values for $x$ and see if you can make a conjecture about the interval
of convergence for this Taylor series.

\ans{The Taylor series converges for $-1\leq x < 1$.}

\sub Explain how you could use the fact that $\ln(1/A)=-\ln{A}$ for
any real number $A>0$ to evaluate $\ln{x}$ for $x>2$.  Use this to
compute $\ln{5}$ to 7 decimals.  How far out in the Taylor series did
you have to go?

\sub If you wanted to calculate $\ln 1.5$, you could use the Taylor
series for $\ln(1-x)$ with either $x=-1/2$, which would lead directly
to $\ln 1.5$, or you could use the series with $x=1/3$, which would
produce $\ln(2/3) = -\ln 1.5$\, .  Which method is faster, and why?

\num We can improve the speed of our calculations of the logarithm
function slightly by the following series of observations: \sub Find
the Taylor series centered at $u=0$ for $\ln(1+u)$.

\ans{$u-u^2/2+u^3/3-u^4/4+u^5/5+\cdots$} 

\sub Find the Taylor series centered at $u=0$ for
\[
\ln\left( \dfrac{1-u}{1+u} \right).
\]
(Remember that $\ln(A/B) = \ln{A} - \ln{B}$.)

\sub Show that {\em any} $x>0$ can be written in the form
$(1-u)/(1+u)$ for some suitable $-1<u<1$.

\sub Use the preceding to evaluate $\ln{5}$ to 7 decimal places. How
far out in the Taylor series did you have to go?

\numa Evaluate $\arctan(.5)$ to 7 decimal places.  \sub Try to use the
Taylor series centered at $x=0$ to evaluate $\arctan(2)$
directly---what happens?  Remembering what the arctangent function
means geometrically, can you figure out a way around this difficulty?

\numa {\bf Calculating $\pi$}\index{$\pi$, calculation of} \quad The
Taylor series for the arctangent function,
\[
\arctan{x} = x - \dfrac{1}{3}x^3 + \dfrac{1}{5} x^5 - \cdots 
       \pm \dfrac{1}{{2n+1}} x^{2n+1} + \cdots ,
\]
lies behind many of the methods for getting lots of decimals of $\pi$
rapidly.  For instance, since $\tan\left(\tfrac{\pi}{4}\right) = 1$, we have
$\tfrac{\pi}{4} = \arctan{1}$.  Use this to get a series expansion for
$\pi$.  How far out in the series do you have to go to evaluate $\pi$
to 3 decimal places?

%  \ans{Out past the term of degree 4911.}

\sub The reason the preceding approximations converged so slowly was
that we were substituting $x=1$ into the series, so we didn't get any
help from the $x^n$ terms in making the successive corrections get
small rapidly.  We would like to be able to do something with values
of $x$ between 0 and 1.  We can do this by using the addition formula
for the tangent function:
\[
\tan(\alpha + \beta ) = \dfrac{{\tan\alpha +\tan\beta}}
    {1 - \tan\alpha \tan\beta}.
\]
Use this to show that
\[
\dfrac{\pi}{4} 
  = \arctan{\left(\dfrac{1}{2}\right)}
    + \arctan{\left(\dfrac{1}{5}\right)} 
    + \arctan{\left(\dfrac{1}{8}\right)}.
\]
Now use the Taylor series for each of these three expressions to
calculate $\pi$ to 12 decimal places.  How far out in the series do
you have to go?  Which series did you have to go the farthest out in
before the 12th decimal stabilized? Why?

\num {\bf Raising $e$ to imaginary powers}\index{$e$, raised to
  imaginary powers} One of the major mathematical developments of the
last century was the extension of the ideas of calculus to {\bf
  complex numbers}\index{complex numbers}---i.e., numbers of the form
$r+s\,i$, where $r$ and $s$ are real numbers, and $i$ is a new symbol,
defined by the property that $i\cdot i=-1$\,.  Thus $i^3=i^2\,i=-i$,
$i^4=i^2\,i^2=(-1)(-1)=1$, and so on. If we want to extend our
standard functions to these new numbers, we proceed as we did in the
previous section and look for the crucial {\em properties} of these
functions to see what they suggest.  One of the key properties of
$e^x$ as we've now seen is that it possesses a Taylor series:
\[
e^x = 1 + x + \dfrac{x^2}{2!} + \dfrac{x^3}{3!} 
                + \dfrac{x^4}{4!} + \cdots .
\]
But this property only involves operations of ordinary arithmetic, and
so makes perfectly good sense even if $x$ is a complex number

\sub Show that if $s$ is any real number, we must define $e^{i\,s}$ to
be $\cos(s)+i\,\sin(s)$ if we want to preserve this property.

\sub Show that $e^{\pi\,i}=-1$.

\sub Show that if $r+s\,i$ is any complex number, we must have
\[
e^{r+s\,i}=e^r(\cos{s}+i\,\sin{s})
\]
if we want complex exponentials to preserve all the right properties.

\sub Find a complex number $r+si$ such that $e^{r+s\,i}= -5$.

\num {\bf Hyperbolic trigonometric functions}\index{function!hyperbolic}\index{$\cosh$}\index{$\sinh$} The hyperbolic trigonometric
functions are defined by the formulas
\[
\cosh(x) = \dfrac{e^x + e^{-x}}{2}, \qquad \quad
\sinh(x) = \dfrac{e^x - e^{-x}}{2}.
\]
(The names of these functions are usually pronounced ``cosh" and
``cinch.")  In this problem you will explore some of the reasons for
the adjectives {\em hyperbolic} and {\em trigonometric}.

\sub Modify the Taylor series centered at $x=0$ for $e^x$ to find a
Taylor series for $\cosh(x)$.  Compare your results to the Taylor
series centered at $x=0$ for $\cos(x)$.

\sub Now find the Taylor series centered at $x=0$ for $\sinh(x)$.
Compare your results to the Taylor series centered at $x=0$ for
$\sin(x)$.

\sub Parts (a) and (b) of this problem should begin to explain the
     {\em trigonometric} part of the story.  What about the {\em
       hyperbolic} part?  Recall that the familiar trigonometric
     functions are called {\em circular functions} because, for any~$t$,
     the point $(\cos t, \sin t)$ is on the unit circle with equation
     $x^2 + y^2 =1$ (cf.\ chapter~7.2).  Show that the point $(\cosh
     t, \sinh t)$ lies on the hyperbola with equation $x^2 - y^2 = 1$.

\num  Consider the Taylor series centered at $x=0$ for $(1+x)^c$.

\sub What does the series give if you let $c=2$?  Is this reasonable?

\sub What do you get if you set $c=3$?

\sub Show that if you set $c=n$, where $n$ is a positive integer, the
Taylor series will terminate. This yields a general formula---the {\bf
  binomial theorem}\index{binomial theorem}---that was discovered by
the 12th century Persian poet and mathematician, Omar
Khayyam\index{Khayyam, Omar} (c.\ 1050--1130), and generalized by
Newton to the form you have just obtained.  Write out the first three
and the last three terms of this formula.

\sub Use an appropriate substitution for $x$ and a suitable value for
$c$ to derive the Taylor series for $1/(1-u)$. Does this agree with
what we previously obtained?

\sub Suppose we want to calculate $\sqrt{17} \,$. We might try letting
$x=16$ and $c=1/2$ and using the Taylor series for $(1+x)^c$.  What
happens when you try this?

\sub We can still use the series to help us, though, if we are a
little clever and write
\[
\sqrt{17} =\sqrt{16+1} = \sqrt{16\left(1+\dfrac{1}{16}\right)} 
    =\sqrt{16}\cdot\sqrt{1+ \dfrac{1}{16}}
    = 4 \cdot \sqrt{1 + \dfrac{1}{16}}.
\]
Now apply the series using $x=1/16$ to evaluate $\sqrt{17}$ to 7
decimal place accuracy.  How many terms does it take?  \sub Use the
same kind of trick to evaluate $\sqrt[3]{30}$.

\vspace{14pt}
\noindent
{\bf Evaluating Taylor series rapidly}\index{polynomial!rapid
  evaluation of} Suppose we wanted to plot the Taylor polynomial of
degree 11 associated with $\sin(x)$.  For each value of $x$, then, we
would have to evaluate
\[
P_{11}(x)= x - \dfrac{x^3}{3!} + \dfrac{x^5}{5!} 
           - \dfrac{x^7}{7!} + \dfrac{x^9}{9!} - \dfrac{x^{11}}{11!}.
\]
Since the length of time it takes the computer to evaluate an
expression like this is roughly proportional to the number of
multiplications and divisions involved (additions and subtractions, by
comparison, take a negligible amount of time), let's see how many of
these operations are needed to evaluate $P_{11}(x)$.  To calculate
$x^{11}$ requires 10 multiplications, while 11! requires 9 (if we are
clever and don't bother to multiply by 1 at the end!), so the
evaluation of the term $x^{11}/11!$ will require a total of 20
operations (counting the final division).  Similarly, evaluating
$x^9/9!$ requires 16 operations, $x^7/7!$ requires 12, on down to
$x^3/3!$, which requires 4.  Thus the total number of multiplications
and divisions needed is
\[
4 + 8 + 12 + 16 + 20 = 60.
\]
This is not too bad, although if we were doing this for many different
values of $x$, which would be the case if we wanted to graph
$P_{11}(x)$, this would begin to add up.  Suppose, though, that we
wanted to graph something like $P_{51}(x)$ or $P_{101}(x)$.  By the
same analysis, evaluating $P_{51}(x)$ for a single value of $x$ would
require
\[
4 + 8 + 12 + 16 + 20 + 24 + \cdots + 96 + 100 = 1300
\]
multiplications and divisions, while evaluation of $P_{101}(x)$ would
require 5100 operations.  Thus it would take roughly 20 times as long
to evaluate $P_{51}(x)$ as it takes to evaluate $P_{11}(x)$, while
$P_{101}(x)$ would take about 85 times as long.

\num Show that, in general, the number of multiplications and
divisions needed to evaluate $P_n(x)$ is roughly $n^2/2$.

\vspace{14pt}
We can be clever, though.  Note that $P_{11}(x)$ can be written as
\[
x \left(1 - \dfrac{x^2}{2 \cdot 3} \left(1 - \dfrac{x^2}{4 \cdot 5}
  \left(1 - \dfrac{x^2}{6 \cdot 7} \left(1 - \dfrac{x^2}{8 \cdot 9}
  \left(1 - \dfrac{x^2}{10 \cdot 11}\right) \right) \right) \right) \right).
\]
 
\num How many multiplications and divisions are required to evaluate
this expression?  \ans{$3+4+4+4+4+1=20$.}

\num Thus this way of evaluating $P_{11}(x)$ is roughly three times as
fast, a modest saving.  How much faster is it if we use this method to
evaluate $P_{51}(x)$?

\ans{The old way takes roughly 13 times as long.}

\num Find a general formula for the number of multiplications and
divisions needed to evaluate $P_n(x)$ using this way of grouping.

\medskip

Finally, we can extend these ideas to reduce the number of operations
even further, so that evaluating a polynomial of degree $n$ requires
only $n$ multiplications, as follows.  Suppose we start with a
polynomial
$$
p(x) = a_0 +a_1\,x+a_2\,x^2+a_3\,x^3 + \cdots +a_{n-1}\,x^{n-1}+a_n\,x^n.
$$
We can rewrite this as
$$
a_0+x\left(a_1+x\left(a_2+\ldots 
   +x\left(a_{n-2}+x\left(a_{n-1}+a_n\,x\right)\right)\ldots\right)\right).
$$
You should check that with this representation it requires only $n$
multiplications to evaluate $p(x)$ for a given $x$.

\numa Write two computer programs to evaluate the 300th degree Taylor
polynomial centered at $x=0$ for $e^x$, with one of the programs being
the obvious, standard way, and the second program being this method
given above. Evaluate $e^1=e$ using each program, and compare the
length of time required.  \sub Use these two programs to graph the
300th degree Taylor polynomial for $e^x$ over the interval $[0, 2]$,
and compare times.

\newpage
%%%%%%%%%%%%%%%%%%%%%%%%%%%%%%%%%%%%%%%%%%%%%%%%%%%%%%%%%%%%%%%%%%%%%%%%
%                                                                      %
%                               Section 4                              %
%                                                                      %
%%%%%%%%%%%%%%%%%%%%%%%%%%%%%%%%%%%%%%%%%%%%%%%%%%%%%%%%%%%%%%%%%%%%%%%%


\section{Power Series and Differential Equations}

So far, we have begun with functions we already know, in the sense of
being able to calculate the value of the function and all its
derivatives at at least one point.  This in turn allowed us to write
down the corresponding Taylor series. Often, though, we don't even
have this much information about a function.  In such cases it is
frequently useful to assume that there is some infinite
polynomial---called a {\bf power series}\index{power series}---which
represents the function, and then see if we can determine the
coefficients of the polynomial.

This technique is especially useful in dealing with differential
equations.  To see why this is the case, think of the alternatives.
If we can approximate the solution $y=y(x)$ to a certain differential
equation to an acceptable degree of accuracy by, say, a 20-th degree
polynomial, then the only storage space required is the insignificant
space taken to keep track of the 21 coefficients.  Whenever we want
the value of the solution for a given value of $x$, we can then get a
quick approximation by evaluating the polynomial at $x$.  Other
alternatives are much more costly in time or in space.  We could use
Euler's method to grind out the solution at $x$, but, as you've
already discovered, this can be a slow and tedious process.  Another
option is to calculate lots of values and store them in a table in the
computer's memory.  This not only takes up a lot of memory space, but
it also only gives values for a finite set of values of $x$, and is
not much faster than evaluating a polynomial.  Until 30 years ago, the
table approach was the standard one---all scientists and
mathematicians had a handbook of mathematical functions containing
hundreds of pages of numbers giving the values of every function they
might need.

 

To see how this can happen, let's first look at a familiar
differential equation whose solutions we already know:
$$
y'=y.
$$
Of course, we know by now that the solutions are $y=ae^x$ for an
arbitrary constant $a$ (where $a=y(0)$).  Suppose, though, that we
didn't already know how to solve this differential equation.  We might
see if we can find a power series of the form
$$
y = a_0 + a_1\,x + a_2\, x^2 + a_3\, x^3 + \cdots + a_n\, x^n + \cdots 
$$ 
that solves the differential equation.  Can we, in fact, determine
values for the coefficients $a_0, a_1, a_2, \cdots , a_n , \cdots $
that will make $y'=y$?

Using the rules for differentiation, we have
$$
y'= a_1 + 2a_2\,x + 3 a_3\,x^2 + 4 a_4\,x^3 + \cdots 
    + n a_n\,x^{n-1} + \cdots .
$$
Two polynomials are equal if and only if the coefficients of
corresponding powers of $x$ are equal.  Therefore, if $y'=y\,,$ it
would have to be true that
\begin{align*}
a_1   &= a_0 \\
2 a_2 &= a_1 \\
3 a_3 &= a_2 \\[-5pt]
      &\ \, \vdots  \\[-5pt]
n a_n &= a_{n-1} \\[-5pt]
      &\ \, \vdots 
\end{align*}
Therefore the values of $a_1$, $a_2$, $a_3$, \ldots are not arbitrary;
indeed, each is determined by the preceding one.  Equations like
these---which deal with a sequence of quantities and relate each term
to those earlier in the sequence---are called {\bf recursion
  relations}\index{recursion relation}.  These recursion relations
permit us to express every $a_n$ in terms of $a_0$:
\addtolength{\arraycolsep}{-3pt}
\[
\begin{array}{rclclcl}
a_1 &=&  a_0  &&&& \\[3pt]
a_2 &=&  \dfrac{1}{2} \, a_1 & &
         &=& \dfrac{1}{2} \, a_0 \\[10pt]
a_3 &=&  \dfrac{1}{3} \, a_2 &=& \dfrac{1}{3}\cdot \dfrac{1}{2} \, a_0 
         &=& \dfrac{1}{3!} \, a_0 \\[10pt]
a_4 &=&  \dfrac{1}{4} \, a_3 &=& \dfrac{1}{4}\cdot\dfrac{1}{3!} \, a_0 
         &=& \dfrac{1}{4!} \, a_0 \\
    &\vdots& &\vdots& &\vdots&  \\
a_n &=&  \dfrac{1}{n} \, a_{n-1}\ &=& 
         \dfrac{1}{n}\cdot\dfrac{1}{(n-1)!} \, a_0 
         &=& \dfrac{1}{n!} \, a_0 \\[-5pt]
    &\vdots& &\vdots& &\vdots&
\end{array}
\]
Notice that $a_0$ remains ``free": there is no equation that
determines its value.  Thus, without additional information, $a_0$ is
{\em arbitrary}.  The series for $y$ now becomes
\[
y = a_0 + a_0\, x + \dfrac{1}{2!} \, a_0 \,x^2 + \dfrac{1}{3!} \, a_0 \,x^3 
   + \cdots + \dfrac{1}{n!} \, a_0 \,x^n + \cdots
\]
or
\[
y = a_0\left[1 + x + \dfrac{1}{2!} \, x^2 + \dfrac{1}{3!} \, x^3 + \cdots +
\dfrac{1}{n!} \, x^2 + \cdots \right].
\]
But the series in square brackets is just the Taylor series for $e^x
\,$---we have derived the Taylor series from the differential equation
alone, without using any of the other properties of the exponential
function.  Thus, we again find that the solutions of the differential
equation $y'=y$ are
\[
y=a_0e^x,
\]
where $a_0$ is an arbitrary constant.  Notice that $y(0)=a_0$, so the
value of $a_0$ will be determined if the initial value of $y$ is
specified.

\noindent
{\bf Note} In general, once we have derived a power series expression
for a function, that power series will also be the Taylor series for
that function.  Although the two series are the same, the term Taylor
series is typically reserved for those settings where we were able to
evaluate the derivatives through some other means, as in the preceding
section.

\subsection{Bessel's Equation}\index{Bessel equation}\index{Bessel equation!order $p$}

For a new example, let's look at a differential equation that arises
in an enormous variety of physical problems (wave motion, optics, the
conduction of electricity and of heat and fluids, and the stability of
columns, to name a few):
$$
x^2 \cdot y'' + x \cdot y' + (x^2 - p^2) \cdot y = 0 \, .
$$
This is called the {\bf Bessel equation of order} {\boldmath $p$}.
Here $p$ is a parameter specified in advance, so we will really have a
different set of solutions for each value of $p$.  To determine a
solution completely, we will also need to specify the initial values
of $y(0)$ and $y'(0)$.  The solutions of the Bessel equation of order
$p$ are called {\bf Bessel functions of order} {\boldmath
  $p$}\index{Bessel function!of order $p$}, and the solution for a
given value of $p$ (together with particular initial conditions which
needn't concern us here) is written $J_p(x)$.  In general, there is no
formula for a Bessel function in terms of simpler functions (although
it turns out that a few special cases like $J_{1/2}(x), J_{3/2}(x) ,
\ldots$ can be expressed relatively simply).  To evaluate such a
function we could use Euler's method, or we could try to find a power
series solution.

\aside{Friedrich Wilhelm Bessel (1784--1846)\index{Bessel, Wilhelm
    Friedrich} was a German astronomer who studied the
  functions that now bear his name in his efforts to analyze the
  perturbations of planetary motions, particularly those of Saturn.}

\enlargethispage*{20pt}
Consider the Bessel equation with $p=0$.  We are thus trying to solve
the differential equation
$$
x^2 \cdot y'' + x \cdot y' + x^2  \cdot y = 0.
$$
By dividing by $x$, we can simplify this a bit to
$$
x \cdot y'' + y' + x \cdot y = 0.
$$

Let's look for a power series expansion
$$
y = b_0 + b_1\,x + b_2\,x^2 +b_3\,x^3 + \cdots .
$$
We have
\begin{align*}
y' &= b_1 + 2b_2x + 3b_3x^2 + 4b_4x^3 + \cdots 
   + (n+1)b_{n+1}x^n +\cdots 
\intertext{and}
y'' &= 2b_2 + 6b_3x + 12b_4x^2 + 20b_5x^3 + \cdots 
   + (n+2)(n+1)b_{n+2}x^n + \cdots , \\
\end{align*}
We can now use these expressions to calculate the series for the
combination that occurs in the differential equation:
\[
\begin{array}{rcccccccl}
xy''&=& & &2b_2\,x&+&6b_3\,x^2&+&\cdots \\
y'&=&b_1 & +&2b_2\,x&+&3b_3\,x^2&+&\cdots \\
xy&=& &  &b_0\,x&+&b_1\,x^2&+&\cdots  \\[1pt] \hline
\\ [-10pt]
xy''+y'+xy&=&b_1&+&(4b_2+b_0 )x&+&(9b_3+b_1)x^2&+&\cdots
\end{array}
\] 
\addtolength{\arraycolsep}{3pt}%
In general, 
%
\mar{Finding the\\ coefficient of $x^n$}
%
the coefficient of $x^n$ in the combination will be 
\[
(n+1)n\, b_{n+1} + (n+1)\, b_{n+1} +b_{n-1} = (n+1)^2\, b_{n+1} +b_{n-1}.
\]

%\enlargethispage*{25pt}
If the power series $y$ is to be a solution to the original
differential equation, the infinite series for $xy''+y'+xy$ must
equal~0. This in turn means that every coefficient of that series must
be~0.  We thus get
\begin{align*}
b_1 &= 0, \\ 
4b_2+b_0 &= 0, \\ 
9b_3+b_1 &= 0, \\[-5pt] 
&\ \, \vdots \\[-5pt] 
n^2b_n +b_{n-2} &= 0. \\[-5pt]
 &\ \, \vdots
\end{align*}
If we now solve these recursively as before, we see first off that
since $b_1=0$, it must also be true that
\[
b_k=0 \qquad  \text{for every odd $k$}.
\]

\newpage

\noindent
For the even coefficients we have
\begin{align*}
b_2&=-\frac{1}{2^2}b_0, \\
b_4&=-\frac{1}{4^2}b_2= \frac{1}{2^24^2}b_0, \\
b_6&=-\frac{1}{6^2}b_4=-\frac{1}{2^2 4^2 6^2}b_0,
\end{align*}
and, in general,
\[
b_{2n} = \pm \frac{1}{2^2 4^2 6^2 \cdots (2n)^2}b_0
    = \pm \frac{1}{2^{2n} (n!)^2}b_0.
\]


Thus any function $y$ satisfying the Bessel equation of order 0 must be of the form
\[
y = b_0\left(1-\frac{x^2}{2^2}+\frac{x^4}{2^4(2!)^2} 
    - \frac{x^6}{2^6(3!)^2} +\cdots \right).
\]
In particular, if we impose the initial condition $y(0)=1$ (which
requires that $b_0=1$), we get the 0-th order Bessel function
$J_0(x)$\index{Bessel function!$J_0$}:
\[
J_0(x) = 1 - \frac{x^2}{4} + \frac{x^4}{64} - \frac{x^6}{2304} 
   + \frac{x^8}{147456}+\cdots .
\]
Here is the graph of $J_0(x)$ together with the 
%
\mar{The graph of the \mbox{Bessel function $J_0$}}
%
polynomial approximations of degree 2, 4, 6, \ldots , 30 over the
interval $[0,14]$:

\vspace*{3in}
\special{psfile=../ps/calc2/Bessel.ps}\index{Bessel function!$J_0$, graph of}

The graph of $J_0$ is suggestive: it appears to be oscillatory, with
decreasing amplitude. Both observations are correct: it can in fact be
shown that $J_0$ has infinitely many zeroes, spaced roughly $\pi$
units apart, and that $\lim_{x \to\infty} J_0(x) = 0$.

\subsection{The $S$-$I$-$R$ Model One More Time}

In exactly the same way, we can find power series solutions when there
are several interacting variables involved. Let's look at the example
we've considered at a number of points in this text to see how this
works.  In the 
%
\mar{The $S$-$I$-$R$ model}
%
$S$-$I$-$R$ model we basically wanted to solve the system of equations
\begin{align*}
S'&= -a SI, \\
I'&= aSI - bI, \\
R'&= bI ,
\end{align*}
where $a$ and $b$ were parameters depending on the specific situation.
Let's look for solutions of the form
\begin{align*}
S&= s_0 + s_1\,t + s_2\,t^2 + s_3\,t^3 + \cdots , \\
I&= i_0 + i_1\,t + i_2\,t^2 + i_3\,t^3 + \cdots , \\
R&= r_0 + r_1\,t + r_2\,t^2 + r_3\,t^3 + \cdots .
\end{align*}

If we put these series in the equation $S'= -a\, S\, I$, we get
\begin{align*}
s_1 + 2s_2t + 3s_3t^2 + \cdots 
  &= -a(s_0 + s_1t + s_2t^2 + \cdots)(i_0 + i_1t + i_2t^2 + \cdots) \\
  &= -a(s_0i_0 + (s_0i_1+s_1i_0)t+(s_0i_2+s_1i_1+s_2i_0)t^2 + \cdots ).
\end{align*}

As before, 
%
\mar{Finding the coefficients of the power series\\ for $S(t)$}
%
if the two sides of the differential equation are to be equal, the
coefficients of corresponding powers of $t$ must be equal:
\begin{align*}
s_1  &= -as_0i_0, \\
2s_2 &= -a(s_0i_1+s_1i_0), \\
3s_3 &= -a(s_0i_2+s_1i_1+s_2i_0), \\[-5pt]
&\ \, \vdots \\[-5pt]
ns_n &= -a(s_0i_{n-1}+s_1i_{n-2}+ \ldots +s_{n-2}i_1+s_{n-1}i_0) \\[-5pt]
&\ \, \vdots \\[-5pt]
\end{align*}

\newpage

\noindent
While this looks messy, it has the crucial recursive feature---each
$s_k$ is expressed in terms of previous terms.
%
\mar{Recursion again}
%
That is, if we knew all the $s$ and the $i$ coefficients out through
the coefficients of, say, $t^6$ in the series for $S$ and $I$, we
could immediately calculate $s_7$.  We again have a {\em recursion
  relation}\index{recursion relation}.

We could expand the equation $I'= aSI - bI$ in the same way,
%
\mar{Finding the \\power series \\for $I(t)$}
%
and get recursion relations for the coefficients $i_k$.  In this
model, though, there is a shortcut if we observe that since $S'=
-aSI$, and since $I'= aSI - bI$, we have $I' = -S' -bI$. If we
substitute the power series in this expression and equate
coefficients, we get
\[
ni_n = -ns_n - bi_{n-1},
\]
which leads to
\[
i_n = -s_n-\frac{b}{n}i_{n-1}
\]
---so if we know $s_n$ and $i_{n-1}$, we can calculate $i_n$.

We are now in a position to calculate the coefficients as far out as
we like.  For we will be given values for $a$ and $b$ when we are
given the model.  Moreover, since $s_0=S(0) = $ the initial
$S$-population, and $i_0=I(0)=$ the initial $I$-population, we will
also typically be given these values as well.  But knowing $s_0$ and
$i_0$, we can determine $s_1$ and then $i_1$.  But then, knowing these
values, we can determine $s_2$ and then $i_2$, and so on.  Since the
arithmetic is tedious, this is obviously a place for a computer.  Here
is a program that calculates the first 50 coefficients in the power
series for $S(t)$ and $I(t)$:\index{program!SIRSERIES}

\begin{center}
{\bf Program: SIRSERIES}
\end{center}
\noindent
\vspace*{-30pt}
\begin{verbatim}
DIM S(0 to 50), I(0 to 50)
a = .00001
b = 1/14
S(0) = 45400
I(0) = 2100
FOR k = 1 TO 50
     Sum = 0
     FOR j = 0 TO k - 1
          Sum = Sum + S(j) * I(k - j - 1)
     NEXT j
     S(k) = -a * SUM/k
     I(k) = -S(k) - b * I(k - 1)/k
NEXT k
\end{verbatim}

\noindent
{\bf Comment:} The opening command in this program introduces a new
feature. It notifies the computer that the variables {\tt S} and {\tt
  I} are going to be {\bf arrays}\index{array}---strings of
numbers---and that each array will consist of 51 elements.  The
element {\tt S(k)} corresponds to what we have been calling $s_k$.
The integer {\tt k} is called the {\bf index}\index{index} of the term
in the array.  The indices in this program run from 0 to 50.

The effect of running this program is thus to create two 51-element
arrays, {\tt S} and {\tt I}, containing the coefficients of the power
series for $S$ and $I$ out to degree 50.  If we just wanted to see
these coefficients, we could have the computer list them.  Here are
the first 35 coefficients for $S$ (read across the rows):

\medskip

\begin{small}
\hspace{-25pt}
\begin{tabular}{rrrrr}
45400    &    $-953.4$  &  $-172.3611$   &  $-17.982061$  &  $-.86969127$ \\
5.4479852e-2 & 1.5212707e-2 &1.4463108e-3  &  4.3532884e-5 &$-7.9100481$e-6\\
$-1.4207959$e-6&$-1.0846994$e-7&$-6.512610$e-10 &9.304633e-10 &1.256507e-10 \\
7.443310e-12&$-2.191966$e-13&$-9.787285$e-14&$-1.053428$e-14&$-4.382620$e-16\\
4.230290e-17 & 9.548369e-18 & 8.321674e-19 & 1.760392e-20  & $-5.533369$e-21\\
$-8.770972$e-22&$-6.101928$e-23&3.678170e-25& 6.193375e-25  & 7.627253e-26 \\ 4.011923e-27&$-1.986216$e-28&$-6.318305$e-29&$-6.271724$e-30&$-2.150100$e-31 \\ \end{tabular}
\end{small}

\medskip

\noindent
Thus the power series for $S$ begins
\small
 $$45400 - 953.4t - 172.3611t^2 - 17.982061t^3 - .86969127t^4 + \\ \cdots - 2.15010\times 10^{-31}t^{34} + \cdots$$
\normalsize

In the same fashion, we find that the power series for $I$ begins
\small
$$2100 +803.4t + 143.66824t^2 + 14.561389t^3 + .60966648t^4 + \cdots + 2.021195\times 10^{-31}t^{34} + \cdots $$
\normalsize

If we now wanted to graph these polynomials over, 
%
\mar{Extending SIRSERIES to graph $S$ and $I$}
%
say, $0\le t\le 10$, we can do it by adding the following lines to
SIRSERIES.  We first define a couple of short subroutines {\tt SUS}
and {\tt INF}\index{subroutines {\tt SUS} and {\tt INF}} to calculate
the polynomial approximations for $S(t)$ and $I(t)$ using the
coefficients we've derived in the first part of the program.  (Note
that these subroutines calculate polynomials in the straightforward,
inefficient way.  If you did the exercises in section~3 which developed
techniques for evaluating polynomials rapidly, you might want to
modify these subroutines to take advantage of the increased speed
available.)  Remember, too, that you will need to set up the graphics
at the beginning of the program to be able to plot.

\newpage

\begin{center}
Extension to \textbf{SIRSERIES}
\end{center}

%this needs revision

\vspace{-20pt}
\noindent
\begin{verbatim}
DEF SUS(x)
     Sum = S(0)
     FOR j = 1 TO 50
         Sum = Sum + S(j) * x^j
     NEXT j
     SUS = Sum
END DEF
DEF INF(x)
     Sum = I(0)
     FOR j = 1 TO 50
         Sum = Sum + I(j) * x^j
     NEXT j
     INF = Sum
END DEF
FOR x = 0 TO 10 STEP .01
\end{verbatim}
\vspace{-9pt}
\rule{24pt}{0pt} {\sl Plot the line from 
                {\tt (x, SUS(x))}} {\sl to 
                {\tt (x + .01, SUS(x + .01))}}\\
\rule{24pt}{0pt} {\sl Plot the line from 
                {\tt (x, INF(x))}} {\sl to 
                {\tt (x + .01, INF(x + .01))}}\\
\vspace{-24pt}
\begin{verbatim}
NEXT x
\end{verbatim}

Here is the graph of $I(t)$ over a 25-day period, together with the polynomial approximations of degree 5, 20, 30, and 70.  

\vspace{180pt}
\special{psfile=../ps/calc2/sircalc.ps}


Note that these polynomials appear to converge to $I(t)$ only out to
values of $t$ around 10.  If we needed polynomial approximations
beyond that point, we could shift to a different point on the curve,
find the values of $I$ and $S$ there by Euler's method, then repeat
the above process. For instance, when $t=12$, we get by Euler's method
that $S(12)=7670$ and $I(12)=27,136$\,. If we now shift our clock to
measure time in terms of $\tau=t-12$, we get the following polynomial
of degree 30:
$$27136 +143.0455\tau-282.0180\tau^2+23.5594\tau^3+.4548\tau^4+\cdots+1.2795\times10^{-25}\tau^{30} \,$$

Here is what the graph of this polynomial looks like when plotted with
the graph of $I$.  On the horizontal axis we list the $t$-coordinates
with the corresponding $\tau$-coordinates underneath.

\vspace*{195pt}
\special{psfile=../ps/calc2/sircalc2.ps}

The interval of convergence seems to be approximately $4<t<20$.  Thus
if we combine this polynomial with the 30-th degree polynomial from
the previous graph, we would have very accurate approximations for $I$
over the entire interval $[0, 20]$.


\subsection{Exercises}


\num  Find power series solutions of the form
$$
y = a_0 + a_1x + a_2 x^2 + a_3 x^3 + \cdots + a_n x^n + \cdots 
$$
for each of the following differential equations.

\sub $y' = 2xy$.

\sub $y'=3x^2y$.

\sub $y'' + xy = 0$.

\sub $y'' + xy' + y = 0$.

\newpage

\numa Find power series solutions to the differential equation
$y''=-y$.  Start with
$$
y = a_0 + a_1x + a_2 x^2 + a_3 x^3 + \cdots + a_n x^n + \cdots .
$$
Notice that, in the recursion relations you obtain, the coefficients
of the even terms are completely independent of the coefficients of
the odd terms.  This means you can get two separate power series, one
with only even powers, with $a_0$ as arbitrary constant, and one with
only odd powers, with $a_1$ as arbitrary constant.  \sub The two power
series you obtained in part a) are the Taylor series centered at $x=0$
of two familiar functions; which ones?  Verify that these functions do
indeed satisfy the differential equation $y''=-y\/$.

\numa Find power series solutions to the differential equation
$y''=y$.  As in the previous problem, the coefficients of the even
terms depend only on $a_0$, and the coefficients of the odd terms
depend only on $a_1$.  Write down the two series, one with only even
powers and $a_0$ as an arbitrary constant, and one with only odd
powers, with $a_1$ as an arbitrary constant.

\sub The two power series you obtained in part a) are the Taylor
series centered at $x=0$ of two {\em hyperbolic trigonometric
  functions\/} (see the exercises in section~3).  Verify that these
functions do indeed satisfy the differential equation
$y''=y$.\index{function!hyperbolic}

\numa Find power series solutions to the differential equation
$y'=xy$, starting with
$$
y = a_0 + a_1x + a_2 x^2 + a_3 x^3 + \cdots + a_n x^n + \cdots .
$$
What recursion relations do you get?  Is $a_1 = a_3 = a_5 = \cdots =
0$?  

\sub Verify that
$$
y = e^{x^2/2}
$$
satisfies the differential equation $y'=xy$.  Find the Taylor series
for this function and compare it with the series you obtained in a)
using the recursion relations.

\num  {\bf The Bessel Equation.}\index{Bessel equation}\index{Bessel function} 

\sub Take $p=1\,$.  The solution satisfying the initial condition $y'
=1/2$ when $x=0$ is defined to be the first order Bessel function
$J_1(x)$\index{Bessel function!$J_1$}. (It will turn out that $y$ has
to be 0 when $x=0$, so we don't have to specify the initial value
of~$y$; we have no choice in the matter.) Find the first five terms of
the power series expansion for $J_1(x)$.  What is the coefficient of
$x^{2n+1}$?

\sub Show by direct calculation from the series for $J_0$ and $J_1$
that
$$
J_0' = - J_1.
$$

\sub To see, from another point of view, that $J_0' = -
J_1$,\index{Bessel function!$J_0$}\index{Bessel function!$J_1$} take
the equation
$$
x \cdot J_0'' + J_0' + x \cdot J_0 = 0
$$
and differentiate it.  By doing some judicious cancelling and
rearranging of terms, show that
$$
x^2 \cdot (J_0')'' + x \cdot (J_0')' + (x^2 -1)(J_0') = 0.
$$
This demonstrates that $J_0'$ is a solution of the Bessel equation
with $p=1$.

\numa When we found the power series expansion for solutions to the
0-th order Bessel equation, we found that all the odd coefficients had
to be 0.  In particular, since $b_1$ is the value of $y'$ when $x=0$,
we are saying that all solutions have to be flat at $x=0$.  This
should bother you a bit.  Why can't you have a solution, say, that
satisfies $y=1$ and $y'=1$ when $x=0$?  \sub You might get more
insight on what's happening by using Euler's method, starting just a
little to the right of the origin and moving left.  Use Euler's method
to sketch solutions with the initial values

\smallskip

\noindent
\begin{tabular}{rllcl}
 i.& $y=2$ & $y'=1$& when &$x=1$, \\
  ii.& $y=1.1$ & $y'=1$ & when & $x=.1$, \\
 iii.& $y=1.01$ &  $y'=1$ & when & $x=.01.$ \\
 \end{tabular} 
 
 \smallskip
\noindent
What seems to happen as you approach the $y$-axis?

\num  {\bf Legendre's differential equation}\index{Legendre's equation}
$$
(1 - x^2)y'' - 2xy' + \ell(\ell + 1)y = 0
$$
arises in many physical problems---for example, in quantum mechanics,
where its solutions are used to describe certain orbits of the
electron in a hydrogen atom.  In that context, the parameter $\ell$ is
called the {\em angular momentum}\index{angular momentum} of the
electron; it must be either an integer or a ``half-integer'' (i.e., a
number like 3/2).  Quantum theory\index{quantum theory} gets its name
from the fact that numbers like the angular momentum of the electron
in the hydrogen atom are ``quantized", that is, they cannot have just
any value, but must be a multiple of some
``quantum"\index{quantum}---in this case, the number 1/2.  \sub Find
power series solutions of Legendre's equation.  \sub Quantization of
angular momentum has an important consequence.  Specifically, when
$\ell$ is an integer it is possible for a series solution to
stop---that is, to be a polynomial.  For example, when $\ell = 1$ and
$a_0 = 0$ the series solution is just $y = a_1x$---all higher order
coefficients turn out to be zero.  Find polynomial solutions to
Legendre's equation for $\ell = 0, 2, 3, 4,$ and $5$ (consider $a_0=0$
or $a_1=0$).  These solutions are called, naturally enough, {\bf
  Legendre polynomials}\index{Legendre polynomials}.

\num It turns out that the power series solutions to the $S$-$I$-$R$
model have a finite interval of convergence.  By plotting the power
series solutions of different degrees against the solutions obtained
by Euler's method, estimate the interval of convergence.

\numa {\bf Logistic Growth}\index{logistic growth} Find the first five
terms of the power series solution to the differential equation
\[
y'=y(1-y).
\]
Note that this is just the logistic equation, where we have chosen our
units of time and of quantities of the species being studied so that
the carrying capacity is 1 and the intrinsic growth rate is 1.

\sub Using the initial condition $y=.1$ when $x=0$, plot this power
series solution on the same graph as the solution obtained by Euler's
method.  How do they compare?

\sub Do the same thing with initial conditions $y=2$ when $x=0$.

\newpage
%%%%%%%%%%%%%%%%%%%%%%%%%%%%%%%%%%%%%%%%%%%%%%%%%%%%%%%%%%%%%%%%%%%%%%%%
%                                                                      %
%                               Section 5                              %
%                                                                      %
%%%%%%%%%%%%%%%%%%%%%%%%%%%%%%%%%%%%%%%%%%%%%%%%%%%%%%%%%%%%%%%%%%%%%%%%


\section{Convergence}

We have written expressions such as
\[
\sin(x) = x - \dfrac{x^3}{3!} + \dfrac{x^5}{5!} 
                - \dfrac{x^7}{7!} + \dfrac{x^9}{9!} - \cdots ,
\] 
meaning that for any value of $x$ the series on the right will
converge to $\sin(x)$.  There are a couple of issues here.  The first
is, what do we even mean when we say the series ``converges'', and how
do we prove it converges to $\sin(x)$?  If $x$ is small, we can
convince ourselves that the statement is true just by trying it.  If
$x$ is large, though, say $x=100^{100}$, it would be convenient to
have a more general method for proving the stated convergence.
Further, we have the example of the function $1/(1+x^2)$ as a
caution---it seemed to converge for small values of $x$ ($|x|<1$), but
not for large values.

Let's first clarify what we mean by convergence\index{convergence}.
It is,
%
\mar{Convergence means\\ essentially that\\ ``decimals stabilize"}
%
essentially, the intuitive notion of ``decimals stabilizing'' that we
have been using all along.  To make explicit what we've been doing,
let's write a ``generic" series
\[
b_0 + b_1 + b_2 + \cdots = \sum_{m=0}^\infty b_m .
\]
When we evaluated such a series, we looked at the {\bf partial
  sums}\index{sum!partial}
\addtolength{\arraycolsep}{-3pt}
\[
\begin{array}{lclcl}
S_1&=&b_0 + b_1 &=&{\displaystyle \sum_{m=0}^1 b_m}, \\[.18in]
S_2&=&b_0 + b_1+b_2 &=&{\displaystyle\sum_{m=0}^2 b_m}, \\[.18in]
S_3&=&b_0 + b_1 + b_2 +b_3&=&{\displaystyle\sum_{m=0}^3 b_m}, \\[-5pt]
&\vdots & &\vdots &  \\[-5pt]
S_n&=&b_0 + b_1 + b_2 + \ldots+b_n&=&{\displaystyle\sum_{m=0}^n b_m}, \\[-5pt]
&\vdots&&\vdots & 
\end{array}
\]
\addtolength{\arraycolsep}{3pt}
\medskip

\enlargethispage*{20pt}
Typically, when we calculated a number of these partial sums, we
noticed that beyond a certain point they all seemed to agree on, say,
the first 8 decimal places.  If we kept on going, the partial sums
would agree on the first 9 decimals, and, further on, on the first 10
decimals, etc. This is precisely what we mean by
convergence:\index{convergence!definition of}\index{sum of a series!definition of}\index{divergence!definition of}

\medskip
\noindent
\begin{center}
\framebox[5.4in]
{
\parbox{4.9in}
{
\vspace{10pt}
The infinite series 
        $$
        b_0 + b_1 + b_2 + \cdots= \sum_{m=0}^\infty b_m
        $$
{\bf converges} if, no matter how many decimal places are specified,
it is always the case that the partial sums eventually agree to at
least this many decimal places.
    
    \smallskip
    
Put more formally, we say the series converges if, given any number
$D$ of decimal places, it is always possible to find an integer $N_D$
such that if $k$ and $n$ are both greater than $N_D$, then $S_k$ and
$S_n$ agree to at least $D$ decimal places.
    
    \medskip

The number defined by these stabilizing decimals is called the {\bf
  sum} of the series.
    
    \medskip

If a series does not  converge, we say it {\bf diverges}.
    \vspace{6pt} 
    }
   }
\end{center}
\medskip

In other words, 
%
\mar{What it means for\\ an infinite sum\\ to converge}
% 
for me to prove to you that the Taylor series for $\sin(x)$ converges
at $x=100^{100}$, you would specify a certain number of decimal
places, say 5000, and I would have to be able to prove to you that if
you took partial sums with enough terms, they would all agree to at
least 5000 decimals.  Moreover, I would have to be able to show the
same thing happens if you specify {\em any} number of decimal places
you want agreement on.

How can this be done?  It seems like an enormously daunting task to be
able to do for any series.  We'll tackle this challenge in stages.
First we'll see what goes wrong with some series that don't
converge---divergent series.  Then we'll look at a particular
convergent series---the {\bf geometric series}---that's relatively
easy to analyze.  Finally, we will look at some more general rules
that will guarantee convergence of series like those for the sine,
cosine, and exponential functions.


\subsection{Divergent Series}

Suppose we have an infinite series
\[
b_0 + b_1 + b_2 + \cdots = \sum_{m=0}^\infty b_m,
\]
and consider two successive partial sums, say
\[
S_{122}=b_0 + b_1 + b_2 + \ldots+b_{122}=\sum_{m=0}^{122} b_m
\]
and
\[
S_{123}=b_0 + b_1 + b_2 + \ldots+b_{122}+b_{123}=\sum_{m=0}^{123} b_m.
\]
Note that these two sums are the same, except that the sum for
$S_{123}$ has one more term, $b_{123}$, added on.  Now suppose that
$S_{122}$ and $S_{123}$ agree to 19 decimal places.  In section~2 we
defined this to mean $|S_{123}-S_{122}|<.5 \times 10^{-19}$.  But
since $S_{123}-S_{122}=b_{123}$, this means that $|b_{123}|<.5 \times
10^{-19}$.  To phrase this more generally,

\noindent
\begin{center}
\framebox[4.3in]
{\parbox{3.85in}{\vspace{4pt}
 Two successive partial sums, $S_n$ and $S_{n+1}$, agree out to $k$
 decimal places if and only if $|b_{n+1}|<.5 \times 10^{-k}$.
 \vspace{4pt} }
 }
\end{center}


But since our definition of convergence required that we be able to
fix any specified number of decimals provided we took partial sums
lengthy enough, it must be true that {\em if the series converges},
%
\mar{A necessary condition for convergence}
%
the individual terms $b_k$ must become arbitrarily small if we go out
far enough. Intuitively, you can think of the partial sums $S_k$ as
being a series of approximations to some quantity.  The term $b_{k+1}$
can be thought of as the ``correction'' which is added to $S_k$ to
produce the next approximation $S_{k+1}$.  Clearly, if the
approximations are eventually becoming good ones, the corrections made
should become smaller and smaller.  We thus have the following
necessary condition for convergence:

\smallskip

\begin{center}
\framebox[4.5in]
{ \parbox{4.5in}{\begin{center}\vspace{-6pt}
If the infinite series $b_0 + b_1 + b_2 + \cdots 
   = \displaystyle \sum_{m=0}^\infty b_m$ converges, \\
then $\displaystyle \lim_{k\rightarrow \infty} b_k = 0.$
\end{center}\vspace{-6pt}}
}
\end{center}

\smallskip

{\bf Remark:} It is important to recognize what this criterion does
and does \mar{{\em Necessary} and {\em sufficient} mean different
  things} not say---it is a {\bf necessary} condition for convergence
(i.e., every convergent sequence has to satisfy the condition
$\lim_{k\rightarrow \infty} b_k = 0$)---but it is not a {\bf
  sufficient} condition for convergence (i.e., there are some
divergent sequences that also have the property that
$\lim_{k\rightarrow \infty} b_k = 0$).  The criterion is usually used
to detect some divergent series, and is more useful in the following
form (which you should convince yourself is equivalent to the
preceding):\index{divergence criterion}

\noindent
\begin{center}
\framebox[5.3in]
{  \parbox{4.8in}{\vspace{-2pt}
\begin{center}
If \quad ${\displaystyle \lim_{k\rightarrow \infty} b_k \ne 0 }$, \qquad 
(either because the limit doesn't exist at all,\\ or it equals something 
besides 0), then the infinite series 
\[
b_0 + b_1 + b_2 + \cdots= \sum_{m=0}^\infty b_m \qquad \text{diverges}.
\]
\end{center}
\vspace{-6pt}
}
}
\end{center}

This criterion allows us to detect a number of divergent series 
%
\mar{Detecting\\ divergent series}
%
right away.  For instance, we saw earlier that the statement
\[
\dfrac{1}{1+x^2} = 1 - x^2 + x^4 - x^6 + \cdots 
\]
appeared to be true only for $|x|<1$. Using the remarks above, we can
see why this series has to diverge for $|x|\ge 1$.  If we write
$1-x^2+x^4-x^6+\cdots$ as $b_0 + b_1 + b_2 + \cdots$, we see that
$b_k=(-1)^kx^{2k}$.  Clearly $b_k$ does not go to 0 for $|x|\ge
1$---the successive ``corrections'' we make to each partial sum just
become larger and larger, and the partial sums will alternate more and
more wildly from a huge positive number to a huge negative number.
Hence the series converges {\em at most} for $-1<x<1$.  We will see in
the next subsection how to prove that it really does converge for all
$x$ in this interval.

Using exactly the same kind of argument, we can show that the
following series also diverge for $|x|>1$:
\[
\begin{array}{ccl}
f(x) &\hspace{.5in} & \hspace{.5in}\mbox{Taylor series for $f(x)$} \\ [2pt] 
\hline & & \\ [-6pt] 
\ln{(1-x)}& & - \left(x + \dfrac{x^2}{2} + \dfrac{x^3}{3} 
                + \dfrac{x^4}{4} + \cdots \right)\\ [.15in]
\dfrac{1}{1-x}& & 1 + x + x^2 + x^3 + \cdots  \\ [.15in]
(1+x)^c & & 1 + cx + \dfrac{c(c-1)}{2!}x^2 
              + \dfrac{c(c-1)(c-2)}{3!}x^3 + \cdots  \\ [.15in]
\arctan{x} & & x - \dfrac{x^3}{3} + \dfrac{x^5}{5} - \cdots
                  \pm  \dfrac{x^{2n+1}}{2n+1}
\end{array}
\]
The details are left to the exercises.  While these common series all
happen to diverge for $|x|>1$, it is easy to find other series that
diverge for $|x|>2$ or $|x|>17$ or whatever---see the exercises for
some examples.

\subsubsection{The Harmonic Series}

We stated earlier in this section that simply knowing that the
individual terms $b_k$ go to 0 for large values of $k$ does not
guarantee that the series
\[
b_0 + b_1 + b_2 + \cdots
\]
will converge.  
%
\mar{An important counterexample}
%
Essentially what can happen is that the $b_k$ go to 0 slowly enough
that they can still accumulate large values.  The classic example of
such a series is the {\bf harmonic series}\index{harmonic series}:
$$
1 + \dfrac{1}{2} + \dfrac{1}{3} + \dfrac{1}{4} + \cdots 
    = \sum_{i=1}^\infty \dfrac{1}{i}.
$$
It turns out that this series just keeps getting larger as you add
more terms. It is eventually larger than 1000, or 1 million, or
$100^{100}$ or \ldots.  This fact is established in the exercises.  A
suggestive argument, though, can be quickly given by observing that
the harmonic series is just what you would get if you substituted
$x=1$ into the power series
\[
x + \dfrac{x^2}{2} + \dfrac{x^3}{3} + \dfrac{x^4}{4} + \cdots .
\]
But this is just the Taylor series for $-\ln(1-x)$, and if we
substitute $x=1$ into this we get $-\ln{0}$, which isn't defined.
Also, $\lim_{x \to 0} -\ln{x}=+\infty$.

\subsection{The Geometric Series}

A series occurring frequently in a wide range of contexts is the {\bf
  geometric series}\index{geometric series}
\[
G(x) = 1 + x + x^2 +x^3 + x^4 +\cdots ,
\]
This is also a sequence we can analyze completely and rigorously in
terms of its convergence.  It will turn out that we can then reduce
the analysis of the convergence of several other sequences to the
behavior of this one.

By the analysis we performed above, 
%
\mar{To avoid divergence, $|x|$ must be less than 1}
%
if $|x|\ge 1$ the individual terms of the series clearly don't go to
0, and the series therefore diverges.  What about the case where
$|x|<1$?

\newpage

The starting point is the partial sums.  A typical partial sum looks like:
        $$
        S_n =  1 +  x +  x^2 +  x^3 + \cdots +  x^n\, .
        $$
This is a finite number; 
%
\mar{A simple expression for the partial sum $S_n$}
%
we must find out what happens to it as $n$ grows without bound.  Since
$S_n$ is finite, we can calculate with it.  In particular,
\[
xS_n =   x +  x^2 +  x^3 + \cdots +  x^n +  x^{n+1}.
\]
Subtracting the second expression from the first, we get
        $$
        S_n - xS_n = 1  -  x^{n+1},
        $$
and thus (if $x \neq 1$)
        $$
        S_n =   \dfrac{1 - x^{n+1}}{1-x}.
        $$
(What is the value of $S_n$ if $x=1$?)

This is a handy, compact form for the partial sum.  Let us see what value it has  for various values of $x$.  For example, if $x = 1/2$, then
\[
\begin{array}{rccccccccc}
n: & 1 & 2 & 3 & 4 & 5 & 6 & \cdots & \to & \infty
                \\[.07in]
S_n: & 1 & \dfrac{3}{2}  & \dfrac{7}{4}  & \dfrac{15}{8}  
                & \dfrac{31}{16}  & \dfrac{63}{32}  & \cdots & \to 
                & 2  
\end{array}
\]


It appears that as $n \to \infty$, $S_n \to 2 \/$.  
%
\mar{Finding the limit of $S_n$ as $n \to \infty$}
%
Can we see this algebraically?
\begin{align*}
S_n &= \dfrac{1 - (1/2)^{n+1}}{1 - \tfrac{1}{2}} \\
    &= \dfrac{1 - (1/2)^{n+1}}{1/2} \\
    &= 2 \cdot (1 - (1/2)^{n+1})
     = 2-(1/2)^n. 
\end{align*}
As $n \to \infty$, $(1/2)^n \to 0$, so the values of $S_n$ become
closer and closer to $2$.  The series converges, and its sum
is 2.

Similarly, 
%
\mar{Summing another geometric series}
%
when $x = -1/2$, the partial sums are
        $$
        S_n = \dfrac{1 - (-1/2)^{n+1}}{3/2} =
                \dfrac{2}{3} \left(
                1 \pm \dfrac{1}{2^{n+1}} \right).
        $$
The presence of the $\pm$ sign does not alter the outcome: since
\mbox{$(1/2^{n+1}) \to 0$}, the partial sums converge to $2/3$.
Therefore, we can say the series converges and its sum is $2/3$.

\newpage

In exactly the same way, though, for any $x$ satisfying $|x|<1$ we have
\[
S_n = \dfrac{1 - x^{n+1}}{1-x} =  \dfrac{1}{1-x}(1-x^{n+1}),
\]
and as $n \to \infty$, $x^{n+1} \to 0$.  Therefore, $S_n \to 1/(1-x)$.
Thus the series converges, and its sum is $1/(1-x)$.

To summarize, we have thus proved that
        \mar{Convergence and divergence of the geometric series}

\noindent
\begin{center}
\framebox[5.0in] {\parbox{4.45in}{\vspace{8pt} 
The geometric
    series\index{geometric series!convergence
      criterion}\index{convergence criterion!geometric series}
\[
G(x) = 1 + x + x^2 +x^3 + x^4 +\cdots
\]
converges for all $x$ such that $|x|<1$. In such cases the sum is
\[
\dfrac{1}{1-x }.
\]
The series diverges for all other values of $x$.
\vspace{6pt}}
}
\end{center}

As a final comment, note that the formula
\[
S_n =  \dfrac{1 - x^{n+1}}{1-x}
\]
is valid for all $x$ except $x=1$.  Even though the partial sums
aren't converging to any limit if $x>1$, the formula can still be
useful as a quick way for summing powers.  Thus, for instance
\[
1 + 3 + 9 + 27 + 81 + 243 = \dfrac{1-3^6}{1-3}
   = \dfrac{1-729}{-2} = \dfrac{-728}{-2}=364,
\]
and
\[
1 - 5 + 25 - 125 + 625 - 3125 = \dfrac{1-(-5)^6}{1-(-5)}
   = \dfrac{1-15625}{6} = -264.
\]


\subsection{Alternating Series}

A large class of common 
%
\mar{Many common series are alternating, at least for some values of $x$}
%
power series consists of the {\bf alternating
  series}\index{alternating series!definition of}---series in which
the terms are alternately positive and negative.  The behavior of such
series is particularly easy to analyze, as we shall see in this
section. Here are some examples of alternating series we've already
encountered :

\newpage

\mbox{}\vspace{-20pt}
\begin{align*}
\sin{x} &= x - \dfrac{x^3}{3!} + \dfrac{x^5}{5!} 
                - \dfrac{x^7}{7!} + \cdots , \\ [.15in]
\cos{x} &= 1 - \dfrac{x^2}{2!} + \dfrac{x^4}{4!} 
                - \dfrac{x^6}{6!} + \cdots , \\ [.15in]
\dfrac{1}{1+x^2}&= 1 - x^2 + x^4 - x^6 + \cdots ,
                 \mbox{\hspace{.5in} (for $|x|<1$)}.
\end{align*}
Other series may be alternating for some, but not all, values of
$x$. For instance, here are two series that are alternating for
negative values of $x$, but not for positive values:
\begin{align*}
e^x &= 1 + x + \dfrac{x^2}{2!} + \dfrac{x^3}{3!} + 
          \dfrac{x^4}{4!} + \cdots ,  \\[.15in]
\ln(1-x) &= -(x + x^2 + x^3 + x^4 + \cdots )
         \mbox{\hspace{.5in} (for $|x|<1$)}.
\end{align*} 

\medskip
\noindent 
{\bf Convergence criterion for alternating series.}  Let us write a generic alternating series as
$$
b_0-b_1 + b_2 - b_3 + \cdots + (-1)^mb_m + \cdots \, ,
$$
where the $b_m$ are positive. It turns out that an alternating series converges if the terms $b_m$ both consistently shrink in size and approach zero:\index{alternating series!convergence criterion}\index{convergence criterion!alternating series}

\medskip
\noindent
\begin{center}
\framebox[4.6in]
{\parbox{4.45in}{\vspace{-4pt}
\begin{center}
$b_0-b_1 + b_2 - b_3 + \cdots + (-1)^mb_m + \cdots$ \hspace{.2in} 
  converges if \\[6pt]
$0<b_{m+1} \leq b_m$ \quad for all $\; m$ \quad and \quad ${\displaystyle \lim_{m \to \infty} b_m = 0}$.\\
\end{center}\vspace{-8pt} 
}
}\index{alternating series!convergence of}
\end{center}

\medskip

It is this property 
%
\mar{Alternating series\\ are easy to test\\ for convergence}
%
that makes alternating series particularly easy to deal with. Recall
that this is {\em not} a property of series in general, as we saw by
the example of the harmonic series.  The reason it is true for
alternating series becomes clear if we view the behavior of the
partial sums geometrically:

\vspace{120pt}

\special{psfile=../ps/calc2/altseries.ps}

\newpage

We mark the partial sums $S_n$ on a number line.  The first sum
$S_0=b_0$ lies to the right of the origin.  To find $S_1$ we go to the
left a distance $b_1$.  Because $b_1 \leq b_0$, $S_1$ will lie between
the origin and $S_0$.  Next we go to the right a distance $b_2$, which
brings us to $S_2$.  Since $b_2 \leq b_1$, we will have $S_2 \leq
S_0$.  The next move is to the left a distance $b_3$, and we find $S_3
\geq S_1$.  We continue going back and forth in this fashion, each
step being less than or equal to the preceding one, since $b_{m+1}
\leq b_m$.  We thus get
\[
0\leq S_1 \leq S_3 \leq S_5 \leq \ldots \leq S_{2m-1} \leq 
   \cdots \leq S_{2m} \leq \ldots \leq S_4 \leq S_2 \leq S_0.
\]

The partial sums oscillate back and forth, 
%
\mar{The partial sums oscillate, with the exact sum trapped between
  consecutive partial sums}
%
with all the odd sums on the left increasing and all the even sums on
the right decreasing. Moreover, since $|S_{n}-S_{n-1}| = b_n$, and
since $\lim_{n \to \infty} b_n = 0$, the difference between
consecutive partial sums eventually becomes arbitrarily small---the
oscillations take place within a smaller and smaller interval.  Thus
given any number of decimal places, we can always go far enough out in
the series so that $S_k$ and $S_{k+1}$ agree to that many decimal
places. But if $n$ is any integer greater than $k$, then, since $S_n$
lies between $S_k$ and $S_{k+1}$, $S_n$ will also agree to that many
decimal places---those decimals will be fixed from $k$ on out.  The
series therefore converges, as claimed---the sum is the unique number
$S$ that is greater than all the odd partial sums and less than all
the even partial sums.

For a convergent alternating
%
\mar{A simple estimate\\ for the accuracy\\ of the partial sums}
%
series, we also have a particularly simple bound for the
error\index{error bound!alternating series}\index{alternating
  series!error bound} when we approximate the sum $S$ of the series by
partial sums.

\noindent
\begin{center}
\framebox[5.3in]
{\parbox{4.8in}{\vspace{8pt}
If \hfill $S_n = b_0-b_1 + b_2 - \cdots \pm b_n$, \hfill\mbox{}\\[5pt]
and if \hfill 
  $0<b_{m+1} \leq b_m \quad \mbox{for all $m$ \quad and} \quad 
     {\displaystyle \lim_{m \to \infty} b_m = 0}$, \hfill\mbox{}\\
(so the series converges), then 
\[
|S-S_n| < b_{n+1}.
\]
In words, the error in approximating $S$ by $S_n$ is less than the
next term in the series.

\vspace{4pt}}
}
\end{center}

\noindent
{\bf Proof}: Suppose $n$ is odd.  Then we have, as above, that
$S_n<S<S_{n+1}$.  Therefore $0<S-S_n<S_{n+1}-S_n=b_{n+1}$.  If $n$ is
even, a similar argument shows $0<S_n-S<S_n-S_{n+1}=b_{n+1}$.  In
either case, we have $|S-S_n| < b_{n+1}$, as claimed.

Note further that we also know whether $S_n$ is too large or too
small, depending on whether $n$ is even or odd.

\medskip
\noindent 
{\bf Example}.  Let's apply the error estimate 
%
\mar{Estimating the error in approximating $\cos(.7)$}
%
for an alternating series to analyze the error if we approximate
$\cos(.7)$ with a Taylor series with three terms:
$$
\cos(.7) \approx        1 - \dfrac{1}{2!} (.7)^2 
                + \dfrac{1}{4!} (.7)^4
                = 0.765004166\ldots .
$$
Since the last term in this partial sum was an addition, this
approximation is too big.  To get an estimate of how far off it might
be, we look at the next term in the series:
$$
\dfrac{1}{6!} (.7)^6 = .0001634\ldots .
$$
We thus know that the correct value for $\cos(.7)$ is somewhere in the
interval
\[
.76484 = .76500 - .00016 \leq \cos(.7) \leq .76501,
\]
so we know that $\cos(.7)$ begins $.76\ldots$ and the third decimal is
either a~4 or a~5.  Moreover, we know that $\cos(.7) = .765$ rounded
to 3~decimal places.

If we use the partial sum with four terms, we get
$$
\cos(.7) \approx        1 - \dfrac{1}{2!} (.7)^2 
                + \dfrac{1}{4!} (.7)^4
                - \dfrac{1}{6!} (.7)^6 
                = .764840765\ldots ,
$$
and the error would be less than
$$
\dfrac{1}{8!}(.7)^8 = .0000014\ldots < .5 \times 10^{-5},
$$
so we could now say that $\cos(.7) =.76484\ldots .$

If we wanted to know in advance 
%
\mar{How many terms\\ are needed\\ to obtain accuracy\\ to 12 decimal places?}
%
how far out in the series we would have to go to determine $\cos(.7)$
to, say, 12 decimals, we could do it by finding a value for $n$ such
that
$$
b_n = \dfrac{1}{n!}(.7)^n \leq .5 \times 10^{-12}.
$$
With a little trial and error, we see that $b_{12} \approx .3 \times
10^{-10}$, while \mbox{$b_{14} < 10^{-13}$}.  Thus if we take the
value of the 12th degree approximation for $\cos(.7)$, we can be
assured that our value will be accurate to 12 places.

We have met this capability of getting an error estimate in a single
step before, in version 3 of Taylor's theorem.  It is in contrast to the
approximations made in dealing with general series, where we typically
had to look at the pattern of stabilizing digits in the succession of
improving estimates to get a sense of how good our approximation was,
and even then we had no guarantee.

\medskip
\noindent
{\bf Computing $e$.}  Because of the fact that we
%
\mar{It may be possible\\ to convert\\ a given problem\\ to one involving\\
  alternating series}
%
can find sharp bounds for the accuracy of an approximation with
alternating series, it is often desirable to convert a given problem
to this form where we can.  For instance, suppose we wanted a good
value for $e$.  The obvious thing to do would be to take the Taylor
series for $e^x$ and substitute $x=1$.  If we take the first 11 terms
of this series we get the approximation
\[
e = e^1 \approx 1 + 1 + \dfrac{1}{2!} + \dfrac{1}{3!} + \cdots 
     + \dfrac{1}{10!} = 2.718281801146\ldots ,
\]
but we have no way of knowing how many of these digits are correct.

Suppose instead, that we evaluate $e^{-1}$:
\[
e^{-1} \approx 1 - 1 + \dfrac{1}{2!} - \dfrac{1}{3!} + \cdots 
     + \dfrac{1}{10!} = .367879464286\ldots .
\]
Since $1/(11!) = .000000025\ldots$, we know this approximation is
accurate to at least 7 decimals.  If we take its reciprocal we get
\[
1/.3678794624286\ldots = 2.718281657666\ldots ,
\]
which will then be accurate to 6 decimals (in the exercises you will
show why the accuracy drops by 1 decimal place), so we can say
$e=2.718281\ldots .$




\subsection{The Radius of Convergence}

We have seen examples of power series that converge for all $x$ (like
the Taylor series for $\sin{x}$) and others that converge only for
certain $x$ (like the series for $\arctan{x}$).  How can we determine
the convergence of an arbitrary power series of the form
        $$
         a_0 + a_1 x + a_2 x^2 + a_3 x^3  + \cdots 
                + a_n x^n + \cdots \ ?
        $$
We must suspect that this series {\em may not\/} converge for all
values of $x$.  For example, does the Taylor series
        $$
        1 + x + \dfrac{1}{2!} x^2 + \dfrac{1}{3!} x^3 + \cdots
        $$
converge for all values of $x$, or only some?  When it converges, does
it converge to $e^x$ or to something else?  After all, this Taylor
series is designed to look like $e^x$ only near $x = 0$; it remains to
be seen how well the function and its series match up far from $x =
0$.

The question of convergence has a definitive answer.  It goes like
this: if the power series
        $$
         a_0 + a_1 x + a_2 x^2 + a_3 x^3  + \cdots 
                + a_n x^n + \cdots \ .
        $$
converges for a particular value of $x$, say $x = s$, then it
automatically
%
\mar{The answer to the convergence question}
%
converges for any {\em smaller\/} value of $x$ (meaning any $x$ that
is closer to the origin than $s$ is; i.e, any $x$ for which $|x| <
|s|$).  Likewise, if the series {\em diverges\/} for a particular
value of $x$, then it also diverges for any value farther from the
origin.  In other words, the values of $x$ where the series converges
are not interspersed with the values where it diverges.  On the
contrary, within a certain distance $R$ from the origin there is only
convergence, while beyond that distance there is only divergence.  The
number $R$ is called the {\bf radius of convergence}\index{radius of
  convergence} of the series, and the range where it converges is
called its {\bf interval of convergence}\index{interval of
  convergence}.

\begin{picture}(360,100)(0,-32)
\put(10,30){\vector(1,0){320}}
\put(100,30){\line(0,1){3}}
\put(170,30){\line(0,1){3}}
\put(240,30){\line(0,1){3}}
\put(334,25){$x$}
\put(90,19){$-R$}
\put(167.5,19){$0$}
\put(235,19){$R$}
\put(100,37){$\overbrace{\rule{140pt}{0pt}}^
        {\mbox{series converges here}}$}
\put(10,16){$\underbrace{\rule{90pt}{0pt}}_
        {\mbox{diverges here}}$}
\put(240,16){$\underbrace{\rule{90pt}{0pt}}_
        {\mbox{diverges here}}$}
\put(170,-20){\makebox[0pt]
        {The radius of convergence of a power series}}
\end{picture}


An obvious example of the radius of convergence is given by
the geometric series
$$
        \dfrac{1}{1-x} = 1 + x + x^2 + x^3 + \cdots
$$
We know that this converges for $|x| < 1$ and diverges for $|x| > 1$.
%
\mar{Radius of convergence of the geometric series}
%
Thus the radius of convergence is $R = 1$ in this case.

It is possible for a power series to converge for all $x$; if that
happens, we take $R$ to be $\infty$.  At the other extreme, the series
may converge only for $x = 0$.  (When $x = 0$ the series collapses to
its constant term $a_0$, so it certainly converges {\em at least\/}
when $x = 0$.) If the series converges only for $x = 0$, then we take
$R$ to be 0.

At $x = R$ the series may diverge or converge; different things happen
for different series.  The same is true when $x = -R$.  The radius of
convergence tells us {\em where\/} the switch from convergence to
divergence happens.  It does not tell us {\em what\/} happens at the
place where the switch occurs.  If we know that the series converges
for $x=\pm R$, then we say that $[-R,R]$ is the interval of
convergence.  If the series converges when $x=R$ but {\em not\/} when
$x=-R$, then the interval of convergence is $(-R,R]$, and so on.


\subsection{The Ratio Test}

There are several ways to determine the radius of convergence of a
power series.  One of the simplest and most useful is by means of the
{\bf ratio test}.\index{ratio test} Because the power series to which
we apply this test need not include {\em consecutive} powers of $x$
(think of the Taylor series for $\cos x$ or $\sin x\/$) we'll write a
``generic" series as
$$
b_0 + b_1 + b_2 + \cdots = \sum_{m=0}^\infty b_m
$$ 
Here are three examples of the use of this notation.

\begin{enumerate}

\item  The Taylor series for $e^x$ is ${\displaystyle \sum_{m=0}^\infty
  b_m}$, where
$$
b_0 = 1, \quad  b_1 = x, \quad b_2 = \dfrac{x^2}{2!}, \cdots , 
   b_m = \dfrac{x^m}{m!}.
$$

\item  The Taylor series for $\cos x$ is ${\displaystyle
  \sum_{m=0}^\infty b_m}$, where
$$
b_0 = 1, \quad b_1 = \dfrac{-x^2}{2!}, \quad b_2 = \dfrac{x^4}{4!}, 
    \cdots , b_m = (-1)^m \dfrac{x^{2m}}{(2m)!} .
$$

\item  We can even describe the series
\[
17 + x + x^2 + x^4 + x^6 + x^8 + \cdots = 17 + x + \sum_{m=2}^\infty x^{2m-2}
\]
in our generic notation, in spite of the presence of the first two
terms ``$17 + x$'' which don't fit the pattern of later ones.  We have
$b_0 = 17$, $b_1 = x$, and then $b_m = x^{2m-2}$ for $m = 2$, 3,
4,~\ldots .  

\end{enumerate}

\newpage

The question of convergence for a power series is
%
\mar{Convergence is determined by the ``tail" of the series}
%
unaffected by the ``beginning" of the series; only the pattern in the
``tail" matters.  (Of course the {\em value\/} of the power series is
affected by all of its terms.)  So we can modify our generic notation
to fit the circumstances at hand.  No harm is done if we don't begin
with $b_0$.


Using this notation we can state the ratio test (but we give no proof).
\begin{center}
\begin{picture}(360,70)
        \begin{thicklines}
        \put(0,0){\framebox(360,70){\shortstack
                {\bf Ratio Test: the series \boldmath 
                        $b_0 + b_1 + b_2 + b_3 + \cdots + b_n + 
                        \cdots$\\[1pt]
                \bf converges if \boldmath ${\displaystyle 
                        \lim_{m \to \infty}\dfrac{|b_{m+1}|}{|b_m|} 
                        < 1.}$ 
                }}}
        \end{thicklines}
\end{picture}
\end{center}
Let's see 
%
\mar{The ratio test for the geometric series $\ldots$}
%
what the ratio test says about the geometric series:
$$
        1 + x + x^2 + x^3 + \cdots \ .
$$
We have $b_m = x^m$, so the ratio we must consider is
$$
\dfrac{|b_{m+1}|}{|b_m|} = \dfrac{|x^{m+1}|}{|x^m|} 
    = \dfrac{|x|^{m+1}}{|x|^m} = |x|.
$$ 
(Be sure you see why $|x^m| = |x|^m\/$.)  Obviously, this ratio has
    the same value for all $m$, so the limit
$$
\lim_{m \to \infty} |x| = |x|
$$
exists and is less than $1$ precisely when $|x| < 1$.  Thus the
geometric series converges for $|x| < 1$---which we already know is
true.  This means that the radius of convergence of the geometric
series is $R=1$.


Look next 
%
\mar{$\ldots$ and for $e^x$}
%
at the Taylor series for $e^x$:
$$
1 + x + \dfrac{x^2}{2!} + \dfrac{x^3}{3!} + \dfrac{x^4}{4!} + \cdots 
    = \sum_{m=0}^{\infty} \dfrac{x^m}{m!}.
$$
For negative $x$ this is an alternating series, so by the criterion
for convergence of alternating series we know it converges for all $x
< 0$.  The radius of convergence should then be $\infty$.  We will use
the ratio test to show that in fact this series converges for {\em
  all\/}~$x$.

In this case
$$
b_m = \dfrac{x^m}{m!},
$$
so the relevant ratio is
\[
\dfrac{|b_{m+1}|}{|b_m|} = 
  \left|\dfrac{x^{m+1}}{(m+1)!}\right| \cdot \left|\dfrac{m!}{x^m}\right| 
  =  \dfrac{|x^{m+1}|}{|x^m|} \cdot \dfrac{m!}{(m+1)!}
  =  |x| \cdot \dfrac{1}{m+1} = \dfrac{|x|}{m+1}.
\]
Unlike the example with the geometric series, the value of this ratio
depends on $m$.  For any particular $x$, as $m$ gets larger and larger
the numerator stays the same and the denominator grows, so this ratio
gets smaller and smaller.  In other words,
$$
\lim_{m \to \infty} \dfrac{|b_{m+1}|}{|b_m|} 
    = \lim_{m \to \infty} \dfrac{|x|}{m+1} = 0.
$$
Since this limit is less than 1 for any value of $x$, the series
converges for {\em all\/}~$x$, and thus the radius of convergence of
the Taylor series for $e^x$ is $R= \infty$, as we expected.


One of the uses of the theory developed so far is that it gives us a
new way of specifying functions.  For example, consider the power
series
$$
\sum_{m=0}^\infty (-1)^m\,\dfrac{2^m}{m^2+1}\,x^m = 
        1 - x + \dfrac{4}{5} x^2 - \dfrac{8}{10} x^3
        + \dfrac{16}{17} x^4 + \cdots .
$$
In this case $b_m=(-1)^m \dfrac{2^m}{m^2+1}x^m$, so to find the radius
of convergence, we compute the ratio
$$
\dfrac{|b_{m+1}|}{|b_m|} = 
        \dfrac{2^{m+1}|x|^{m+1}}{(m+1)^2 + 1} \cdot 
        \dfrac{m^2 + 1}{2^m |x|^m}  
      = 2|x| \dfrac{m^2 + 1}{m^2 + 2m + 2}.
$$
To figure out what happens to this ratio as $m$ grows large, 
%
\mar{Finding the limit of $|b_{m+1}|/|b_m|$ may require some algebra}
%
it is helpful to rewrite the factor involving the $m$'s as
$$
\dfrac{m^2 \cdot (1 + 1/m^2)}{m^2 \cdot (1+2/m +2/m^2)}
= \dfrac{1 + 1/m^2}{1 + 2/m + 2/m^2}.
$$
Now we can see that
$$
 \lim_{m \to \infty} \dfrac{|b_{m+1}|}{|b_m|} 
    = 2 |x| \cdot \dfrac{1}{1} = 2|x|.
$$
The limit value is less than $1$ precisely when $2|x| < 1\/$, or,
equivalently, $|x| < 1/2,$ so the radius of convergence of this series
is $R=1/2\/$.  It follows that for $|x| < 1/2$, we have a new function
$f(x)$ defined by the power series:
\[
f(x) = \sum_{m=0}^\infty (-1)^m \dfrac{2^m}{m^2+1} \,x^m .
\]


We can also discuss the radius of convergence of a power series
\[
a_0 + a_1 (x - a) + a_2 (x - a)^2 + \cdots + a_m (x - a)^m + \cdots 
\]
centered at a location $x = a$ other than the origin.  The radius of
convergence of a power series of this form can be found by the {\bf
ratio test} in exactly the same way it was when $a = 0$.
\medskip

\noindent 
{\bf Example.} Let's apply the ratio test to the Taylor series
centered at $a=1$ for $\ln(x)$:
$$
\ln(x) = \sum_{m=1}^\infty \dfrac{(-1)^{m-1}}{m} \, (x-1)^m.
$$
We can start our series with $b_1\/$, so we can take $b_m =
\dfrac{(-1)^{m-1}}{m} \, (x-1)^m\/$.  Then the ratio we must consider
is
$$
\dfrac{|b_{m+1}|}{|b_m|} = \dfrac{|x-1|^{m+1}}{m+1} \cdot \dfrac{m}{|x-1|^m} 
   = |x-1| \cdot \dfrac{m}{m+1} = |x-1| \cdot \dfrac{1}{1 + 1/m}.
$$
Then
$$
 \lim_{m \to \infty} \dfrac{|b_{m+1}|}{|b_m|} = |x-1| \cdot 1 = |x-1|.
$$
From this we conclude that this series converges for $|x-1| < 1$.
This inequality is equivalent to
$$
-1 < x-1 < 1,
$$ 
which is an interval of ``radius" $1$ about $a=1$, so the radius of
convergence is $R=1$ in this case.  We may also write the interval of
convergence for this power series as
$$
0 < x < 2.
$$

More generally, using the ratio test we find that a power series
centered at $a$ converges in an interval of ``radius" $R$ (and width
$2R$) around the point $x = a$ on the $x$-axis.  Ignoring what happens
at the endpoints, we say the {\bf interval of
  convergence\/}\index{interval of convergence} is
\[
        a - R < x < a + R.
\]
Here is a picture of what this looks like:
\begin{center}
\begin{picture}(360,90)(-10,-30)
\put(10,30){\vector(1,0){320}}
\put(100,30){\line(0,1){3}}
\put(170,30){\line(0,1){3}}
\put(240,30){\line(0,1){3}}
\put(334,25){$x$}
\put(87,19){$a-R$}
\put(167.5,19){$a$}
\put(228,19){$a+R$}
\put(100,37){$\overbrace{\rule{140pt}{0pt}}^
        {\mbox{convergence}}$}
\put(10,16){$\underbrace{\rule{90pt}{0pt}}_
        {\mbox{divergence}}$}
\put(240,16){$\underbrace{\rule{90pt}{0pt}}_
        {\mbox{divergence}}$}
\put(170,-20){\makebox[0pt]
        {The convergence of a power series centered at $a$}}
\end{picture}
\end{center}


\subsection{Exercises}

\num Find a formula for the sum of each of the following power series
by performing suitable operations on the geometric series and the
formula for {\em its\/} sum.

\medskip
\noindent
\begin{tabular}{ll}
a) $1-x^3 + x^6 - x^9 + \cdots .$ & d) $x+2x^2 +3x^3 + 4x^4 + \cdots .$ \\[3pt]
b) $x^2 + x^6 + x^{10} + x^{14} + \cdots .$ & e) $x +
  \dfrac{x^2}{2} + \dfrac{x^3}{3} + \dfrac{x^4}{4} + \cdots .$ \\ [7pt]
c) $1-2x+3x^2-4x^3 + \cdots .$ &
\end{tabular}


\num Determine the value of each of the following infinite sums.
(Each os these sums is a geometric or related series evaluated at a
particular value of $x$.)

\medskip
\noindent
\begin{tabular}{ll}
  a) $\dfrac{1}{4} + \dfrac{1}{16} + \dfrac{1}{64} + \dfrac{1}{256} + \cdots
  .$ & d) $\dfrac{1}{1} - \dfrac{2}{2} + \dfrac{3}{4} - \dfrac{4}{8} + \dfrac{5}{16} - \dfrac{6}{32} + \cdots .$ \\ [7pt]
b) .02020202 \ldots . & e)
  $\dfrac{1}{1\cdot 10} + \dfrac{1}{2\cdot10^2} + \dfrac{1}{3\cdot10^3} +
  \dfrac{1}{4\cdot10^4} + \cdots .$
\\ [7pt]
c) $- \dfrac{5}{2} + \dfrac{5}{4} - \dfrac{5}{8} + \cdots .$ & 
\end{tabular}


\num {\bf The Multiplier Effect}.  Economists know that the effect on
a local economy of tourist spending is greater than the amount
actually spent by the tourists.  The {\em multiplier effect}
quantifies this enlarged effect.  In this problem you will see that
calculating the multiplier effect involves summing a geometric series.

Suppose that, over the course of a year, tourists spend a total of $A$
dollars in a small resort town.  By the end of the year, the
townspeople are therefore $A$ dollars richer.  Some of this money
leaves the town---for example, to pay state and federal taxes or to
pay off debts owed to ``big city" banks.  Some of it stays in town but
gets put away as savings.  Finally, a certain fraction of the original
amount is spent in town, by the townspeople themselves.  Suppose
3/5-ths is spent this way.  The tourists and the townspeople {\em
  together} are therefore responsible for spending
$$
S = A + \dfrac{3}{5} A \quad \mbox{dollars}
$$
in the town that year.  The second amount---$\tfrac{3}{5}A$
dollars---is {\em recirculated} money.

Since one dollar looks much like another, the recirculated money
should be handled the same way as the original tourist dollars: some
will leave the town, some will be saved, and the remaining 3/5-ths
will get recirculated a {\em second} time.  The twice-recirculated
amount is
$$
\dfrac{3}{5} \times \dfrac{3}{5}A \quad \mbox{dollars},
$$
and we must revise the calculation of the total amount spent in the
town to
$$
S = A + \dfrac{3}{5}A + \left(\dfrac{3}{5}\right)^2\!\! A \quad \mbox{dollars}.
$$
But the twice-recirculated dollars look like all the others , so
3/5-ths of them will get recirculated a {\em third} time.  Revising
the total dollars spent yet again, we get
$$
S = A + \dfrac{3}{5}A + \left(\dfrac{3}{5}\right)^2\!\! A + 
\left(\dfrac{3}{5}\right)^3\!\! A \quad \mbox{dollars}.
$$
This process never ends: no matter how many times a sum of money has
been recirculated, 3/5-ths of it is recirculated once more.  The total
amount spend in the town is thus given by a {\em series.}

\sub Write the series giving the total amount of money spent in the
town and calculate its sum.

\sub Your answer in a) is a certain multiple of $A$---what is the
multiplier?

\sub Suppose the recirculation rate is $r$ instead of 3/5.  Write the
series giving the total amount spent and calculate its sum.  What is
the multiplier now?

\sub Suppose the recirculation rate is 1/5; what is the multiplier in
this case?

\sub Does a lower recirculation rate produce a smaller multiplier
effect?

\newpage

\num Which of the following alternating series converge, which diverge? Why? \\

\begin{minipage}{2.4in}
\sub ${\displaystyle \sum_{n=1}^{\infty} (-1)^n\dfrac{n}{n^2+1}}$
\sub ${\displaystyle \sum_{n=1}^{\infty} (-1)^n\dfrac{1}{\sqrt{3n+2}}}$
\sub ${\displaystyle \sum_{n=2}^{\infty} (-1)^n\dfrac{n}{\ln{n}}}$
\sub ${\displaystyle \sum_{n=1}^{\infty} (-1)^n\dfrac{n}{5n-4}}$
\sub ${\displaystyle \sum_{n=1}^{\infty} (-1)^n\dfrac{\arctan{n}}{n}}$
\end{minipage} \ \
\begin{minipage}{2.4in}
\sub ${\displaystyle \sum_{n=1}^{\infty} (-1)^n\dfrac{(1.0001)^n}{n^{10}+1}}$
\sub ${\displaystyle \sum_{n=1}^{\infty} (-1)^n\dfrac{1}{n^{1/n}}}$
\sub ${\displaystyle \sum_{n=2}^{\infty} (-1)^n\dfrac{1}{\ln{n}}}$
\sub ${\displaystyle \sum_{n=1}^{\infty} (-1)^n\dfrac{n!}{n^n}}$
\sub ${\displaystyle \sum_{n=1}^{\infty} (-1)^n\dfrac{n!}{1\cdot 3\cdot 5\cdots (2n-1)}}$
\end{minipage}

\num For each of the sums in the preceding problem that converges, use
the alternating series criterion to determine how far out you have to
go before the sum is determined to 6 decimal places.  Give the sum for
each of these series to this many places.

\num Find a value for $n$ so that the $n$th degree Taylor series for
$e^x$ gives at least 10 place accuracy for all $x$ in the interval
$[-3, 0]$.

\num We defined the harmonic series as the infinite sum
$$
1 + \dfrac{1}{2} + \dfrac{1}{3} + \dfrac{1}{4} + \cdots 
   = \sum_{i=1}^\infty \dfrac{1}{i} .
$$

\sub Use a calculator to find the partial sums
$$
S_n = 1 + \dfrac{1}{2} + \dfrac{1}{3} + \dfrac{1}{4} + \cdots \dfrac{1}{n}
$$
for  $n = 1$, 2, 3, \ldots , 12.

\sub Use the following program to find the value of $S_n$ for $n =
100$.  Modify the program to find the values of $S_n$ for $n = 500$,
1000, and $5000$.

\newpage

\index{program!HARMONIC}
\begin{center}
{\bf Program: HARMONIC}
\end{center}
\begin{verbatim}
  n = 100
  sum = 0
  FOR i = 1 TO n
       sum = sum + 1/i
  NEXT i
  PRINT  n, sum
\end{verbatim}
\sub Group the terms in the harmonic series as indicated by the parentheses:
\begin{gather*}%multline*}
1 + \left(\dfrac{1}{2}\right) + 
                \left(\dfrac{1}{3} + \dfrac{1}{4}\right) + 
                \left(\dfrac{1}{5} + \dfrac{1}{6} + \dfrac{1}{7} + 
                      \dfrac{1}{8}\right) +  \qquad\\ 
\qquad \mbox{} + \left(\dfrac{1}{9} + \cdots + 
                \dfrac{1}{16}\right) + \left(\dfrac{1}{17} + 
                \cdots + \dfrac{1}{32}\right) + \cdots .
\end{gather*}%multline*}
Explain why each parenthetical grouping totals at least $1/2$.

\sub Following the pattern in part (c), if you add up the terms of the
harmonic series forming $S_n$ for $n=2^k$, you can arrange the terms
as $1 + k$ such groupings.  Use this fact and the result of c) to
explain why $S_n$ exceeds $1+k\cdot \tfrac{1}{2}$.

\sub Use part (d) to explain why the harmonic series {\em diverges}.

\sub You might try this problem if you've studied physics---enough to
know how to locate the center of mass of a system. Suppose you had $n$
cards and wanted to stack them on the edge of a table with the top of
the pile leaning out over the edge. How far out could you get the pile
to reach if you were careful?  Let's choose our units so the length of
each card is 1. Clearly if $n=1$, the farthest reach you could would
be $\frac{1}{2}$.  If $n=2$, you could clearly \linebreak

\vspace{-9pt}

\special{psfile=../ps/calc2/cards.ps} 
\hfill
\begin{minipage}{3in}
place the top card to extend half a unit beyond the bottom
card.  For the system to be stable, the center of mass of the two
cards must be to the left of the edge of the table. Show that for this
to happen, the bottom card can't extend more than 1/4 unit beyond the
edge.  Thus with $n=2$\,, the maximum extension of the pile is
$\frac{1}{2} + \frac{1}{4} = \frac{3}{4}$.  The picture at the right
shows 10 cards stacked carefully.
\end{minipage}

\noindent
Prove that if you have $n$ cards, the stack can be built to extend a
distance of
\[
\frac{1}{2}+\frac{1}{4}+\frac{1}{6}+\frac{1}{8}+\cdots +\frac{1}{2n}=\frac{1}{2}\left(1+\frac{1}{2}+\frac{1}{3}+\frac{1}{4}+\cdots +\frac{1}{n}\right).
\]

In the light of what we have just proved about the harmonic series, this
shows that if you had enough cards, you could have the top card
extending 100 units beyond the edge of the table!

\num {\bf An estimate for the partial sums of the harmonic
  series}.\index{harmonic series!logarithmic approximation} \quad You
may notice in part (d) of the preceding problem that the sum $S_n = 1
+ 1/2 + 1/3 + \cdots + 1/n$ grows in proportion to the exponent $k$ of
$n=2^k$; i.e., the sum grows like the logarithm of~$n$.  We can make
this more precise by comparing the value of $S_n$ to the value of the
integral
$$
\int_{1}^{n} \dfrac{1}{x} \, dx = \ln(n).
$$

\sub Let's look at the case $n=6$ to see what's going on.  Consider
the following picture:

\vspace*{130pt}

\special{psfile=../ps/calc2/harmonic.ps}

Show that the lightly shaded region plus the dark region has area
equal to $S_5$\,, which can be rewritten as $S_6 - \tfrac{1}{6}$.
Show that the dark region alone has area $S_6 - 1$. Hence prove that
\[
S_6 - 1 < \int_{1}^{6} \dfrac{1}{x} \, dx < S_6 - \dfrac{1}{6},
\]
and conclude that
\[
\dfrac{1}{6}<S_6 - \ln(6)<1.
\]

\sub Show more generally that \quad $\dfrac{1}{n}<S_n-\ln(n)<1$.

\sub Use part (b) to get upper and lower bounds for the value of $S_{10000}$.

\ans{$9.21044 < S_{10000} < 10.21035$.}

\sub Use the result of part (b) to get an estimate for how many cards
you would need in part (f) of the preceding problem to make the top of
the pile extend 100 units beyond the edge of the table.

\ans{It would take approximately $10^{87}$ cards---a number which is
  on the same order of magnitude as the number of atoms in the
  universe!}

\medskip

\aside{Remarkably, partial sums of the harmonic series exceed $\ln(n)$
  in a very regular way.  It turns out that
        $$
        \lim_{n \rightarrow \infty} \, 
                \left\{ S_n - \ln(n) \right\} = \gamma,
        $$
where $\gamma = .5772\ldots$ is called {\bf Euler's
  constant.}\index{Euler's constant} (You have seen another constant
named for Euler, the base $e= 2.7183\ldots $.)  Although one can find
the decimal expansion of $\gamma$ to any desired degree of accuracy,
no one knows whether $\gamma$ is a rational number or not.}

\num Show that the power series for $\arctan{x}$ and $(1+x)^c$ diverge
for $|x|>1 \,$.  Do the series converge or diverge when $|x|=1$?


\num  Find the radius of convergence of each of the following power series.

\sub $1 + 2 x + 3 x^2 + 4 x^3 + \cdots$.

\sub $x + 2 x^2 + 3 x^3 + 4 x^4 + \cdots$.

\sub $1 + \dfrac{1}{1^2} x + \dfrac{1}{2^2} x^2 + \dfrac{1}{3^2} x^3 +
     \dfrac{1}{4^2} x^4 + \cdots$. \hfill [Answer: $R = 1$]

\sub  $x^3 + x^6 + x^9 + x^{12} + \cdots$.

\sub $1 + (x + 1) + (x+1)^2 + (x+1)^3 + \cdots$.

\sub $17 + \dfrac{1}{3} x + \dfrac{1}{3^2} x^2 + \dfrac{1}{3^3} x^3 +
\dfrac{1}{3^4} x^4 + \cdots$.


\num Write out the first five terms of each of the following power
series, and determine the radius of convergence of each.

\sub  ${\displaystyle \sum_{n = 0}^{\infty}\, n x^n}$.
\hfill [Answer: $R = 1$]

\sub ${\displaystyle \sum_{n = 0}^{\infty} \dfrac{n^2}{2^n} x^n}$. \hfill
     [Answer: $R = 2$]

\sub ${\displaystyle \sum_{n = 0}^{\infty} (n+5)^2\, x^n}$. \hfill
     [Answer: $R = 1$]

\sub  ${\displaystyle \sum_{n = 0}^{\infty} 
\dfrac{99}{n^n} x^n}$. \hfill [Answer: $R = \infty$]

\sub  ${\displaystyle \sum_{n = 0}^{\infty} n!\, x^n}$.
\hfill [Answer: $R = 0$]

\num Find the radius of convergence of the Taylor series for $\sin{x}$
and for $\cos{x}$.  For which values of $x$ can these series therefore
represent the sine and cosine functions?

\num Find the radius of convergence of the Taylor series for $f(x) =
1/(1 + x^2)$ at $x = 0$.  (See the table of Taylor series in section~3.)
What is the radius of convergence of this series?  For which values of
$x$ can this series therefore represent the function $f$?  Do these
$x$ values constitute the {\em entire\/} domain of definition of $f$?

\num In the text we used the alternating series for $e^x$, $\; x < 0$,
to approximate $e^{-1}$ accurate to 7 decimal places.  The claim was
made that in taking the reciprocal to obtain an estimate for $e$, the
accuracy drops by one decimal place.  In this problem you will see why
this is true.

\sub Consider first the more general situation where two functions are
reciprocals, $g(x) = 1/f(x)$.  Express $g'(x)$ in terms of $f(x)$ and
$f'(x)$.

\sub Use your answer in part a) to find an expression for the {\em
  relative error} in $g$, $\Delta g/g(x) \approx g'(x) \Delta x/g(x)$,
in terms of $f(x)$ and $f'(x)$.  How does this compare to the relative
error in $f$?

\sub Apply your results in part b) to the functions $e^x$ and $e^{-x}$
at $x=1$.  Since $e$ is about 7 times as large as $1/e$, explain why
the error in the estimate for $e$ should be about 7 times as large as
the error in the estimate for~$1/e$.


\newpage
%%%%%%%%%%%%%%%%%%%%%%%%%%%%%%%%%%%%%%%%%%%%%%%%%%%%%%%%%%%%%%%%%%%%%%%%
%                                                                      %
%                               Section 6                              %
%                                                                      %
%%%%%%%%%%%%%%%%%%%%%%%%%%%%%%%%%%%%%%%%%%%%%%%%%%%%%%%%%%%%%%%%%%%%%%%%


\section{Approximation Over Intervals}

A powerful result in mathematical analysis is the {\bf
  Stone--Weierstrass Theorem}\index{Stone--Weierstrass Theorem}, which
states that given any continuous function $f(x)$ and any
interval $[a, b]$, there exist polynomials that fit $f$ over this
interval to any level of accuracy we care to specify. In many cases,
we can find such a polynomial simply by taking a Taylor polynomial of
high enough degree.  There are several ways in which this is not a
completely satisfactory response, however. First, some functions (like
the absolute value function) have corners or other places where they
aren't differentiable, so we can't even build a Taylor series at such
points. Second, we have seen several functions (like $1/(1+x^2)$)
that have a finite interval of convergence, so Taylor polynomials may
not be good fits no matter how high a degree we try.  Third, even for
well-behaved functions like $\sin(x)$ or $e^x$, we may have to take a
very high degree Taylor polynomial to get the same overall fit that a
much lower degree polynomial could achieve.

In this section we will develop the general machinery for finding
polynomial approximations to functions over given intervals.  In
chapter 12.4 we will see how this same approach can be adapted to
approximating periodic functions by {\bf trigonometric polynomials}.

\subsection{Approximation by polynomials}

{\bf Example}.  Let's return to the problem introduced at the
beginning of this chapter: find the second degree polynomial which
best fits the function $\sin(x)$ over the interval $[0, \pi]$.
%
\mar{Two possible criteria for best fit}
%
Just as we did with the Taylor polynomials, though, before we can
start we need to agree on our criterion for the best fit. Here are two
obvious candidates for such a criterion:

\begin{enumerate}

\item The second degree polynomial $Q(x)$ is the best fit to $\sin(x)$
  over the interval $[0, \pi]$ if the {\em maximum} separation between
  $Q(x)$ and $\sin(x)$ is smaller than the maximum separation between
  $\sin(x)$ and any other second degree polynomial:
\[
\max_{0\le x\le \pi}|\sin(x)-Q(x)| \qquad \text{is the smallest possible}.
\]

\item The second degree polynomial $Q(x)$ is the best fit to $\sin(x)$
  over the interval $[0, \pi]$ if the {\em average} separation between
  $Q(x)$ and $\sin(x)$ is smaller than the average separation between
  $\sin(x)$ and any other second degree polynomial:
\[
\dfrac{1}{\pi }\int_0^{\pi} |\sin(x)-Q(x)|\, dx \qquad
     \text{is the smallest possible}.
\]

\end {enumerate}
Unfortunately, even though their clear meanings make these two
criteria very attractive, they turn out to be largely \mar{Why we
  don't use either criterion} unusable---if we try to apply either
criterion to a specific problem, including our current example, we are
led into a maze of tedious and unwieldy calculations.

Instead, we use a criterion that, while slightly less obvious than
either of the two we've already articulated, still clearly measures
some sort of ``best fit'' and has the added virtue of behaving well
mathematically.  We accomplish this by modifying criterion 2 slightly.
It turns out that the major difficulty with this criterion is the
presence of absolute values.  If, instead of considering the average
separation between $Q(x)$ and $\sin(x)$, we consider the average of
the {\em square} of the separation between $Q(x)$ and $\sin(x)$, we
get a criterion we can work with.  (Compare this with the discussion
of the best-fitting line in the exercises for chapter~9.3.)  Since
this is a definition we will be using for the rest of this section, we
frame it in terms of arbitrary functions $g$ and $h$, and an arbitrary
interval $[a, b]$:

\noindent
\begin{center}
\framebox[4.8in] {\parbox{4.45in}{\vspace{8pt} 
Given two functions $g$ and $h$ defined over an interval $[a, b]$, 
we define the {\bf mean square separation} between $g$ and $h$ over 
this interval to be
\[
\dfrac{1}{(b-a)} \int_a^b \left(g(x) -h(x)\right)^2 \, dx.
\]}
}

\end{center}\index{mean square separation}
Note: In this setting the word {\bf mean}\index{mean} is synonymous
with what we have called ``average''. It turns out that there is often
more than one way to define the term ``average''---the concepts of
median and mode are two other natural ways of capturing
``averageness'', for instance---so we use the more technical term to
avoid ambiguity.

We can now rephrase our original 
%
\mar{The criterion\\ we shall use}
%
problem as: find the second degree polynomial $Q(x)$ whose mean
squared separation from $\sin(x)$ over the interval $[0, \pi]$ is as
small as possible.  In mathematical terms, we want to find
coefficients $a_0$, $a_1$, and $a_2$ which such that the integral
\[
\int_0^{\pi} \left( \sin(x) - (a_0 + a_1\,x + a_2\,x^2)\right)^2 \, dx
\]
is minimized. The solution $Q(x)$ is called the quadratic {\bf least
  squares approximation}\index{approximation!least squares} to
$\sin(x)$ over $[0,\pi]$.

The key to solving this problem is to observe that $a_0$, $a_1$, and $a_2$ can take on any values we like and that this integral can thus be considered a function of these three variables. For instance, if we couldn't think of anything cleverer to do, we might simply try various combinations of $a_0$, $a_1$, and $a_2$ to see how small we could make the given integral.  Therefore another way to phrase our problem is
\mar{\vspace{24pt}A mathematical formulation of the problem}

\noindent
\begin{center}
\framebox[4.8in]
{\parbox{4.45in}{\vspace{8pt}
Find values for $a_0$, $a_1$, and $a_2$ that minimize the function
$$F(a_0, a_1, a_2) = \int_0^{\pi} \left( \sin(x) - (a_0 + a_1\,x + a_2\,x^2)\right)^2 \, dx \, .$$\vspace{-8pt}}
}
\end{center}

\medskip

We know how to find points where functions take on their extreme
values---we look for the places where the partial derivatives are
0. But how do we differentiate an expression involving an integral
like this?  It turns out that for all continuous functions, or even
functions with only a finite number of breaks in them, we can simply
interchange integration and differentiation.  Thus, in our example,
\begin{align*}
\dfrac{\partial}{\partial a_0}F(a_0, a_1, a_2)
  &= \dfrac{\partial}{\partial a_0}\int_0^{\pi} 
     \left( \sin(x) - (a_0 + a_1\,x + a_2\,x^2)\right)^2 \, dx \\[.15in]
  &= \int_0^{\pi} \dfrac{\partial}{\partial a_0}
     \left( \sin(x) - (a_0 + a_1\,x + a_2\,x^2)\right)^2 \, dx \\[.15in]
  &= \int_0^{\pi}2\left( \sin(x) - (a_0 + a_1\,x + a_2\,x^2)\right)(-1) \, dx.
\end{align*}
Similarly we have
\begin{align*}
\dfrac{\partial}{\partial a_1}F(a_0, a_1, a_2)
  &= \int_0^{\pi}2\left( \sin(x) 
            - (a_0 + a_1\,x + a_2\,x^2)\right)(-x) \, dx, \\[.15in]
\dfrac{\partial}{\partial a_2}F(a_0, a_1, a_2)
  &= \int_0^{\pi}2\left( \sin(x) 
             - (a_0 + a_1\,x + a_2\,x^2)\right)(-x^2) \, dx.
\end{align*}

We now want 
        \mar{Setting the partials equal to zero gives equations for $a_0, a_1, a_2$}
to find values for $a_0$, $a_1$, and $a_2$ that make these partial derivatives simultaneously equal to 0.  That is, we want
\begin{align*}
\int_0^{\pi}2\left( \sin(x) - (a_0 + a_1\,x + a_2\,x^2)\right)(-1) \, dx 
   &= 0, \\[.15in]
\int_0^{\pi}2\left( \sin(x) - (a_0 + a_1\,x + a_2\,x^2)\right)(-x) \, dx 
   &= 0, \\[.15in]
\int_0^{\pi}2\left( \sin(x) - (a_0 + a_1\,x + a_2\,x^2)\right)(-x^2) \, dx 
   &= 0,
\end{align*}
which can be rewritten as
\begin{align*}
\int_0^{\pi}\sin(x) \, dx
  &= \int_0^{\pi}\left( a_0 + a_1\,x + a_2\,x^2\right)\, dx,  \\[.15in]
\int_0^{\pi}x\,\sin(x) \, dx
  &= \int_0^{\pi}\left(a_0\,x + a_1\,x^2 + a_2\,x^3\right) \, dx,  \\[.15in]
\int_0^{\pi}x^2\,\sin(x) \, dx
  &= \int_0^{\pi}\left( a_0\,x^2 + a_1\,x^3 + a_2\,x^4\right) \, dx.
\end{align*}

All of these integrals 
        \mar{Evaluating the integrals gives three linear equations}
can be evaluated relatively easily (see the exercises for a hint on evaluating the integrals on the left--hand side).  When we do so, we are left with 
\begin{align*}
        2 &= \pi a_0 + \dfrac{\pi^2}{2}a_1 + \dfrac{\pi^3}{3}a_2,  \\[3pt]
      \pi &=\dfrac{\pi^2}{2}a_0 + \dfrac{\pi^3}{3}a_1 + 
             \dfrac{\pi^4}{4}a_2, \\[5pt]
\pi^2 - 4 &= \dfrac{\pi^3}{3}a_0 + \dfrac{\pi^4}{4}a_1 
             + \dfrac{\pi^5}{5}a_2. 
\end{align*}

But this is simply a set of three linear equations in the unknowns
$a_0$, $a_1$, and $a_2$, and they can be solved in the usual ways.  We
could either replace each expression in $\pi$ by a corresponding
decimal approximation, or we could keep everything in terms of $\pi$.
Let's do the latter; after a bit of tedious arithmetic we find
\addtolength{\tabcolsep}{-3pt}%

\begin{center}
\begin{tabular}{rclcl}
$a_0 $&=& $\dfrac{12}{\pi}-\dfrac{120}{\pi^3}$&=& $-.050465\ldots ,$ \\ [.15in]
$a_1 $&=& $\dfrac{-60}{\pi^2}+\dfrac{720}{\pi^4}$&=& $1.312236\ldots ,$ \\ [.15in]
$a_2 $&=& $\dfrac{60}{\pi^3}-\dfrac{720}{\pi^5}$&=& $-.417697\ldots ,$ \\ [.15in]
\end{tabular}
\end{center}
\addtolength{\tabcolsep}{3pt}%
and we have
\[
Q(x) = -.050465 +1.312236\,x-.417698\,x^2,
\]
which is the equation given in section~1 at the beginning of the chapter.

The analysis 
%
\mar{\vspace{50pt}How to find least squares polynomial approximations in general}
%
we gave for this particular case can clearly be generalized to apply
to any function over any interval.  When we do this we get
\begin{center}
\framebox[5.4in]
{\vspace{4pt}
\parbox{4.6in}{\vspace{8pt}
Given a function $g$ over an interval $[a, b]$, then the $n$-th degree
polynomial
\[
P(x) = c_0 +c_1\,x+c_2\,x^2 + \cdots +c_n\,x^n
\]
whose mean square distance from $g$ is a minimum has coefficients that
are determined by the following $n+1$ equations in the $n+1$ unknowns
$c_0$, $c_1$, $c_2$, \ldots , $c_n$:
\begin{align*}
\int_a^b g(x)\,dx 
  &= c_0\int_a^b dx +c_1\int_a^b x\,dx  + \cdots 
            + c_n\int_a^b x^n\,dx, \\[6pt]
\int_a^b x\,g(x)\,dx 
  &= c_0\int_a^b x\,dx +c_1\int_a^b x^2\,dx + \cdots 
             + c_n\int_a^b x^{n+1}\,dx, \\[-6pt]
&\ \, \vdots \\[-6pt]
\int_a^b x^n\,g(x)\,dx 
  &= c_0\int_a^b x^n\,dx +c_1\int_a^b x^{n+1}\,dx  + \cdots 
             + c_n\int_a^b x^{2n}\,dx.
\end{align*}
 }
}
\end{center}\index{least squares polynomial approximation, general rule}

\medskip

All the integrals on the right-hand side can be evaluated immediately.
The integrals on the left-hand side will typically need to be
evaluated numerically, although simple cases can be evaluated in
closed form.  Integration by parts is often useful in these cases.
The exercises contain several problems using this technique to find
approximating polynomials.

\newpage

The real catch, though, is not in obtaining the 
%
\mar{Solving the equations is a job for the computer}
%
equations---it is that solving systems of equations by hand is
excruciatingly boring and subject to frequent arithmetic mistakes if
there are more than two or three unknowns involved.  Fortunately,
there are now a number of computer packages available which do all of
this for us.  Here are a couple of examples, where the details are
left to the exercises.

\medskip
\noindent
{\bf Example}.  Let's find polynomial approximation for
\mbox{$1/(1+x^2)$} over the interval $[0, 2]$.  We saw earlier that
the Taylor series for this function converges only for $|x|<1$, so it
will be no help.  Yet with the above technique we can derive the
following approximations of various degrees (see the exercises for
details):

\vspace{190pt}

\special{psfile=../ps/calc2/leastsq1.ps}

Here are the corresponding equations of the approximating polynomials:
\begin{center}
\begin{tabular}{cl}
degree & \hspace{1in} polynomial \\[2pt] 
\hline  \\ [-8pt] 
1 & $1.00722 - .453645\,x $ \\[3pt]
2& $1.08789 -.695660\,x +.121008\,x^2 $\\[3pt]
3 & $1.04245 -.423017\,x-.219797\,x^2+.113602\,x^3$ \\[3pt]
4 & $1.00704 -.068906\,x-1.01653\,x^2+.733272\,x^3-.154916\,x^4$
\end{tabular}
\end{center}

\noindent
{\bf Example}.  We can even use this new technique to find 
%
\mar{The technique works\\ even when differentiability fails}
%
polynomial approximations for functions that aren't differentiable at
some points.  For instance, let's approximate the function $h(x) =
|x|$ over the interval $[-1,1]$.  Since this function is symmetric
about the $y$--axis, and we are approximating it over an interval that
is symmetric about the $y$--axis, only even powers of $x$ will appear.
(See the exercises for details.)  We get the following approximations
of degrees 2, 4, 6, and~8:

\vspace{220pt}

\special{psfile=../ps/calc2/leastsq2.ps}

\medskip
Here are the corresponding polynomials:
\begin{center}
\begin{tabular}{cl}
degree & \hspace{1in} polynomial \\[2pt] 
\hline \\ [-8pt]
2 & $.1875 +.9375\,x^2 $ \\[3pt]
4& $.117188 +1.64062\,x^2 -.820312\,x^4 $\\[3pt]
6 & $.085449 +2.30713\,x^2-2.81982\,x^4+1.46631\,x^6$ \\[3pt]
8 & $.067291 +2.960821\,x^2-6.415132\,x^4+7.698173\,x^6-3.338498\,x^8$
\end{tabular}
\end{center}

\medskip
\noindent
{\bf A Numerical Example}.  If we have some function 
%
\mar{The technique is useful for data functions}
%
which exists only as a set of data points---a numerical solution to a
differential equation, perhaps, or the output of some laboratory
instrument---it can often be quite useful to replace the function by
an approximating polynomial.  The polynomial takes up much less
storage space and is easier to manipulate.  To see how this works,
let's return to the $S$-$I$-$R$ model we've studied before
\begin{align*}
S'&= -.00001\, SI, \\
I'&= .00001\, SI -  I/14, \\
R'&= I/14,
\end{align*}
with initial values $S(0)=45400$ and $I(0)=2100$.

Let's find an 8th degree polynomial $Q(t)=i_0+i_1t+i_2t^2+\cdots
+i_8t^8$ approximating $I$ over the time interval $0\le t\le 40$.  We
can do this by a minor modification of the Euler's method programs
we've been using all along.  Now, in addition to keeping track of the
current values for $S$ and $I$ as we go along, we will also need to be
calculating Riemann sums for the integrals
\[
\int_0^{40} t^kI(t)\,dt \qquad \text{for $k =0$, 1, 2, \ldots , 8},
\]
as we go through each iteration of Euler's method.  

Since the numbers involved become enormous very quickly, we open
ourselves to various sorts of computer
%
\mar{The importance\\ of using the\\ right-sized units}
%
roundoff error.  We can avoid some of these difficulties by {\bf rescaling}\index{rescaling} our equations---using units that keep the numbers involved more manageable.  Thus, for instance, suppose we measure $S$, $I$, and $R$ in units of 10,000 people, and suppose we measure time in ``decadays'', where 1 decaday = 10 days.  When we do this, our original differential equations become
\begin{align*}
S'&= -SI, \\
I'&= SI -  I/1.4, \\
R'&= I/1.4,
\end{align*}
with initial values $S(0)=4.54$ and $I(0)=0.21$.  The integrals we
want are now of the form
\[
\int_0^{4} t^kI(t)\,dt \qquad \text{for $k = 0$, 1, 2, \ldots , 8}.
\]

The use of Simpson's rule\index{Simpson's rule} (see chapter~11.3)
will also reduce errors.
%
\mar{Using Simpson's rule helps reduce errors}
%
It may be easiest to calculate the values of $I$ first, using perhaps
2000 values, and store them in an array.  Once you have this array of
$I$ values, it is relatively quick and easy to use Simpson's rule to
calculate the 9 integrals needed.  If you later decide you want to get
a higher-degree polynomial approximation, you don't have to re-run the
program.

Once we've evaluated these integrals, we set up and solve the
corresponding system of 9 equations in the 9 unknown coefficients
$i_k$. We get the following 8-th degree approximation
\begin{align*}
Q(t) &= .3090 -.9989\, t + 7.8518\,t^2 - 3.6233\,t^3 - 3.9248\,t^4 
            + 4.2162\,t^5 \\
& \mbox{}\quad  -1.5750\,t^6 +.2692\,t^7-.01772\,t^8 \, .
\end{align*}

When we graph $Q$ and $I$ together over the interval $[0, 4]$
(decadays), we get

\vspace{180pt}

\special{psfile=../ps/calc2/sireuler.ps}

\noindent
---a reasonably good fit.

\medskip
\noindent
{\bf A Caution}.  Numerical least-squares fitting of the sort
performed in this last example fairly quickly pushes into regions
where the cumulative effects of the inaccuracies of the original data,
the inaccuracies of the estimates for the integrals, and the immense
range in the magnitude of the numbers involved all combine to produce
answers that are obviously wrong. Rescaling the equations and using
good approximations for the integrals can help put off the point at
which this begins to happen.


\subsection{Exercises}


\num To find polynomial approximations for $\sin(x)$ over the interval
$[0,\pi]$, we needed to be able to evaluate integrals of the form
\[
\int_0^{\pi} x^n \sin(x)\,dx.
\]
The value of this integral clearly depends on the value of $n$, so
denote it by~$I_n$.

\enlargethispage*{20pt}
\sub Evaluate $I_0$ and $I_1$.  Suggestion: use integration by parts
(Chapter~11.3) to evaluate~$I_1$.

\ans{$I_0=2$, and $I_1=\pi$.}

\sub Use integration by parts twice to prove the general {\bf
  reduction formula}\index{reduction formula}:
\[
I_{n+2} = \pi^{n+2}-(n+2)(n+1)I_n \qquad \mbox{for all $n \ge 0$}.
\]

\newpage

\sub Evaluate $I_2$, $I_3$, and $I_4$.

\ans{$I_2=\pi^2-4$, $I_3=\pi^3-6\pi$, and $I_4=\pi^4-12\pi^2+48$.}


\sub If you have access to a computer package that will solve a system
of equations, find the 4-th degree polynomial that best fits the sine
function over the interval $[0,\pi]$.  What is the maximum difference
between this polynomial and the sine function over this interval?

\ans{$.00131 +.98260\,x+ .05447\,x^2-.23379\,x^3+.03721\,x^4$, with
  maximum difference occuring at the endpoints.}

\num To find polynomial approximations for $|x|$ over the interval
$[-1,1]$, we needed to be able to evaluate integrals of the form
\[
\int_{-1}^{1} x^n |x|\,dx.
\]
As before, let's denote this integral by $I_n$.

\sub Show that
\[
I_n  = \begin{cases}
\dfrac{2}{n+2} &\text{if $n$ is even}, \\[3pt]
0              &\text{if $n$ is odd}.
\end{cases}
\]

\sub Derive the quadratic least squares approximation to $|x|$ over
$[-1,1]$.

\sub If you have access to a computer package that will solve a system
of equations, find the 10-th degree polynomial that best fits $|x|$
over the interval $[-1,1]$.  What is the maximum difference between
this polynomial and $|x|$ over this interval?

\num To find polynomial approximations for $1/(1+x^2)$ over the
interval $[0,2]$, we needed to be able to evaluate integrals of the
form
\[
\int_0^2 \dfrac{x^n}{1+x^2}\,dx.
\]
Call this integral $I_n$.

\sub Evaluate $I_0$ and $I_1$.

\ans{$I_0 = \arctan(1)=\pi/4$, and $I_1=(\ln{2})/2=.3465736$.}

\sub Prove the general reduction formula:
\[
I_{n+2} = \dfrac{2^{n+1}}{n+1} -I_n \qquad \text{for $n = 0$, 1, 2, \ldots }.
\]

\sub Evaluate $I_2$, $I_3$, and $I_4$.

\ans{$I_2=2-\dfrac{\pi}{4}, I_3=2-\dfrac{\ln{2}}{2}, I_4= \dfrac{2}{3}
  + \dfrac{\pi}{4}$.}

\sub If you have access to a computer package that will solve a system
of equations, find the 4-th degree polynomial that best fits $1/(1 +
x^2)$ over the interval $[0, 2]$.  What is the maximum difference
between this polynomial and the function over this interval?

\num Set up the equations (including evaluating all the integrals) for
finding the best fitting 6-th degree polynomial approximation to
$\sin(x)$ over the interval $[-\pi, \pi]$.

\num In the $S$-$I$-$R$ model, find the best fitting 8-th degree
polynomial approximation to $S(t)$ over the interval $0\le t\le 40$.


\section{Chapter Summary}

\subsection{The Main Ideas}

\begin{itemize}

\item {\bf Taylor polynomials} approximate functions at a point.  The
  Taylor polynomial $P(x)$ of degree $n$ is the {\bf best fit} to
  $f(x)$ at $x=a$; that is, $P$ satisfies the following conditions:
  $P(a)=f(a)$, $P'(a)=f'(a)$, $P''(a)=f''(a)$, \ldots ,
  $P^{(n)}(a)=f^{(n)}(a)$.

\item {\bf Taylor's theorem} says that a function and its Taylor
  polynomial of degree $n$ agree to order $n + 1$ near the point where
  the polynomial is centered.  Different versions expand on this idea.

\item If $P(x)$ is the Taylor polynomial approximating $f(x)$ at
  $x=a$, then $P(x)$ approximates $f(x)$ for values of $x$ near $a$;
  $P'(x)$ approximates $f'(x)$; and $\int \! P(x) \, dx$ approximates
  $\int \! f(x) \, dx$.

\item A {\bf Taylor series} is an infinite sum whose partial sums are
  Taylor polynomials.  Some functions {\em equal} their Taylor series;
  among these are the sine, cosine and exponential functions.

\item  A {\bf power series} is an ``infinite" polynomial
\[
a_0 + a_1x + a_x x^2 + \cdots + a_n x^n + \cdots .
\]  
If the solution of a differential equation can be represented by a
power series, the coefficients $a_n$ can be determined by {\bf
  recursion relations} obtained by substituting the power series into
the differential equation.


\item An infinite series {\bf converges} if, no matter how many
  decimal places are specified, all the partial sums eventually agree
  to at least this many decimal places.  The number defined by these
  stabilizing decimals is called the {\bf sum} of the series.  If a
  series does not converge, we say it {\bf diverges}.

\item If the series $\sum_{m=0}^\infty b_m$ converges, then $\lim_{m
  \to \infty} b_m=0$.  The important counter-example of the {\bf
  harmonic series} $\sum_{m=1}^{\infty} 1/m$ shows that $\lim_{m \to
  \infty} b_m = 0$ is a necessary but not sufficient condition to
  guarantee convergence.

\item The {\bf geometric series} $\sum_{m = 0}^{\infty} x^m$ converges
  for all $x$ with $|x| < 1$ and diverges for all other $x$.

\item An {\bf alternating series} $\sum_{m=0}^{\infty} (-1)^mb_m$
  converges if $0 < b_{m+1} \leq b_m$ for all $m$ and $\lim_{m \to
    \infty} b_m = 0$.  For a convergent alternating series, the error
  in approximating the sum by a partial sum is less than the next term
  in the series.

\item A convergent power series converges on an {\bf interval of
  convergence} of width $2R$; $R$ is called the {\bf radius of
  convergence}.  The {\bf ratio test} can be used to find the radius
  of convergence of a power series: $\sum_{m=0}^{\infty} b_m$
  converges if $\lim_{m \to \infty} |b_{m+1}|/|b_m| < 1$.

\item A polynomial $P(x)=a_0+a_1 x+a_2 x^2 + \cdots + a_n x^n$ is the
  {\bf best fitting} approximation to a function $f(x)$ on an interval
  $[a,b]$ if $a_0, a_1, \cdots, a_n$ are chosen so that the {\bf mean
    squared separation} between $P$ and~$f$
\[
\frac{1}{b-a} \int_a^b (P(x)-f(x))^2 \, dx
\]
is as small as possible. The polynomial $P$ is also called the {\bf
  least squares approximation} to $f$ on $[a,b]$.


\end{itemize}


\subsection{Expectations}

\begin{itemize}

\item Given a differentiable function $f(x)$ at a point $x=a$, you
  should be able to write down any of the {\bf Taylor polynomials} or
  the {\bf Taylor series} for $f$ at $a$.

\item You should be able to use the program TAYLOR to graph Taylor
  polynomials.

\item You should be able to obtain new Taylor polynomials by
  substitution, differentiation, anti-differentiation and
  multiplication.

\item You should be able to use Taylor polynomials to find the value
  of a function to a specified degree of accuracy, to approximate
  integrals and to find limits.

\item You should be able to determine the order of magnitude of the
  agreement between a function and one of its Taylor polynomials.

\item You should be able to find the power series solution to a
  differential equation.

\item You should be able to test a series for divergence; you should
  be able to check a series for convergence using either the {\bf
    alternating series test} or the {\bf ratio test}.

\item You should be able to find the sum of a {\bf geometric series}
  and its interval of convergence.

\item You should be able to estimate the error in an approximation
  using partial sums of an alternating series.

\item You should be able to find the {\bf radius of convergence} of a
  series using the ratio test.

\item You should be able to set up the equations to find the {\bf
  least squares} polynomial approximation of a particular degree for a
  given function on a specified interval.  Working by hand or, if
  necessary, using a computer package to solve a system of equations,
  you should be able to find the coefficients of the least squares
  approximation.

\end{itemize}
 


\cleardoublepage


